% YM4_BA, standalone manuscript source.
% Standardized from manuscript/ym-rotation/main_reorganized.tex.
% Revised mathematical manuscript; source and audit records are kept separately.
\documentclass[11pt,letterpaper,reqno]{amsart}
% BEGIN INLINE manuscript-style
\makeatletter
\NeedsTeXFormat{LaTeX2e}[2020/10/01]
% Use with \documentclass[11pt,letterpaper,reqno]{amsart}.
% This file controls presentation only. It supplies no mathematical content.

\RequirePackage[T1]{fontenc}
\RequirePackage{lmodern}
\RequirePackage{microtype}
\RequirePackage[letterpaper,margin=1in,headheight=15pt]{geometry}
\RequirePackage{amsmath,amsthm,amssymb,mathtools}
\RequirePackage{enumitem}
\RequirePackage{perpage}
\RequirePackage{longtable}
\RequirePackage{xcolor}
\definecolor{manuscriptred}{rgb}{0.6,0.1,0.1}
\definecolor{manuscriptblue}{rgb}{0.1,0.3,0.5}
\definecolor{manuscriptmagenta}{HTML}{60445C}
\RequirePackage[breaklinks=true,colorlinks=true,
  linkcolor=manuscriptred,citecolor=manuscriptblue,urlcolor=manuscriptmagenta]{hyperref}
\RequirePackage{cleveref}
\hypersetup{pdfauthor={}}

% Retain the AMS paragraph, title, heading and running-head conventions.
% Equation counters restart each section: 2.i, 2.ii, ... .
% Appendix sections keep their letters: A.i, A.ii, ... .
\numberwithin{equation}{section}
\renewcommand{\theequation}{\thesection.\roman{equation}}
% Prose references include "eq." and omit parentheses around the number.
\renewcommand{\eqref}[1]{\textup{eq.~\ref{#1}}}
\crefname{equation}{eq.}{eqs.}
\Crefname{equation}{eq.}{eqs.}
\creflabelformat{equation}{#2#1#3}
\setcounter{tocdepth}{1}
\renewcommand{\thefigure}{\Roman{figure}}
\renewcommand{\thetable}{\Roman{table}}
\MakePerPage{footnote}
\renewcommand{\thefootnote}{%
  \ifnum\value{footnote}>9\arabic{footnote}\else\fnsymbol{footnote}\fi}

% Upright statements. Main results are A, B, C; supporting statements
% share a section-local sequence. These end marks indicate typography,
% never proof completeness, validation, or a deferred-proof status.
\theoremstyle{definition}
\newtheorem{ms@theorem}{Theorem}
\renewcommand{\thems@theorem}{\Alph{ms@theorem}}
\newtheorem*{ms@theoremstar}{Theorem}
\newtheorem*{ms@theoremzero}{Theorem 0}
\newtheorem{ms@proposition}{Proposition}[section]
\newtheorem{ms@lemma}[ms@proposition]{Lemma}
\newtheorem{ms@definition}[ms@proposition]{Definition}
\newtheorem{ms@corollary}[ms@proposition]{Corollary}
\newtheorem*{ms@corollarystar}{Corollary}

\newcommand{\ms@majorstart}{%
  \pushQED{\qed}\renewcommand{\qedsymbol}{$\blacksquare\square$}}
\newcommand{\ms@minorstart}{%
  \pushQED{\qed}\renewcommand{\qedsymbol}{$\blacksquare$}}
\newenvironment{theorem}[1][]{\ms@majorstart\ms@theorem[#1]}
  {\popQED\endms@theorem}
\newenvironment{theorem*}[1][]{\ms@majorstart\ms@theoremstar[#1]}
  {\popQED\endms@theoremstar}
\newenvironment{theoremzero}[1][]{\ms@majorstart\ms@theoremzero[#1]}
  {\popQED\endms@theoremzero}
\newenvironment{proposition}[1][]{\ms@minorstart\ms@proposition[#1]}
  {\popQED\endms@proposition}
\newenvironment{lemma}[1][]{\ms@minorstart\ms@lemma[#1]}
  {\popQED\endms@lemma}
\newenvironment{definition}[1][]{\ms@minorstart\ms@definition[#1]}
  {\popQED\endms@definition}
\newenvironment{corollary}[1][]{\ms@majorstart\ms@corollary[#1]}
  {\popQED\endms@corollary}
\newenvironment{corollary*}[1][]{\ms@majorstart\ms@corollarystar[#1]}
  {\popQED\endms@corollarystar}
% The source's lemma/proposition variants are deliberately not carried over:
% choose one marker for each statement type rather than encode proof status.

\theoremstyle{remark}
\newtheorem{remark}[ms@proposition]{Remark}
\newtheorem*{remark*}{Remark}
% The ordinary amsthm proof environment retains its open square.

\crefname{ms@theorem}{theorem}{theorems}
\Crefname{ms@theorem}{Theorem}{Theorems}
\crefname{ms@proposition}{proposition}{propositions}
\Crefname{ms@proposition}{Proposition}{Propositions}
\crefname{ms@lemma}{lemma}{lemmas}
\Crefname{ms@lemma}{Lemma}{Lemmas}
\crefname{ms@definition}{definition}{definitions}
\Crefname{ms@definition}{Definition}{Definitions}
\crefname{ms@corollary}{corollary}{corollaries}
\Crefname{ms@corollary}{Corollary}{Corollaries}
\crefname{remark}{remark}{remarks}
\Crefname{remark}{Remark}{Remarks}

% Bibliography configuration for this manuscript.
\makeatother
% END INLINE manuscript-style
% Begin front.tex
\title[Block averaging and gauge geometry]{Block averaging and gauge geometry for four-dimensional lattice Yang--Mills theory}
\date{}
\hypersetup{pdftitle={Block averaging and gauge geometry for four-dimensional lattice Yang--Mills theory},pdfauthor={}}

\begin{document}
\begin{abstract}
We construct the centered contour average of \(\operatorname{SU}(2)\)
connections on periodic four-dimensional lattices. A positive flux
formula gives a Wilson-energy estimate whose leading constant is one.
We also give the local gauge coordinates and the differentiated
averages, with bounds uniform in the number of scales on specified
small-field balls. A transverse choice of fine links yields analytic
coordinates and a Haar density on a fixed-average fibre.
\end{abstract}
\maketitle
\tableofcontents

\section{Introduction}\label{sec:introduction}

A four-dimensional lattice \(\operatorname{SU}(2)\) connection
admits a local, gauge-covariant block average. In small curvature its
Wilson energy increases by at most a factor \(1+O_L(\varepsilon)\) in
one step, where \(\varepsilon\) is the largest plaquette angle.
Small-field coordinates give controlled derivatives of the iterated
average and an analytic parametrization of the local fixed-average
fibre.

See \cref{prop:avg-map,prop:avg-energy,prop:avg-iteration,thm:fibre-chart}
for the precise statements.

The average of the matrix entries of several transports generally
leaves the group. The contour construction replaces this operation
by two averages of logarithms. The first identifies the frame at
each site of a block. The second compares each fine transport across
the two blocks with the central straight transport. This is the
centered prescription of Ba\l aban \cite[pp.~252--253,
eqs.~(0.10)--(0.12)]{Bal87}. His earlier work develops small-curvature
estimates, straightened coordinates and axial gauges for a
single-comb, corner-based prescription
\cite[\S\S B--F]{Bal85a}\cite[\S B]{Bal85g}.
The proofs below establish the estimates for the centered prescription.

The results concern finite averaging maps, their local coordinates,
differentials and measures. Constructing a constrained critical
background also requires inverse estimates and a nonlinear field
equation, beyond the local fibre parametrization proved here.

At the identity the endpoint terms in the link differential are a
coarse gradient. The curl therefore averages the fluxes through
translated squares with positive coefficients. The row and column
totals of this averaging matrix have product one in dimension four.
A separate local tree gauge for each square lets the same estimate
control the nonlinear field. Straightening removes the endpoint
terms before the link differential is iterated. For a fixed average
we instead use the plain map and solve for one fine link per coarse
constraint. These are the two coordinate constructions used below.

\section{Lattices and conventions}\label{sec:lattice}

Fix an odd integer \(L\ge3\) and put \(m=(L-1)/2\). The fine lattice
at spacing \(h>0\) is
\(\Lambda_h=h\mathbb Z^4/(hN_1\mathbb Z\times\cdots\times hN_4\mathbb Z)\).
A step is allowed when \(L\) divides each \(N_\mu\) and
\(N_\mu/L\ge2\). The coarse lattice \(\Lambda_{Lh}\) consists of the
images of \(Lh\mathbb Z^4\) in the same torus. For several steps we
write \(h_j=L^jh_0\), and require at least two sites in every
direction on the last lattice. We use the partition
\[
 B(y)=y+h\{-m,\ldots,m\}^4,\qquad y\in\Lambda_{Lh}.
\]
The block \(B^j(y)\) at depth \(j\) has \(L^{4j}\) fine sites and is
formed by composing these partitions. All centered paths and boxes
are formed in lattice lifts before projection. Counting a repeated
occurrence in a lift means counting it with its multiplicity.

A positive link \(b=(x,\mu)\) joins \(b_-=x\) to \(b_+=x+he_\mu\).
Positive links are indexed by their base point and direction,
including on a circle with two sites. The reverse \(\bar b\)
is the same edge traversed backwards. A plaquette \(p=(x,\mu,\nu)\)
with \(\mu<\nu\) is based at \(p_-=x\) and is oriented first in
direction \(\mu\), then in direction \(\nu\). Unless otherwise
specified, link and plaquette sums run over these independent
indices. We use \([x,x']\) for a straight segment with its indicated
orientation, and \(\mathrm d\) for the ordinary oriented cochain
boundary sum acting on link variables.

Let \(G=\operatorname{SU}(2)\),
\(\mathfrak g=\mathfrak{su}(2)\) in its anti-Hermitian convention,
and \(G_{\mathbb C}=\operatorname{SL}(2,\mathbb C)\), with Lie
algebra \(\mathfrak g_{\mathbb C}=\mathfrak{sl}_2(\mathbb C)\).
We denote the matrix operator norm by \(|\cdot|\) or \(\|\cdot\|\).
On \(\mathfrak g\) this is the Euclidean norm from
\(\langle X,Y\rangle=-\tfrac12\operatorname{tr}(XY)\).
A connection \(U\) assigns a group matrix to each positive link
and has \(U(\bar b)=U(b)^{-1}\). For a path \(\gamma\),
\(U(\gamma)\) is the product in traversal order and equals \(I\)
on the empty path. A site gauge \(g\) acts by
\[
 U^g(b)=g(b_-)\,U(b)\,g(b_+)^{-1}.
\]
Thus \(U^g(\gamma)=g(\gamma_-)U(\gamma)g(\gamma_+)^{-1}\)
and \((U^g)^u=U^{ug}\). We write
\(\operatorname{Ad}(g)X=gXg^{-1}\).
We take \(\log\) near \(I\) with \(\log I=0\), using the
convergent series for \(|U-I|<1\). When the branch is defined,
\(F_p=\log U(\partial p)\) is the unscaled plaquette curvature.
A link increment \(A\) means \(e^AU\) in the same unscaled
convention. A potential per unit length means an increment \(hA\).
Supremum norms of fields use the matrix norm above. Counting
norms use counting measure, and any lattice-volume weights are
specified where used. Constants subscripted by \(L\) may change
from line to line and depend only on the blocking factor.
Empty sums are zero and empty products are the relevant identity.
The maximum of an empty family of nonnegative bounds is zero.

The affine cube-centre lattices $h_j(\mathbb Z+\tfrac12)^4$ can
also be aligned simultaneously with the lattices used here. Translation
by $\tfrac12h_k(1,1,1,1)$ does so at every level $0\le j\le k$:
$h_k/h_j=L^{k-j}$ is odd, so its half differs from $1/2$ by an
integer. This is one translation of the ambient torus between the
two affine lattice families.

% End front.tex
% Begin average-map.tex
\section{The centered block average}\label{sec:avg-map}

A link average must combine transports with common endpoints. We use
Ba\l aban's symmetric contour construction
\cite[(0.4), (0.10)--(0.12)]{Bal87}. The two averages in its definition
serve different purposes. A mean combines alternative paths within each
block. The logarithmic average of loops then compares the paths across
two blocks with a single straight transport.

\subsection{Means of nearby group elements}

The series logarithm on \(\|V-I\|<1\) takes values in
\(\mathfrak g_{\mathbb C}\) on \(G_{\mathbb C}\), since
\(\exp(\operatorname{tr}\log V)=\det V=1\) and the two eigenvalue
logarithms have imaginary parts in \((-\pi/2,\pi/2)\). It takes values
in \(\mathfrak g\) on \(G\) by unitary diagonalization. We also use the
intrinsic logarithm of \(V\in G\setminus\{-I\}\), obtained by assigning
the eigenvalues \(e^{\pm i\theta}\), \(0\leq\theta<\pi\), the
logarithms \(\pm i\theta\). At \(V=I\) it is zero. This logarithm is
continuous and respects conjugation and inversion.

For a determinant-one two-by-two matrix \(X\), the identity
\(X^{-1}-I=\operatorname{adj}(X-I)\) gives
\begin{equation}\label{eq:avg-inverse-defect}
 \|X^{-1}-I\|=\|X-I\|.
\end{equation}
Indeed the adjugate of a two-by-two matrix has the same singular
values in reverse order. The series logarithm therefore admits
both \(X\) and \(X^{-1}\) when either is in its norm-one disk.
It satisfies \(\log X^{-1}=-\log X\), by the scalar identity
for the holomorphic logarithm on the right half-plane.

\begin{lemma}[Local group mean]\label{lem:avg-mean}
Let \(V_1,\ldots,V_N\in G_{\mathbb C}\), \(N\geq1\), and suppose
\(d=\max_{i,j}\|V_iV_j^{-1}-I\|<1/40\). There is a unique
\(V\in G_{\mathbb C}\) satisfying
\begin{equation}
 \max_i\|V_iV^{-1}-I\|\leq 1/10,
 \qquad \frac1N\sum_i\log(V_iV^{-1})=0 .
 \label{eq:avg-mean}
\end{equation}
The mean \(V=\operatorname{av}(V_i)\) is holomorphic and symmetric,
and the maximum in \eqref{eq:avg-mean} is at most \(4d\). For
\(V_i\in G\) it lies in \(G\), and uniqueness holds among real
candidates with that maximum less than one. For real families the
mean commutes with inversion and with simultaneous left and right
multiplication by elements of \(G\). For complex families it commutes with simultaneous left and right
multiplication whenever both families have diameter less than \(1/40\). At a family of identities,
its derivative in left logarithmic coordinates is the arithmetic mean.
\end{lemma}

\begin{proof}
Put \(T_i=V_iV_1^{-1}\) and write the candidate as \(e^X V_1\).
On \(\|X\|\leq1/8\), the map
\[
 X\longmapsto X+\frac1N\sum_i\log(T_i e^{-X})
\]
is a contraction with constant at most \(2/5\).
Indeed, \(\|T_i e^{-X}-I\|\leq(1+d)e^{1/8}-1<1/6\).
On that ball the nonlinear part of the logarithm has Lipschitz
constant at most \(1/5\), by summing its series after one
differentiation. The same argument for the exponential gives a
Lipschitz constant \(e^{1/8}-1\) for \(e^{-X}-I+X\).
The linear term in the displayed map is bounded by \(d\) after
cancellation. The sum of these three bounds is at most
\[
 d+(1+d)(e^{1/8}-1)+\tfrac15(1+d)e^{1/8}<2/5 .
\]
At zero the map has norm at most \(d/(1-d)\), so it takes the ball
into itself and its fixed point satisfies \(\|X\|\leq 2d\).
The iteration starting at zero stays in the traceless algebra and,
for unitary inputs, in its skew-Hermitian real form. The fixed
point is holomorphic by the implicit function theorem, since the
derivative of the fixed point map has norm less than one.
Its relative differences are at most \(e^{3d}-1\leq4d\).
A candidate within \(1/10\) of the first family member in the
relative norm has \(X=-\log(V_1V^{-1})\) of norm less than \(1/8\),
so is the same fixed point. This uniqueness also gives symmetry
under permutations.

For real uniqueness on the larger ball, identify \(G\) with the
unit quaternions, whose Euclidean norm agrees with the operator norm
on their matrix representatives. For any real candidate \(V'\) in
that ball, write \(V_i(V')^{-1}=\cos\theta_i+\sin\theta_i\,q_i\)
with \(q_i\) a unit imaginary quaternion and \(0\leq\theta_i<\pi/3\).
The logarithmic equation says that
\(\sum_i a_i V_i(V')^{-1}\) is real, where
\(a_i=\theta_i/\sin\theta_i>0\) and \(a_i=1\) at zero. Its real
part is positive. Thus \(V'\) is the unit normalization of the
convex combination \(\sum_i a_i V_i/\sum_i a_i\). The combination
is within \(d\) of \(V_1\) and normalization moves it by at most
another \(d\). This puts \(V'\) inside the smaller uniqueness ball.

If \(V_i'=uV_iv\), their normalized matrices satisfy
\(T_i'=uT_i u^{-1}\). Start the two fixed-point iterations at zero.
Each stays in its own contraction ball, so each logarithm argument
stays in the norm-\(1/6\) disk. Whenever the current iterates are
conjugate, the logarithm and exponential series give conjugate
next iterates. Passing to the limits and multiplying by the first
family member proves
\(\operatorname{av}(V_i')=u\operatorname{av}(V_i)v\). For inversion in the real case use
\(V_i^{-1}V=V^{-1}(V_i V^{-1})^{-1}V\).
Finally differentiation of \eqref{eq:avg-mean} at the constant
family gives the derivative formula.
\end{proof}

Complex inversion is also admitted under a quantitative condition.
If $d<1/80$ and $\|V_i\|\|V_i^{-1}\|\le2$ for every $i$, the
inverse family has diameter at most $2d<1/40$, and
\[
 \operatorname{av}(V_i^{-1})=\operatorname{av}(V_i)^{-1}.
\]
Indeed, if $V=\operatorname{av}(V_i)$, then
$\|V_i^{-1}V-I\|\le2\|VV_i^{-1}-I\|\le8d<1/10$.
Equation~\ref{eq:avg-inverse-defect} gives the second inequality
from the mean bound. Moreover
$V_i^{-1}V=V^{-1}(V_iV^{-1})^{-1}V$.
The logarithms therefore sum to zero, and uniqueness applies to
the inverse family.

At any constant family $V_i=V_0$, write the inputs as $e^{A_i}V_0$.
On a universal small supremum ball their mean is $e^M V_0$, where
\begin{equation}\label{eq:mean-quadratic}
 \left\|M-N^{-1}\sum_i A_i\right\|
                    \le C\max_i\|A_i\|^2.
\end{equation}
The constant is independent of $N$. The contraction proof bounds
$M$ uniformly on a larger supremum ball. Cauchy's formula in two
scalar directions gives a uniform second derivative on the smaller ball; the
computed first derivative and Taylor's formula give the estimate.

\subsection{Contours and the three branches}

For \(x\in B(y)\), let \(\mathcal G(y,x)\) be the paths that
complete all steps in one coordinate before starting the next,
using each order of the nonzero coordinates once. The empty path
is used when \(x=y\). Define
\[
 T_U(y,x)=\operatorname{av}\bigl(U(\gamma):\gamma\in\mathcal G(y,x)\bigr),
 \qquad T_U(x,y)=T_U(y,x)^{-1},
\]
where the mean is defined. For a coarse link \(c\) write \(D_c=B(c_-)\cup B(c_+)\),
\(x'=x+c_+-c_-\) for \(x\in B(c_-)\), and
\(S_c=U([c_-,c_+])\). On the torus the displacement in these
expressions is the directed displacement of the link. All path
links are in the rectangular lift of \(D_c\). The first real
branch applies if every physical fine plaquette with all its
vertices in the projected \(D_c\) has defect
\(\|U(\partial p)-I\|<1/(200L^2)\). At the shortest period this
test can include plaquettes outside the rectangular graph.
Define the analytic path formula on positive links by
\begin{equation}
 \begin{split}
 Z_x&=T_U(c_-,x)\,U([x,x'])\,T_U(c_+,x')^{-1}S_c^{-1},\\
 \mathcal M_{\rm an}(U)(c)&=\exp\!\left(L^{-4}\sum_{x\in B(c_-)}\log Z_x\right)S_c.
 \end{split}
 \label{eq:avg-map}
\end{equation}
Its row chart requires contour diameters less than \(1/40\) and
\(\max_x\|Z_x-I\|<1\). The first branch of the real map
\(\mathcal M\) uses \(\mathcal M_{\rm an}\).
Outside the first real branch, let
\[
 Z_{x,\gamma,\gamma'}
   =U(\gamma)U([x,x'])U(\gamma')^{-1}S_c^{-1},
 \qquad
 (\gamma,\gamma')\in\mathcal G(c_-,x)\times\mathcal G(c_+,x').
\]
If all of these elements differ from \(-I\), define the real value
\(\mathcal M(U)(c)\) by replacing each \(\log Z_x\) on the right
side of \eqref{eq:avg-map} by the arithmetic mean of the intrinsic
logarithms of \(Z_{x,\gamma,\gamma'}\) for that \(x\). If any equals
\(-I\), use \(\mathcal M(U)(c)=S_c\). These two rules define
\(\mathcal M\) outside its first real branch; they do not change
\(\mathcal M_{\rm an}\). A conjugation-equivariant traceless logarithm at \(-I\)
cannot be assigned: it would commute with all of \(G\), hence be
scalar and traceless, with exponential \(I\).

Superscript \(j\) denotes composition through \(j\) successive
lattices for each map. The analytic iterate is defined where its
successive row charts are admitted. An individual row is defined
locally if the charts in its backward path dependence are admitted,
without conditions on unused rows. The zeroth iterates are the identity.

\begin{proposition}[Properties of the centered map]\label{prop:avg-map}
On the periodic lattices of Section~\ref{sec:lattice}, the three
real rules define a Borel map. The value and branch on \(c\)
depend only on the fine field on its projected double block.
For real configurations the path formula on the reverse link
gives the inverse value. For a real site gauge \(g\),
\[
 \mathcal M(U^g)(c)=g(c_-)\mathcal M(U)(c)g(c_+)^{-1}.
\]
The map commutes with translations and signed coordinate
permutations preserving the coarse lattice and its centered
blocks. Each real iterate is onto. A row of \(\mathcal M^j\)
reads only the two depth-\(j\) blocks of its endpoints.

On the row charts of \eqref{eq:avg-map}, \(\mathcal M_{\rm an}\)
is holomorphic, including at real inputs outside the first real
branch. Extend its positive-link values to reverse links by inversion.
If the reverse path formula is admitted in its own chart, it gives
the same value. The endpoint covariance identity for
\(\mathcal M_{\rm an}\) holds for complex site gauges whenever both
the original and transformed fields satisfy the row-chart bounds.
\end{proposition}

\begin{proof}
An interchange of two successive steps in distinct coordinate
directions multiplies the path transport on the left by a conjugate
of an oriented plaquette. Steps in the same coordinate commute as
paths, with backtracks cancelled.
For unitary transports the norm of the change is bounded by the
plaquette's distance from the identity. A permutation of the steps
in a contour takes at most \(2L^2\) adjacent interchanges, because
its length is at most \(2(L-1)\). Hence, writing \(\varepsilon\) for the
maximum plaquette distance on \(D_c\), each contour diameter is at
most \(2L^2\varepsilon<1/100\) in the first branch.

Choose a contour \(\gamma\) and its translate \(\gamma'\). Moving
the \(L\) steps of \([x,x']\) to the front of the concatenated path
takes at most \(2L(L-1)\) interchanges. All interchanges stay within
the double block. It follows that
\(\|Z_{x,\gamma,\gamma'}-I\|\leq2L^2\varepsilon\).
Replacing each contour by its mean costs at most \(4\cdot2L^2\varepsilon\)
by \cref{lem:avg-mean} and unitarity. Thus
\(\|Z_x-I\|\leq18L^2\varepsilon<1/10\), which admits the
logarithm and gives \(\|\mathcal M(U)(c)S_c^{-1}-I\|\leq
C L^2\varepsilon\). All expressions within a branch are continuous
on its domain, and its domain is Borel by the displayed branch tests.
The series and the mean lemma give holomorphy of
\(\mathcal M_{\rm an}\) on its stated chart.

Gauge transformation conjugates the loop by the transformation at
\(c_-\). The open contour means transform at their two endpoints
by \cref{lem:avg-mean}. In the complex case a conjugation by \(g\)
increases a norm bound by at most \(\|g\|\|g^{-1}\|\); the series
transforms termwise as soon as the transformed arguments lie in the
chart. Real endpoint conjugation preserves every branch test. Reversing
a link exchanges \(x\) and \(x'\) and replaces each loop \(Z\)
by \(S_c^{-1}Z^{-1}S_c\). For a real field this preserves the
norm test and the exclusion of \(-I\). Its logarithm is
\(-\operatorname{Ad}(S_c^{-1})\log Z\), which gives the inverse
coarse transport. The last branch uses the reversed straight
transport, also the inverse. In the complex case, when both
outer logarithm charts are admitted, the same logarithm identity
follows from \eqref{eq:avg-inverse-defect} and termwise similarity
covariance of the series. Translations and signed coordinate
permutations permute the contour families, centered blocks and
plaquette tests. The real reversal identity covers those
permutations which reverse a link.

The literal paths are in the rectangular double block, and the
real tests use physical plaquettes with vertices in its projection.
Each depth-\(j\) block is a union of the preceding blocks. Thus
both types of reads stay in the projected union of the two
depth-\(j\) blocks when the map is composed.

For surjectivity fix a field on level \(j\). Assign identity to
fine links inside its depth-\(j\) blocks and its prescribed link
value to every directed fine crossing of the corresponding face.
Use inverses on reversed links. We claim that the same description
holds after each step, with fine links replaced by links of the
current level \(r\le j\). For \(r<j\), the centered level-\(r\)
blocks below level \(r+1\) are contained in the top blocks.
Their internal contour paths therefore have identity transport.
The two endpoints of every averaging segment are in the same
top blocks as the corresponding endpoints of the central
segment. In the direction of the segment a top block is an
interval of \(L^{j-r}\) consecutive level-\(r\) sites. The
segment is monotone and has length \(Lh_r\le h_j\), so it crosses
at most one top face. The transverse coordinates of the segment
remain in the transverse top block of the central segment.
Hence every averaging segment has the same directed face
transport as the central one, or they are all internal. This
is a statement on the lift. At a two-site top period the two
positive coarse bonds with the same endpoint set still have
distinct initial sites and distinct lifted faces. Every closed
loop used by either logarithmic rule is now \(I\). Both rules
give the common segment transport. This proves the induction
and gives exactly the prescribed field at level \(j\).
\end{proof}

For a finite measure \(\nu\) on fine connections, the pushforward
\(\mathcal M_*\nu\) preserves total mass and preserves positivity.
Gauge invariance also passes to this measure because every coarse
gauge extends to a fine gauge and \(\mathcal M\) is covariant.

The path count has a general form. Reordering a path of $m$ steps
by adjacent swaps takes at most $m(m-1)/2$ swaps. If every
intermediate path and elementary plaquette lies in the chosen
region, each swap in distinct coordinates inserts a conjugate
plaquette; a swap in one coordinate leaves the transport unchanged.
For a real field
with plaquette defect at most $\epsilon$, unitary telescoping
therefore bounds the change in the path transport by
$m(m-1)\epsilon/2$. The containment assumption prevents applying
this estimate across a missing part of the carrier.

\subsection{Local comb coordinates}

A local gauge compares each edge with the coordinate paths from
a root to its endpoints. The comparison uses the curvature only
inside the box.

\begin{lemma}[Comb estimate on a box]\label{lem:avg-comb}
Let \(D\) be a rectangle in a lift of the fine lattice, with
at most \(s\) sites in each direction and a root \(y\in D\).
Use its literal rectangular edges and plaquettes. For a field
\(U\) on these edges, let \(\gamma_x\) be the path from \(y\)
to \(x\) in a fixed coordinate order. Put
\[
 g(x)=U(\gamma_x),\qquad
 \varepsilon=\max_{p\subset D}\|U(\partial p)-I\|,
\]
with \(\varepsilon=0\) when the box has no plaquettes.
If \(U\) is real, then
\[
 \|U^g(b)-I\|\le C s\varepsilon\qquad (b\subset D).
\]
If also \(C s\varepsilon\le1/2\), the logarithms
\(a_b=\log U^g(b)\) obey the same pointwise bound up to an
absolute factor and satisfy
\[
 \left(\sum_{b\subset D}\|a_b\|^2\right)^{1/2}
 \le C_s\left(\sum_{p\subset D}\|F_p\|^2\right)^{1/2}.
\]
The first constant is absolute, and \(C_s\) depends only on \(s\).
For a complex field, if every path conjugating a plaquette to
the root in the comb comparison has condition number at most
\(K\), then
\(\|U^g(b)-I\|\le \exp(C sK\varepsilon)-1\).
All conclusions concern the lift when a periodic vertex has
more than one occurrence.
\end{lemma}

\begin{proof}
Append \(b\) to the comb to \(b_-\). Move this one extra step
to its coordinate position and cancel an opposite step if
needed. This gives the comb to \(b_+\) in at most \(Cs\)
interchanges inside \(D\). Each interchange contributes a
plaquette transported to the root by a prefix path. If the
plaquette has reverse orientation, write it as the inverse
of the positively oriented plaquette at its standard base
point. Include any change of base point in the transporting
path. The product of these transported plaquette factors is
\(U(\gamma_{b_-})U(b)U(\gamma_{b_+})^{-1}\).

In the real case all conjugations are isometric and the factors
are unitary. Telescoping the product gives \(Cs\varepsilon\).
In the complex case \eqref{eq:avg-inverse-defect} handles
inverse factors, so the deviation of each transported factor
is at most \(K\varepsilon\). Multiplying the bounds
\(1+K\varepsilon\) yields \(\exp(CsK\varepsilon)-1\).
The transporting paths have length \(O(s)\), since each is a
prefix of a reordered comb with the extra step and at most
the rebase to its plaquette.

On the norm-half disk the logarithm has bounded derivative.
In the real case this bounds \(\|a_b\|\) by a constant times
the sum of \(\|F_p\|\) over its plaquette occurrences. The
finite incidence matrix has row and column counts bounded
in terms of \(s\). Cauchy--Schwarz in each row followed by
the column count proves the squared-sum estimate.
\end{proof}

\subsection{The flat derivative and its curl}

For a fine cochain \(a\) at the identity use \(a(\bar b)=-a(b)\)
and let \(a(\gamma)\) denote the sum along a path. Put
\[
 H_a(y)=L^{-4}\sum_{x\in B(y)}
       \frac1{|\mathcal G(y,x)|}\sum_{\gamma\in\mathcal G(y,x)}a(\gamma).
\]
Let \(Q a\) be the derivative of \(\log\mathcal M_{\rm an}(e^{ta})\) at \(t=0\),
with the fine and coarse perturbations measured as raw link
logarithms. The linearization of the mean and the product rule give
\begin{equation}
 (Qa)(c)=L^{-4}\sum_{x\in B(c_-)}a([x,x'])
                           +H_a(c_-)-H_a(c_+).
 \label{eq:avg-flat-link}
\end{equation}
The derivative of the central straight transport cancels against
its occurrence inside the logarithm in \eqref{eq:avg-map}.
In particular, if all \(a(b)\) are real multiples of a fixed element
of \(\mathfrak g\) and \(\|a\|_\infty\) is sufficiently small
depending on \(L\), then \(\mathcal M(e^a)=e^{Qa}\) exactly.
The products commute, the small logarithms are additive, and
the defining equation of the mean reduces to an arithmetic average.

Let \(\mathrm d\) and \(\mathrm d_c\) be the oriented boundary
sums around fine and coarse plaquettes. If \(C=(y,\mu,\nu)\) is a
coarse plaquette with initial vertex \(y\) and oriented coordinate
directions \(\mu,\nu\), let \(S_x(C)\) be the \(L\)-by-\(L\) fine
surface parallel to \(C\) starting at \(x\in B(y)\). Then
\begin{equation}
 (\mathrm d_cQa)(C)=L^{-4}\sum_{x\in B(y)}\sum_{p\in S_x(C)}
                         (\mathrm da)(p).
 \label{eq:avg-flat-curl}
\end{equation}
The terms \(H_a\) cancel in the oriented boundary sum. The four
straight segments left for each \(x\) bound \(S_x(C)\), and its
interior links cancel in the sum of fine boundaries. The resulting
positive flux weights have row sum \(L^2\). Their column sum
enters the energy estimate in \eqref{eq:energy-schur}.
After \(j\) iterations the flux weights still have nonnegative
entries, with row sum \(L^{2j}\).


% End average-map.tex
% Begin energy.tex
\section{Curvature and Wilson energy}\label{sec:energy}

Write \(\varepsilon(U)=\max_p |F_p|\) when all the plaquette logarithms
are defined in the small branch. For real fields put
\[
 \mathcal E(U)=\sum_p\left(1-\tfrac12\operatorname{tr}U(\partial p)\right).
\]
This sum is dimensionless and counts each plaquette once. Since the
eigenvalues of \(F\in\mathfrak g\) are \(\pm i|F|\), the corresponding
summand is \(1-\cos |F|\). In particular, \(\mathcal E(U)\) is half the
sum of \(\lvert U(\partial p)-I\rvert^2\) in the operator norm.

\begin{theorem}[One-step energy estimate]\label{prop:avg-energy}
There are positive constants \(\varepsilon_L,C_L\), depending
only on \(L\), with the following property. Let \(U\) be a real
connection on a torus admitting one blocking step. Suppose its
fine plaquette holonomies are in the series logarithm chart and
\(\varepsilon(U)=\max_p\|\log U(\partial p)\|\le\varepsilon_L\).
Then each output link uses the first real branch, the output
plaquette holonomies lie in the series logarithm chart, and
\begin{align}
 \varepsilon(\mathcal M U)&\le L^2\varepsilon(U)+C_L\varepsilon(U)^2,
       \label{eq:energy-curv}\\
 \mathcal E(\mathcal M U)^{1/2}&\le
      (1+C_L\varepsilon(U))\mathcal E(U)^{1/2}.
       \label{eq:energy-step}
\end{align}
\end{theorem}

\begin{proof}
We work in fine lattice units. Let \(C=(y,\mu,\nu)\) denote a coarse
plaquette. The differential of its logarithmic holonomy is the curl
of the flat link differential. By \eqref{eq:avg-flat-curl}, this curl is
a sum of fine plaquette curls with weights \(L^{-4}\) indexed by
\[
 (x,i,j)\in B(y)\times\{0,\ldots,L-1\}^2,\qquad
 p=(x+ie_\mu+je_\nu,\mu,\nu).
\]
Write \(w_{C,p}\) for the resulting nonnegative coefficients,
retaining separate occurrences when a lift uses the same physical
plaquette more than once. Each coarse row has total weight \(L^2\).
For a fixed physical \(p\), each pair \((i,j)\) determines the fine
site \(x\) and its unique block origin \(y\). Thus each physical
column has total weight \(L^{-2}\). For any orthogonal maps \(O_{C,p}\)
of \(\mathfrak g\) and any real plaquette field \(f\), weighted
Cauchy--Schwarz gives
\begin{equation}\label{eq:energy-schur}
 \sum_C\left|\sum_p w_{C,p}O_{C,p}f_p\right|^2
 \le L^2\sum_{C,p}w_{C,p}|f_p|^2
 =\sum_p |f_p|^2.
\end{equation}
The norm here comes from the inner product
\(-\tfrac12\operatorname{tr}(XY)\) on \(\mathfrak g\). The product
of the row and column totals is one in dimension four.

To apply this calculation at \(U\), lift the four endpoint blocks
of \(C\) together to a rectangular box \(D_C\). It has \(2L\)
sites in each plaquette direction and \(L\) in each transverse
direction. Take the coordinate tree gauge
of \cref{lem:avg-comb}, rooted at \(y\), and write the links in
this gauge as \(e^{a_b}\). For a constant depending only on \(L\),
\[
 \|a\|_{\infty,D_C}\le C_L\varepsilon(U),\qquad
 \|a\|_{2,D_C}\le C_L\|F\|_{2,D_C}.
\]
The norms count the edges and plaquettes of the rectangular
lift. Its vertices have distinct projections, since each
fine period has at least \(2L\) sites. Extend the real gauge
to the other sites. This preserves the defects of every
physical plaquette, including tests across an opposite face
when a direction has exactly two coarse sites. The global
curvature hypothesis puts all four output links in the first
real branch. Thus their plaquette product is the one obtained
from the analytic path formula in the literal box. On a fixed
small ball of the \(a_b\), its contour means and logarithms are
analytic and its coarse plaquette product stays in the series
logarithm chart. Reducing \(\varepsilon_L\) puts the actual
local logarithms in that ball.

In this gauge the fine logarithms have the form
\((\mathrm da)_p+r_p\), where
\(\|r\|_{2,D_C}\le C_L\|a\|_{\infty,D_C}\|a\|_{2,D_C}\).
Indeed each plaquette product has four factors and its first
derivative is the oriented boundary sum. Taylor's formula bounds
the error by a constant times the maximum of the four log norms
times their sum. Each edge occurs in a bounded number of faces.
The coarse plaquette logarithm has the flat curl as its first
derivative, with Taylor error at most
\(C_L\|a\|_{\infty,D_C}\|a\|_{2,D_C}\), by analyticity on a
fixed smaller ball of this finite box. Substituting the fine
logarithms into the flat formula and conjugating back to the
original frames therefore gives
\begin{equation}\label{eq:energy-flux}
 F_C(\mathcal M U)=\sum_p w_{C,p}O_{C,p}F_p(U)+r_C,\qquad
 |r_C|\le C_L\varepsilon(U)\|F\|_{2,D_C}.
\end{equation}
For the box gauge \(g_C\), write the coarse Taylor error in
that frame as \(r_C^{\rm out}\), and the fine errors as \(r_p\).
Then the terms after returning to the original frame are
\[
 O_{C,p}=\operatorname{Ad}\bigl(g_C(y)^{-1}g_C(p_-)\bigr),
 \qquad
 r_C=\operatorname{Ad}(g_C(y)^{-1})
       \left(r_C^{\rm out}-\sum_pw_{C,p}r_p\right).
\]
Here \(p_-\) is its occurrence in the box. Since the
nonnegative row weights sum to \(L^2\), their squared sum
has square root at most \(L^2\). Cauchy--Schwarz bounds the
second error term by \(L^2\|r\|_{2,D_C}\). No cancellation
between different box frames is used. For each of the six
coarse plaquette orientations, a given fine plaquette base
can lie in at most two translated boxes along each of the two
long directions and at most one along each other direction.
Thus a physical fine plaquette occurs in at most \(24\) boxes,
including periodic multiplicities.

The first inequality follows from the row total and the finite
size of each box. The family of lifted boxes has bounded overlap,
counted with multiplicity on the torus. Indeed each box is a
translate of one of six shapes of diameter \(O(L)\) by a coarse
lattice vector. Consequently \(\|r\|_2\le C_L\varepsilon(U)\|F\|_2\).
Using \eqref{eq:energy-schur} in \eqref{eq:energy-flux} yields
\(\|F(\mathcal M U)\|_2\le (1+C_L\varepsilon(U))\|F(U)\|_2\).
For the input use
\(1-\cos t\ge \tfrac12t^2(1-Ct^2)\) on the small interval.
For the output use \(1-\cos t\le \tfrac12t^2\). The resulting
factor \((1-C\varepsilon(U)^2)^{-1/2}\) is absorbed in the
right side of \eqref{eq:energy-step} by reducing \(\varepsilon_L\).
\end{proof}

\begin{corollary}[Iteration]\label{cor:avg-real-iterate}
There is \(\alpha_L>0\) such that, for every torus admitting \(k\ge0\)
steps and every real field whose fine plaquettes admit the series
logarithms and satisfy
\(\varepsilon(U)\le\alpha L^{-2k}\) for \(0\le\alpha\le\alpha_L\),
all \(k\) steps use the first real branch and
\begin{equation}\label{eq:energy-depth}
 \varepsilon(\mathcal M^j U)\le 2\alpha L^{-2(k-j)},\qquad
 \mathcal E(\mathcal M^j U)\le 2\mathcal E(U)
 \quad(0\le j\le k).
\end{equation}
\end{corollary}

\begin{proof}
Let \(\varepsilon_j=\varepsilon(\mathcal M^j U)\).
As long as \(\varepsilon_i\le 2\alpha L^{-2(k-i)}\) for \(i<j\),
their sum is at most \(2\alpha/(L^2-1)\). Choose
\(2\alpha_L\le\varepsilon_L\), so that every step through \(j\)
is defined. Dividing \eqref{eq:energy-curv} by the factor \(L^2\)
at each step and using \(1+s\le e^s\) gives
\[
 L^{-2j}\varepsilon_j
 \le \varepsilon_0 \exp\!\left(C_L L^{-2}\sum_{i<j}\varepsilon_i\right)
 \le 2\alpha L^{-2k}
\]
after reducing \(\alpha_L\). This proves the induction, including the
case \(\varepsilon_0=0\). The product of the factors in
\eqref{eq:energy-step} is at most
\(\exp(2C_L\alpha/(L^2-1))\). Choose \(\alpha_L\) so that its
square is at most two.
\end{proof}

\begin{corollary}[Curvature on complete local stocks]
\label{cor:avg-local-real-iterate}
On a torus admitting \(k\) steps, let \(Y\) be a union of
depth-\(k\) blocks. Let \(U\) be real, and suppose every
physical fine plaquette with all its vertices in \(Y\) has
its series logarithm with norm at most \(\alpha L^{-2k}\),
where \(0\le\alpha\le\alpha_L\). For \(0\le j\le k\), the
level-\(j\) plaquettes with vertices in \(Y\) have series
logarithms of norm at most \(2\alpha L^{-2(k-j)}\).
Every level-\(j\) averaging row with endpoints in \(Y\),
\(1\le j\le k\), uses the first real branch and depends only
on fine links with endpoints in \(Y\).
\end{corollary}

\begin{proof}
At every level at most \(k\), the block of a site in \(Y\)
is contained in \(Y\) by alignment. For a row whose endpoints
are in \(Y\), all its paths and all its real test plaquettes
at the preceding level therefore have vertices in \(Y\).
This includes projected opposite-face tests. The rectangular
lift for any plaquette with corners in \(Y\) projects into
the four endpoint blocks, which also lie in \(Y\).
The pointwise argument of \cref{prop:avg-energy} thus applies
using only the largest defect at the preceding level on \(Y\).
At steps below \(k\), the inductive bound keeps the real tests
inside their cutoff after a fixed decrease of \(\alpha_L\).
The same curvature recurrence and geometric sum as in
\cref{cor:avg-real-iterate} give the claim. This argument is
local and does not require \(Y\) to be connected.\end{proof}

\subsection{A planar interface}

The exact flux formula also describes the loss when the curvature
hypothesis changes at a block face. In fine units put the tail
centre at \(0\), take \(C\) in directions \(1,2\), and set
\(a_2(x)=f(x_1)X\), with \(|X|=1\) in \(\mathfrak g\), and
\(a_\mu=0\) for \(\mu\ne2\). Choose the first difference of \(f\)
to equal \(\phi\) on \(-m,\ldots,m-1\), \(-\phi\) on the preceding
\(L-1\) sites and zero elsewhere. This has a periodic primitive
on a sufficiently large torus. Take \(|\phi|\) small enough
depending on \(L\) so all the logarithms are on their analytic
branches. All factors commute, and the coarse curvature is
exactly the linear flux in \eqref{eq:avg-flat-curl}.

The first direction of a surface starting at \(x\in B(0)\)
contributes \(m-x_1\) copies of \(\phi\) and no negative copies.
Its second direction contributes a factor \(L\). Averaging
over \(x_1\) gives the signed coarse curvature coefficient \(L(L-1)\phi/2\), although
the curvature already vanishes in the head block. If instead
the fine curvature is constant on the full flux support,
the signed coarse curvature coefficient is \(L^2\phi\). This second configuration can
also be realized with a periodic primitive by compensating the
first differences outside the support. These examples locate
the curvature cost of using the full read of a row across a face.

% End energy.tex
% Begin gauges.tex
\section{Relative axial gauges and straightening}\label{sec:gauge-straightening}

A field on one block can be put in a fixed gauge by specifying how each
site is transported to the block centre. A second choice fixes the value of
the gauge transformation at that centre. For composing averages, the
useful choice prescribes a transported average of the gauge transformation.

\subsection{Local gauges and their normalization}

On a tree with root \(y\), the gauge \(g(x)=W(\gamma_{yx})\), with
\(\gamma_{yx}\) the tree path, sets every tree link to \(1\). It is unique
among gauges with \(g(y)=1\). Indeed each tree edge from a parent \(x\)
to a child \(z\) gives the recursion \(g(z)=g(x)W(x,z)\).
For real fields this change of variables also preserves the product
Haar measure in the following sense. Use the \(g(x)\) for \(x\ne y\) as
coordinates for the tree links and the transformed links as coordinates
off the tree. First orient each tree coordinate from parent to child;
reversing an independent positive link preserves Haar measure by inversion.
In parent-before-child order each child-gauge substitution is then a
fixed left translation when the previous coordinates are held fixed.
Once the tree coordinates are fixed, each other link is changed by left
and right translations. The resulting measure is product normalized Haar
in these coordinates. The same proof works on a forest with a prescribed
root gauge on each component.

For a block \(B(y)\), let \(T_W(y,x)\) be the mean transport in
\eqref{eq:avg-map}. For two fields \(W,V\), the relative block axial
condition is \(T_{V^g}(y,x)=T_W(y,x)\). This condition fixes the
transformation inside the block once \(g(y)\) is given. If some sites are
pinned, one can impose the condition on the free sites while keeping
\(g(y)\) fixed. At a pinned site the condition is an equation on the
existing transport, so it must be included in the datum whenever it is
needed. In particular the gauge \(g(x)=T_V(y,x)\) fixes the mean
transports to \(1\) with the block centre held fixed. This choice and
the coordinate tree gauge are defined by their respective paths.

The following estimate explains why the relative conditions provide
local coordinates even in the presence of coarse holonomy. Let \(D\)
be the rectangular lift of \(B(c_-)\cup B(c_+)\) along \(c\).
Write
\[
  \epsilon_D(W)=\max_{p\subset D}\|W(\partial p)-1\|.
\]
Here \(p\subset D\) and \(b\subset D\) refer to the rectangular
graph. The real branch test may also read plaquettes that project
across the opposite face of a short periodic lattice.

\begin{lemma}[Relative transport estimate]\label{lem:gauge-rigidity}
Let \(W,V\) be real fields defined on \(D\) and on the plaquette
test stock for \(c\). Suppose
\(\epsilon_D(W),\epsilon_D(V)\le\epsilon\) for sufficiently small
\(\epsilon\) depending on \(L\). If their mean transports in the two
blocks differ by at most \(\delta\) and their averages on \(c\)
differ by at most \(d\), then
\[
 \max_{b\subset D}\|V(b)-W(b)\|\le d+2\delta+C_L\epsilon .
\]
\end{lemma}

\begin{proof}
For an edge \(b\) internal to \(B(y)\), the product
\(T_W(y,b_-)W(b)T_W(y,b_+)^{-1}\) differs from \(1\) by \(C_L\epsilon\).
To see this, compare both mean transports with their coordinate combs
by the path interchange estimate, then close the combs with \(b\) and
use Lemma~\ref{lem:avg-comb}. For a crossing edge with
\(b_-\in B(c_-)\), the corresponding product
\[
 T_W(c_-,b_-)W(b)T_W(c_+,b_+)^{-1}
\]
differs from \(W([c])\) by \(C_L\epsilon\). With combs in place of the
means, the path and \([c]\) bound a closed walk in \(D\) of bounded
length. In a coordinate tree gauge every link on this walk is within
\(C_L\epsilon\) of \(1\), giving the estimate by the unitary product
inequality. The mean transport comparison restores the same bound.
The same closed walk argument bounds every loop in the first two
rules defining \(\mathcal M\) by \(C_L\epsilon\) in distance
from \(I\). In particular all the fallback path loops avoid \(-I\).
Either rule that applies gives
\(\|\mathcal M(W)(c)-W([c])\|\le C_L\epsilon\).
Thus internal links are within
\(C_L\epsilon\) of the product of their endpoint transports, and crossing
links are within \(C_L\epsilon\) of that product with
\(\mathcal M(W)(c)\) inserted between them. Compare the corresponding
formulas for \(V\) and telescope the three unitary factors.
\end{proof}

For a site function \(g\) with values in \(G\) or \(G_{\mathbb C}\), define
its one-step average relative to \(W\) by
\begin{equation}\label{eq:gauge-site-nonlinear}
 R_Wg(y)=g(y)\exp\!\left\{L^{-4}\sum_{x\in B(y)}
 \log\!\left[g(y)^{-1}T_W(y,x)g(x)T_W(y,x)^{-1}\right]\right\}.
\end{equation}
All logarithms are in the branch at \(1\) and the formula is used on
its open chart. For \(V=e^aW\), with \(a\) a raw link logarithm, set
\begin{equation}\label{eq:gauge-block-factor}
 u_W(a)(y)=\exp\!\left\{L^{-4}\sum_{x\in B(y)}
 \log\big[T_V(y,x)T_W(y,x)^{-1}\big]\right\}.
\end{equation}
For contour covariance in complex frames, require the original and
transformed contour families to lie in the mean charts of
Lemma~\ref{lem:avg-mean}.
Prescribe \(R_Wg(y)=s(y)\) at the centre. Suppose the relative contour
logarithms exist. For complex data also require the contour families
of the candidate below after gauge transformation to lie in their
mean charts. The unique admitted relative block axial gauge is then
\begin{equation}\label{eq:gauge-relative-axial}
 g(y)=s(y)u_W(a)(y)^{-1},\qquad
 g(x)=T_W(y,x)^{-1}g(y)T_V(y,x).
\end{equation}
Covariance of \(T\) gives the second formula and makes each argument of
the logarithm in \eqref{eq:gauge-site-nonlinear} equal to
\(T_V(y,x)T_W(y,x)^{-1}\). Thus the prescribed average fixes the value
at the centre. These formulas adapt the transported gauge average and
the axial normalization of Ba{\l}aban
\cite[(78)--(92)]{Bal85a}\cite[\S B]{Bal85g} to the mean transports.

\subsection{Straightened link coordinates}

On the admitted analytic charts the one-step plain coordinate is
\[
 P_W(a)=\log[\mathcal M_{\rm an}(e^aW)\mathcal M_{\rm an}(W)^{-1}].
\]
Define the one-step straightened
coordinate \(F_W\) by
\begin{equation}\label{eq:gauge-straighten}
 e^{F_W(a)(c)}\mathcal M_{\rm an}(W)(c)
 =u_W(a)(c_-)^{-1}\mathcal M_{\rm an}(e^aW)(c)\,u_W(a)(c_+) .
\end{equation}
Take \(F_W(a)(c)\) to be the series logarithm of the right side
multiplied by \(\mathcal M_{\rm an}(W)(c)^{-1}\). Its chart requires
this product to lie in the norm-one disk about \(I\).
If the relative axial gauge in \eqref{eq:gauge-relative-axial}
exists for \(s=1\) and its transformed field lies in the analytic
row chart, covariance identifies the right side with the analytic
average of that transformed field. The two
factors of \(u_W\) subtract the infinitesimal block-transport terms
from the plain linearization. Set \(J_W=D_aF_W(0)\). For a fine
Lie algebra site field \(\lambda\), define the site differential by
\[
 q_W\lambda=\left.\frac{d}{dt}R_W(e^{t\lambda})\right|_{t=0}.
\]
Its explicit row is computed in \eqref{eq:gauge-site-one}.

For later estimates use row and column sums of the norms of colour
matrices. These are counting-coordinate sums. Define the covariant
straight-segment row by
\begin{equation}\label{eq:gauge-straight-row}
 (S_Wa)(c)= L^{-4}\sum_{x\in B(c_-)}
 \sum_{b\in[x,x']}
 \operatorname{Ad}\!\big(T_W(c_-,x)W([x,b_-])\big)a(b),
 \qquad x'=x+(c_+-c_-).
\end{equation}
Here the formula is written on positive \(c\), with segment edges
oriented along it. If \(H\) is a field on oriented links, its raw
multiplicative tangent at \(H\) extends from positive edges by
\begin{equation}\label{eq:gauge-reverse-tangent}
 a(\bar b)=-\operatorname{Ad}(H(b)^{-1})a(b).
\end{equation}
This follows by inverting \(e^{a(b)}H(b)\). It specifies the induced
action of all lattice symmetries, including those that reverse links.
In particular \(u_W(a)\) pulls back as a site field, by equivariance
of the contour family. The field on the right of
\eqref{eq:gauge-straighten} then pulls back as a link field, so its
logarithm pulls back by \eqref{eq:gauge-reverse-tangent} at the
transformed reference. Differentiation gives the same symmetry
for \(J_W\) and its iterates.

Let \(s\) be a lattice isometry preserving the fine and coarse
grids and the centered block partition. Write \(s^*\) for its
pullback and \(s_W^*\) for its induced map on logarithmic tangents.
When both the original and pulled-back path formulas and their
logarithmic coordinates lie in the admitted analytic charts,
\begin{equation}\label{eq:gauge-isometry-laws}
 D(\mathcal M_{\rm an})_{s^*W}\,s_W^*
  =s_{\mathcal M_{\rm an}(W)}^*D(\mathcal M_{\rm an})_W,
 \qquad q_{s^*W}s^*=s^*q_W.
\end{equation}
On the open first real branch the analytic differential restricted
to real tangents is the Borel map's differential. The corresponding
analytic identity holds for the straightened derivative \(J_W\).
Pullbacks compose in the
ordinary pullback order; differentiating that exact identity gives
the composition law for the background-dependent tangent maps.
For real $W$, a reversed edge contributes only a sign and a unitary
adjoint in \eqref{eq:gauge-reverse-tangent}. The maps are therefore
isometries in the counting supremum and physical $L^p$ norms,
$1\le p\le\infty$. This argument uses the permutation of the
underlying links and their unchanged lengths.

\begin{proposition}[One-step derivative]\label{prop:gauge-one-step}
There are positive constants \(\epsilon_L,r_L,C_L\), depending
only on \(L\), with the following property on a torus admitting
one averaging step. Let \(W\) be real, with plaquette defect at most
\(0\le\epsilon\le\epsilon_L\) on the determining double blocks and
their projected first-branch tests. For a complex Lie algebra
increment \(a\) on the fine positive links, \(F_W(a)\) and
\(\log u_W(a)\) are
holomorphic for \(\|a\|_\infty<r_L\), with bounds
\(C_L\|a\|_\infty\). Then
\[
 \|S_W\|_{\rm row}\le L,\qquad
 \|S_W\|_{\rm col}\le L^{-3},\qquad
 \|J_W-S_W\|_{\rm row}+\|J_W-S_W\|_{\rm col}
 \le C_L\epsilon .
\]
On a smaller fixed ball, \(D_aF_W(a)-J_W\) has both sums bounded
by \(C_L\|a\|_\infty\). The second derivative obeys
\[
 \sup_b\sum_{c,b'}\|D_bD_{b'}F_W(a)(c)\|\le C_L.
\]
These estimates are uniform in the volume. For \(W'=e^{a'}W\) with
\(\|a'\|_\infty<r_L\) after decreasing \(r_L\), the same assertions
hold holomorphically in \((a',a)\) with
\(\epsilon+\|a'\|_\infty\) in the comparison bound and with
\(L(1+C_L\|a'\|_\infty)\) and
\(L^{-3}(1+C_L\|a'\|_\infty)\) in the bounds for \(S_{W'}\).
\end{proposition}

\begin{proof}
Work on one determining double block in the coordinate tree gauge of
Lemma~\ref{lem:avg-comb}. In this gauge \(W\) has link logarithms
bounded by \(C_L\epsilon\). Conjugation by this real gauge preserves the
norms of the input and output entries. A fixed sufficiently small
joint ball in the reference logarithms and the increments admits all
path products, relative logarithms and mean equations in
\eqref{eq:gauge-block-factor} and \eqref{eq:gauge-straighten}.
The paths have length bounded in terms of \(L\). The mean transport and
the remaining operations are holomorphic and bounded on this ball,
uniformly in the choice of the row. Cauchy's formula on a smaller ball
bounds each first and second partial derivative and each mixed
reference and increment derivative by \(C_L\). A row contains at most
\(C_L\) fine bonds, and a fine bond belongs to at most a fixed number
of determining double blocks. These two counts give the claimed
sums, once the derivative at the trivial field is identified.
The same chart contains \(e^{a'}W\) and its increments, since
\(\log(e^{a'}W)\) in the local gauge has size at most
\(C_L\epsilon+C_L\|a'\|_\infty\).
For real increments, a tested plaquette changes by at most
\(4(\|a'\|_\infty+\|a\|_\infty)\), by telescoping its four
unitary factors. Choose the joint radii so that this change and
the reference defect fit below the first-branch cutoff. The test
stock includes projected opposite-face plaquettes on a shortest
period. Thus the analytic formula agrees with the real map there.
The same analytic argument applies on embedded covers of periodic carriers. Counting a repeated occurrence
on a cover and adding its entry on descent only increases the bounds.

The derivative of \(\log u_1\) is \(H_a\), the block mean of contour
sums defined before \eqref{eq:avg-flat-link}. Its two endpoint terms cancel
those in \(DP_1(0)\), leaving \(J_1=S_1\).
Both \(J_W\) and \(S_W\) are holomorphic functions of
the finite set of reference links. Integrating their mixed derivative
from the trivial reference in the local gauge gives the comparison
bound. Integrating the second increment derivative gives the bound
for \(D_aF_W(a)-J_W\). The values \(F_W(0)=0\) and
\(\log u_W(0)=0\) give the first-order size bounds.

Each real conjugation in \eqref{eq:gauge-straight-row} has norm one.
There are \(L^4\) segments of \(L\) links in one row. For a fixed input
edge parallel to the row direction, its possible positions in a
segment number \(L\), and each such position determines the start
and its block. Its total column weight is \(L^{-4}L\).
For the complex reference comparison the transported matrices differ
from the corresponding real ones by \(C_L\|a'\|_\infty\), by
holomorphy on the same finite ball. This multiplies both counting
bounds by at most \(1+C_L\|a'\|_\infty\).
\end{proof}

\subsection{Composition}

For a reference \(W\) on the fine lattice set
\(W_j=\mathcal M_{\rm an}^j(W)\) on the admitted iteration charts.
For a field \(e^{a_0}W\), define \(a_j\) and the accumulated site factors
\(v_j\) by
\begin{equation}\label{eq:gauge-recursion}
 v_0=1,\qquad a_{j+1}=F_{W_j}(a_j),\qquad
 v_{j+1}=(v_j|_{\mathrm{level}\ j+1})u_{W_j}(a_j).
\end{equation}
Require the plain iterates and the endpoint conjugations by \(v_j\)
to remain in the analytic row charts of Proposition~\ref{prop:avg-map}.
The small balls in Proposition~\ref{prop:avg-iteration} supply this
admission. Covariance and \eqref{eq:gauge-straighten} then give
\begin{equation}\label{eq:gauge-factorization}
 \mathcal M_{\rm an}^j(e^{a_0}W)=(e^{a_j}W_j)^{v_j}.
\end{equation}
Indeed applying \(\mathcal M_{\rm an}\) to the right side restricts \(v_j\) to
the next lattice. Applying \eqref{eq:gauge-straighten} there gives
\(e^{a_{j+1}}W_{j+1}\) with gauge \(u_{W_j}(a_j)\), and the two
gauge transformations compose in the order of \eqref{eq:gauge-recursion}.
For real fields the same identity holds with the Borel map when the
relevant rows use the first branch at each level. The local identity
uses the backward determining blocks of the output, including the
projected tests for the Borel map.

The factor \(v_1\) can contribute already at first order. At \(W=1\),
perturb a single edge leaving a block centre \(y\) by \(e^{tX}\),
where \(X\in\mathfrak g\) and \(t\) is small. A coordinate path
can use that edge only as its first step. The total weight of first
steps, averaged over the endpoints \(x\in B(y)\), is \(1-L^{-4}\).
Signed coordinate permutations distribute it equally among the eight
edges leaving \(y\). Since all perturbed products commute, the contour
mean averages their logarithms exactly on this chart. Hence
\[
 v_1(y)=\exp\bigl(t(1-L^{-4})X/8\bigr).
\]
All other block factors are \(1\). In particular the site terms in
\eqref{eq:gauge-straighten} affect the linearized link coordinate.

% End gauges.tex
% Begin average-estimates.tex
\section{Estimates through several scales}\label{sec:avg-iteration}

We measure all spacings in units of the final spacing \(h_k\) for
the first two results of this section. Put
\(\eta=L^{-k}\) and \(t_j=L^j\eta\) for \(0\leq j\leq k\).
An input \(A\) gives a raw fine increment \(a_0=\eta A\). The
coarse variables \(a_j\) in \eqref{eq:gauge-recursion} are raw
increments on the level-\(j\) lattice. This normalization lets the
linear amplification of a straight segment, of size \(L\), appear
explicitly at each step.

\subsection{The perturbation ball}

\begin{proposition}[Depth-uniform iteration]\label{prop:avg-iteration}
On a torus admitting \(k\) steps, let \(A,B\) be complex Lie
algebra fields on the fine positive links. Let \(W\) be real with
\(\sup_p\|W(\partial p)-I\|\leq\alpha\eta^2\),
and let \(W'=e^{\eta B}W\) with \(\|B\|_\infty\leq\rho\).
There are \(c_L>0\) and \(C_L<\infty\), uniform in the depth and
volume, such that the following holds if
\(\alpha,\rho,\|A\|_\infty<c_L\). Use \(W'\) as the reference in
\eqref{eq:gauge-recursion}, and write \(W'_j=\mathcal M_{\rm an}^j(W')\).
The iterates and factorization are admitted on the regular analytic
charts, with holomorphic dependence on \(B,A\). They satisfy
\begin{equation}
 \|a_j(A)\|_\infty\leq 2t_j\|A\|_\infty,\qquad
 \|v_j(A)-I\|_\infty\leq C_L t_j\|A\|_\infty .
 \label{eq:avg-raw-radii}
\end{equation}
Define \(K_j(A)=D_Aa_j(A)\) and let
\(Q_j^{\rm st}=K_j(0)\). With
\(\Delta=\alpha+\rho+\|A\|_\infty\), the counting-coordinate bounds are
\begin{equation}
 \begin{split}
 \|K_j(A)\|_{\rm row}&\leq t_j(1+C_L\Delta),\\
 \|K_j(A)\|_{\rm col}&\leq t_j L^{-4j}(1+C_L\Delta),\\
 \sup_b\sum_{c,b'}\|D_{A(b')}K_j(A)(c,b)\|
        &\leq C_L t_j^2 L^{-4j}.
 \end{split}
 \label{eq:avg-derivative}
\end{equation}
Each coefficient is supported in the backward double block of its
output link. In particular the column norm of
\(K_j(A)-K_j(\widetilde A)\) is bounded by
\(C_L t_j^2 L^{-4j}\|A-\widetilde A\|_\infty\) on the ball.

The plain log coordinate
\[
 p_j(A)(c)=\log\big[\mathcal M_{\rm an}^j(e^{\eta A}W')(c)W'_j(c)^{-1}\big]
\]
has derivative \(Q_j^{\rm pl}\) at zero. Both Taylor remainders satisfy
\begin{equation}
 \|a_j(A)-Q_j^{\rm st}A\|_\infty+
 \|p_j(A)-Q_j^{\rm pl}A\|_\infty
 \leq C_L t_j^2\|A\|_\infty^2.
 \label{eq:avg-remainders}
\end{equation}
The derivative of the straightened coordinate is exactly
\(Q_j^{\rm st}=\eta J_{W'_{j-1}}\cdots J_{W'_0}\),
where \(J_V=D_aF_V(0)\) as in \cref{prop:gauge-one-step}.
At \(B=0\) its difference from
\(\eta S_{W_{j-1}}\cdots S_{W_0}\), \(W_j=\mathcal M^j W\), has
row norm at most \(C_L\alpha t_j\) and column norm at most
\(C_L\alpha t_j L^{-4j}\). For complex \(B\) the same comparison
with this real-reference product has \(\alpha+\rho\) in place
of \(\alpha\).
\end{proposition}

\begin{proof}
Write \(W_j=\mathcal M^j W\). By \cref{cor:avg-real-iterate},
after decreasing \(c_L\), the real reference plaquette defects
are bounded by \(C_L\alpha t_j^2\) throughout the iteration.
The real reference uses the first branch at every step, so
\(W_j=\mathcal M_{\rm an}^j(W)\) as well. First take \(B=0\). By \cref{prop:gauge-one-step}, each
\(D_aF_{W_s}(a)\) has row and column sums at most \(L\) and
\(L^{-3}\), respectively, multiplied by
\(1+C_L(\alpha t_s^2+\|a\|_\infty)\). The estimates are valid
on a fixed one-step ball. In particular, setting
\(r_s=\|a_s\|_\infty/t_s\) and integrating the derivative along
the segment from zero to \(a_s\) gives
\[
 r_{s+1}\leq r_s\bigl(1+C_L(\alpha t_s^2+t_s r_s)\bigr).
\]
The geometric sums satisfy
\(\sum_{s<j} t_s\leq t_j/(L-1)\) and
\(\sum_{s<j}t_s^2\leq t_j^2/(L^2-1)\).
As long as \(r_s\leq 2\|A\|_\infty\), the product of the
displayed factors is bounded by
\(\exp(C_L(\alpha+2\|A\|_\infty))<2\).
Since \(r_0=\|A\|_\infty\), induction gives the first
bound in \eqref{eq:avg-raw-radii}. The one-step site log
bound \(\|\log u_{W_s}(a_s)\|_\infty\leq C_L\|a_s\|_\infty\)
and the geometric sum bound the accumulated site factors by
\(\exp(C_Lt_j\|A\|_\infty)-1\) in distance from \(I\).

The one-step contour families, mean differences and outer logs
lie inside their charts with a fixed reserve. This follows from
the local real comb coordinates in the proof of
\cref{prop:gauge-one-step} by decreasing \(c_L\).
The accumulated factors have condition numbers at most
\(\exp(C_Lt_j\|A\|_\infty)\), so the reserve admits their
endpoint conjugations in \cref{prop:avg-map}. Thus
\eqref{eq:gauge-factorization} holds on the entire ball, including
complex increments. If the increment is real, the same reserve
can be chosen so that the plaquette defect at each input to
\(\mathcal M\) is below the real branch cutoff. Hence the
holomorphic formulas agree with the real average there.

To treat \(B\ne0\), first apply this argument to the increment
\(B\) about \(W\). The exact factorization gives
\(W'_s=e^{b_s}W_s\), where
\(b_s=\log(W'_sW_s^{-1})\) satisfies
\(\|b_s\|_\infty\leq C_L\rho t_s\). Indeed the ratio is a
product of \(e^{a_s(B)}\), one site factor, and the inverse
site factor transported by the real \(W_s\).
Apply the shifted-reference clause of \cref{prop:gauge-one-step}
to \(W'_s\). This adds \(C_L\rho t_s\) to the distortion factor
at step \(s\). The same summation, with this term included,
proves \eqref{eq:avg-raw-radii} and chart admission for \(W'\).
When \(A,B\) are real, the same one-step cutoff estimate applies
with \(a'=b_s\) and \(a=a_s\). The accumulated real gauges preserve
the test defects, so the iterates equal the Borel iterates there.
This step uses the size of the reference link shift, which may
change the fine curvature scale.

For the derivatives, put \(E_s=D_aF_{W'_s}(a_s(A))\) and
\(P_{v,u}=E_{v-1}\cdots E_u\) for \(0\leq u\leq v\leq k\),
with the empty product equal to the identity. The row and column
sums of a product are bounded by products of the corresponding
sums. The total distortion of these products is at most
\[
 \exp\!\left(C_L\sum_{u\leq s<v}
   [\alpha t_s^2+(\rho+\|A\|_\infty)t_s]\right)
 \leq 1+C_L\Delta .
\]
The chain rule gives \(K_j(A)=\eta P_{j,0}\), which proves the
first two bounds in \eqref{eq:avg-derivative}. At zero all the
\(a_s\) vanish. Replace each factor \(J_{W'_s}\) by
\(S_{W_s}\) in turn. Its error, in both norms, is bounded by
\(C_L(\alpha t_s^2+\rho t_s)\) by the one-step comparisons
in a real comb gauge. The preceding product bounds on the other
factors give the stated errors, since these radii sum to
\(O_L(\alpha+\rho)\).

Differentiating the recursion twice gives
\begin{equation}
 D_A^2a_j(A)
  =\eta^2\sum_{s<j}P_{j,s+1}\,
     D_a^2F_{W'_s}(a_s)[P_{s,0}\,\cdot,P_{s,0}\,\cdot].
 \label{eq:avg-second-chain}
\end{equation}
Fix the fine link of one input leg and sum the norms over the
other input and the output. Use a column sum of \(P_{s,0}\)
on the fixed input and a row sum of \(P_{s,0}\) on the
summed input. The one-step Hessian bound of
\cref{prop:gauge-one-step} sums its output and its second
input with its first input fixed. Finally sum the output of
\(P_{j,s+1}\) by its column bound. The total is at most
\[
 C_L\eta^2\sum_{s<j}L^{-3(j-s-1)}L^{-3s}L^s
 \leq C_L t_j^2 L^{-4j}.
\]
Integrating this coefficient bound along the line between
the inputs gives the column Lipschitz estimate.
Each factor and one-step Hessian reads only the relevant double
block, so the products and derivatives have the asserted support.

Using row bounds for all factors of \eqref{eq:avg-second-chain}
gives
\[
 \|D_A^2a_j(A)[X,Y]\|_\infty
 \leq C_L\eta^2\sum_{s<j} L^{j-s-1} L^{2s}
                         \|X\|_\infty\|Y\|_\infty
 \leq C_L t_j^2\|X\|_\infty\|Y\|_\infty .
\]
Taylor's formula proves the straightened remainder.
The quadratic part is the sum in \eqref{eq:avg-second-chain}
evaluated at zero with both input legs \(A\), divided by two.
At final depth \(k\), its term from scale \(s\) is bounded by
\(C_L t_s\|A\|_\infty^2\). If the reference used in the recursion
is \(W'=I\), this term vanishes when every link increment is a
multiple of one fixed \(X\in\mathfrak g\). All path factors then
commute, so each one-step logarithm is exactly linear on that subspace.

The derivative of \(\log v_j\) at zero, denoted by \(\dot v_j\),
is the sum of the linear parts of the successive block factors
restricted to level \(j\). The Taylor remainder of the logarithm
of the factor from level \(s\) is bounded by
\(C_L t_s^2\|A\|_\infty^2\), using the straightened remainder
at level \(s\) and the analytic one-step site map. For products
of exponentials with total log norm at most a fixed small
number, the log of the product differs from the sum by at
most a constant times the square of that total norm. This
follows by expanding the product and the logarithm and
summing terms of degree at least two. Consequently
\(\|\log v_j-\dot v_j\|_\infty\leq C_Lt_j^2\|A\|_\infty^2\).
By \eqref{eq:gauge-factorization},
\[
 e^{p_j(A)(c)}
 =v_j(c_-) e^{a_j(A)(c)}
      \bigl[W'_j(c)v_j(c_+)^{-1}W'_j(c)^{-1}\bigr].
\]
The references \(W'_j\) have uniformly bounded condition numbers,
as shown by their shifts from \(W_j\).
The product log estimate gives the plain remainder and the exact
identity
\begin{equation}
 (Q_j^{\rm pl}A)(c)
  =(Q_j^{\rm st}A)(c)+\dot v_j(A)(c_-)
               -\operatorname{Ad}(W'_j(c))\dot v_j(A)(c_+).
 \label{eq:avg-plain-st}
\end{equation}
The last two terms can be nonzero even at a flat reference.
\end{proof}

The plain coordinate also has a uniform first derivative:
$\|D_Ap_j(A)\|_{\infty\to\infty}\le C_Lt_j$ on a smaller
fixed input ball. The factorization and the site-factor bounds
give $\|p_j(A)\|_\infty\le C_Lt_j$ on a larger such ball;
Cauchy's formula along a unit input direction proves the assertion.
Dividing by the output length and converting the input to physical
units gives a bounded per-unit-length derivative.

There is also a useful positive majorant. Let $S_j^+$ be the
per-unit-length straight row with every real parallel transport
replaced by its norm one. Let $D(c)=B^j(c_-)\cup B^j(c_+)$ be
the backward double block of the output row, and put
$H_j|A|(c)=L^{-4j}\sum_{b\subset D(c)}|A(b)|$.
For the real reference of Proposition~\ref{prop:avg-iteration},
write $\widehat Q_j^{\rm st}=t_j^{-1}Q_j^{\rm st}$ for the
derivative with both coordinates per unit length. It obeys
\begin{equation}\label{eq:avg-positive-majorant}
 |\widehat Q_j^{\rm st}(W)A(c)|
       \le S_j^+|A|(c)+C_L\alpha H_j|A|(c).
\end{equation}
The comparison with the product of the $S_{W_i}$ has each entry
bounded by $C_L\alpha L^{-4j}$, by its column estimate, and
support in $D(c)$. The product itself is bounded by $S_j^+$ by
unitarity and the triangle inequality. These facts give the display.

For a prescribed external complex gauge $g$, take the real reference
$W$ (thus $B=0$ in Proposition~\ref{prop:avg-iteration}). In its
final-scale units the factorization becomes
\begin{equation}\label{eq:avg-external-gauge}
 \mathcal M_{\rm an}^j\bigl((e^{\eta A}W)^g\bigr)
       =(e^{a_j}W_j)^{(g|_{\mathrm{level}\ j})v_j}.
\end{equation}
Besides the small-field hypotheses, require the contour and logarithm
charts and their conjugates to be admitted. A sufficient condition on
the smaller balls used above is
\[
 \max_x\operatorname{cn}(g(x))\max_{i,x}\operatorname{cn}(v_i(x))\le2,
 \qquad \operatorname{cn}(g)=\|g\|\|g^{-1}\|.
\]
The proof of chart admission has a fixed reserve, so decreasing the
input ceiling accommodates this factor two at every level. Endpoint
covariance and \eqref{eq:gauge-factorization} then give the identity.

\subsection{Curvature and open transports}

For the real reference in \cref{prop:avg-iteration}, let \(d_j\) be
the supremum of the distance between \(W_j(c)\) and the straight
bare transport between the endpoints of \(c\), over all level-\(j\)
links. The bound on the loop correction in the proof of
\cref{prop:avg-map} and unitary telescoping give
\[
 d_{j+1}\leq Ld_j+C_L\alpha t_j^2,\qquad d_0=0.
\]
Hence \(d_j\leq C_L\alpha t_j^2\). In this comparison the error
at the finer levels is multiplied only by the number of links
in a straight segment.

\begin{proposition}[Complex curvature at its own scale]\label{prop:avg-complex-curvature}
Under the assumptions of \cref{prop:avg-iteration}, assume also
\(\sup_p\|W'(\partial p)-I\|\leq\alpha'\eta^2\) with
\(0\leq\alpha'\leq c_L\). After a fixed decrease of \(c_L\),
\[
 \|\mathcal M_{\rm an}^j(W')(\partial C)-I\|
         \leq 2\alpha't_j^2
\]
for each level-\(j\) plaquette \(C\).
For an individual plaquette, the additional curvature bound on
\(W'\) is needed only on its complete backward blocks. The real
reference hypotheses of \cref{prop:avg-iteration} remain in force.
\end{proposition}

\begin{proof}
The result at \(j=0\) is the assumption. For \(j>0\) take the
box \(Y\) formed by the four order-\(j\) blocks at the corners of
\(C\). Work on a lift if necessary. Put \(n=L^j\), \(r=\alpha't_j^2\).
Use the comb gauge \(g\) in \(Y\) formed with the complex field
\(W'\). Each path and inverse path of at most \(C n\) bare links
has norm at most \(e^{C\rho t_j}\). This follows by multiplying
the bounds \(e^{\eta\rho}\) on complex increments or their
inverses and using unitarity of the real reference. We take
\(K=e^{C\rho t_j}\), increasing the absolute constant to bound
all condition numbers in the comb estimate and those of \(g\).
With \(W'^g=e^a\) on \(Y\), \cref{lem:avg-comb} gives
\(\|a\|_\infty\leq C_L r/n\). For inverse plaquette factors
one may use \(\|P^{-1}-I\|=\|P-I\|\) for \(P\in G_{\mathbb C}\):
\(P^{-1}-I\) is the adjugate of \(P-I\), whose singular
values are those of \(P-I\) in reverse order.

Extend \(a\) by zero outside \(Y\) on a torus on which the read
embeds. This permits use of \cref{prop:avg-iteration} at the flat
reference with final scale \(j\), since the normalized input
has norm \(n\|a\|_\infty\leq C_L r\).
The four output links read \(a\) only in \(Y\). The extension
serves to apply the flat supremum-ball estimate.
Let \(P_j(a)=\log\mathcal M_{\rm an}^j(e^a)\) and let \(Q^{(j)}\) be
the composition of the \(j\) raw flat link differentials
\eqref{eq:avg-flat-link}. Then
\[
 \|P_j(a)\|_\infty\leq C_L r,\qquad
 \|P_j(a)-Q^{(j)}a\|_\infty\leq C_L r^2 .
\]
At intermediate levels \(s\leq j\) the corresponding plain
logarithms have size at most \(C_L r L^{s-j}\).
The contours on their small branch have a fixed reserve.
Covariance under \(g^{-1}\) is therefore admitted at each
step by \cref{prop:avg-map}, since the conjugations cost at
most \(K\). This proves the relation between the four output
links of the gauged field and those of \(W'\) on the read.

On fine plaquettes in \(Y\), product expansion gives
\[
 \|\mathrm da(p)\|\leq K\alpha'\eta^2+C_L\|a\|_\infty^2.
\]
Apply the iterated version of \eqref{eq:avg-flat-curl} at \(C\).
It reads only plaquettes in \(Y\), has nonnegative weights,
and its total weight is \(n^2\). The coarse plaquette product
differs from its linear term by \(O_L(r^2)\). The displayed
Taylor bound for the four coarse links and \(n\|a\|_\infty
\leq C_L r\) therefore give
\[
 \|\mathcal M_{\rm an}^j(e^a)(\partial C)-I\|\leq Kr+C_L r^2.
\]
Conjugating back gives at most \(K^2r+C_LKr^2\leq2r\)
for the stated ceilings. The path and flux identities on a
cover descend with the same directed links and multiplicities.
\end{proof}

\subsection{Rows on a layered region}

We return to physical spacings \(h_j=L^j h_0\). For a level-\(j\)
link \(c\), let \(D(c)=B^j(c_-)\cup B^j(c_+)\). A straightened
row uses the raw fine perturbation \(h_0A\), and has output
\(q_c(A)=a_j(c)/h_j\), per unit length. Fix a family of rows
without repetitions and assign a length \(\xi(b)\) to each fine
link in their carriers. Assume that \(h_j/L\leq\xi(b)\leq Lh_j\)
for \(b\subset D(c)\). For rows on adjacent layers this allows the
smaller spacing on the part of the carrier outside the inner
region. A fixed fine link can then enter rows of at most three
levels, since their spacings lie in \([\xi(b)/L,L\xi(b)]\).
At each level any such row is incident to the unique block
containing \(b_-\). The degree bound gives bounded overlap.

Assume on each \(D(c)\) the hypotheses of
\cref{prop:avg-iteration} with final spacing \(h_j\). In physical
units these include a real plaquette defect bounded by
\(\alpha(h_0/h_j)^2\) and a bound \(h_j\|A\|_{\infty,D(c)}<c_L\).
With a complex reference shift \(h_0B\), impose
\(h_j\|B\|_{\infty,D(c)}<c_L\) there as well.
The proof of \cref{prop:avg-iteration} applies to the row on
this stock. Here \(D(c)\) in the hypotheses includes every fine
link and plaquette whose vertices lie in the projected union.
On a torus with two sites in the row direction at level \(j\),
the real branch test can include plaquettes at the other periodic
face of this union. The curvature hypothesis includes those
plaquettes. The analytic derivatives use just the lifted path
formula. Every intermediate double block in that formula lies in
the carrier. The real curvature recurrence on four endpoint
blocks bounds the plaquettes needed at each subsequent step,
including those tested across a periodic face.

With \(A\) measured per length, the coefficients satisfy
\begin{equation}
 \begin{split}
 \sum_b\|D_bq_c(A)\|&\leq 2,\\
 \|D_bq_c(A)\|&\leq 2(h_0/h_j)^4,\\
 \sup_b\sum_{b'}\|D_bD_{b'}q_c(A)\|
             &\leq C_L h_j(h_0/h_j)^4 .
 \end{split}
 \label{eq:avg-layered-coeffs}
\end{equation}
Indeed use \(h_j A\) for the normalized input of
\cref{prop:avg-iteration} with final depth \(j\) and divide
the output by \(h_j\). Each derivative multiplies by
\(h_j\), giving the factor \(h_j\) for the second derivative.
The fine-link support of both legs is \(D(c)\).
For example, the raw row \(a_j(c)\) has weighted row sum
\(\sum_b\|D_ba_j(c)\|/\xi(b)\leq 2L\).
The sum of its derivatives over all the chosen rows at fixed
input link is bounded by \(C_L h_0^4/\xi(b)^3\).
The single adjacent-scale comparison in the first bound costs
at most \(L\), independent of how many levels enter the row.

The same density controls a weighted squared-sum estimate.
All fine sums below run over the union of the row carriers,
where the positive length $\xi$ is assigned.
For any fine input \(X\), Cauchy--Schwarz in one row and then
the carrier overlap give
\begin{equation}
 \sum_c h_j^4 h_j^{-2}|Dq_c(A)X|^2
 \leq C_L\sum_b h_0^4\xi(b)^{-2}|X(b)|^2,
 \label{eq:avg-layered-trace}
\end{equation}
where \(j\) on the left is the level of \(c\).
In fact the factor multiplying \(|X(b)|^2\) from one row is
bounded by \(C_Lh_0^4 h_j^{-2}\) by
\eqref{eq:avg-layered-coeffs}, and there are only boundedly
many rows containing \(b\).

For an assigned row current \(\omega(c)\) satisfying
\(|\omega(c)|\leq\gamma h_j^{-3}\), the second derivative
also satisfies
\begin{equation}\label{eq:avg-layered-second}
 \left|\sum_c h_j^4
     \langle\omega(c),D^2q_c(A)[X,Y]\rangle\right|
 \leq C_L\gamma
 \left(\sum_b h_0^4\xi(b)^{-2}|X(b)|^2\right)^{1/2}
 \left(\sum_b h_0^4\xi(b)^{-2}|Y(b)|^2\right)^{1/2}.
\end{equation}
The pairing is the bilinear extension of the real colour pairing.
Fix the first fine link and sum the absolute coefficients over
the second link and the rows containing the first one. The current
and row measure cost at most \(\gamma h_j\). Multiplying by the
last bound in \eqref{eq:avg-layered-coeffs} gives
\(C_L\gamma h_0^4 h_j^{-2}\) for each carrier. The scale comparison
and overlap bound replace this by \(C_L\gamma h_0^4\xi(b)^{-2}\)
after summation. Symmetry of the second derivative gives the
same estimate with the input links exchanged. Weighted
Cauchy--Schwarz then gives \eqref{eq:avg-layered-second}.

% End average-estimates.tex
% Begin rows.tex
\subsection{Determining rows in separated regions}\label{sec:layered-geometry}

Let $\Omega_0\supset\cdots\supset\Omega_k$ be regions on the
ambient periodic lattice and put $\Omega_j=\varnothing$ for $j>k$. Use
half-open lattice cells for boundary membership. Suppose each
$\Omega_j$ is a union of centered depth-$j$ blocks and
$\Omega_{j+1}$ is a union of centered depth-$(j+1)$ blocks in the
same hierarchy. In particular, the hypothesis holds for regions
formed from aligned cubes of side $M_1 h_j$, with $M_1$ an odd
power of $L$. Impose the separation
\begin{equation}\label{eq:layered-separation}
 \operatorname{dist}_\infty(\Omega_j^c,\Omega_{j+1})>\rho_{\rm gap} h_j
 \qquad (0\le j<k),\qquad \rho_{\rm gap}\ge 4L.
\end{equation}
The distance is infinite when one of the sets is empty. For a
level-$j$ site $y\in\Omega_j\setminus\Omega_{j+1}$ call $y$ a
top of level $j$ and $B^j(y)$ its pillar. These pillars partition
the fine sites of $\Omega_0$. Indeed, a fine site in layer $j$
belongs to a depth-$j$ block contained in $\Omega_j$ and disjoint
from $\Omega_{j+1}$, by alignment. Its centre is a level-$j$ top.

Use the determining rows
\begin{equation}\label{eq:layered-rows}
 \Gamma_j=\{c:\ c\text{ is a level-}j\text{ link with both ends in }
 \Omega_0,\ \text{one end in }\Omega_j,\ \text{and no end in }
 \Omega_{j+1}\}.
\end{equation}
If both ends are in $\Omega_j$, the row is interior. If only one
end is in $\Omega_j$, the row is outward. We use one positive
orientation for each link. A link from layer $j$ into
$\Omega_{j+1}$ is an inward link and is excluded from $\Gamma_j$.
The scale at a fine site is the spacing of the level of its pillar.
For a fine link or plaquette $e$, let $\xi(e)$ be the largest of
these spacings at its endpoints or corners.

\begin{lemma}[Available rows and their carriers]\label{lem:layered-carriers}
Let \(\Omega_0\supset\cdots\supset\Omega_k\) be nested regions,
each a union of aligned depth-\(j\) blocks at its respective level,
satisfying \eqref{eq:layered-separation}. For the determining rows
\(\Gamma_j\) of \eqref{eq:layered-rows}, the endpoints of every row
$c\in\Gamma_j$ are tops of level $j$ or $j-1$, with only level
zero allowed when $j=0$. If $j\ge1$, the double block
$B^j(c_-)\cup B^j(c_+)$ is contained in $\Omega_{j-1}$ and
disjoint from $\Omega_{j+1}$. For an outward row the block at the
outer end is entirely in layer $j-1$. All level-$(j-1)$ links of
that outer block's rectangular graph are interior determining rows
at that level.
Each top of positive level has an incident determining row in
every coordinate direction, allowing either sign. The row carriers,
also with a collar of one fine cell intersected with $\Omega_0$,
have bounded overlap independent of $k$ and the volume. The carrier
scale is comparable to all its fine site scales within a factor $L$.
\end{lemma}

\begin{proof}
An interior endpoint is in layer $j$ by \eqref{eq:layered-rows}.
For $j\ge1$, the other endpoint of an outward row is at distance
at most $h_j$ from $\Omega_j$. Its block and its one-cell collar are within
$2h_j$ of $\Omega_j$, so \eqref{eq:layered-separation} places
them in $\Omega_{j-1}$. Alignment makes the entire outer block
disjoint from $\Omega_j$. This proves the claim about its rows.
The two endpoint blocks miss $\Omega_{j+1}$ by alignment and the
absence of inward rows. The same distance estimate places both
blocks in $\Omega_{j-1}$. A one-cell collar can enter
$\Omega_{j+1}$ but cannot reach $\Omega_{j+2}$. If $j+2\le k$, this
follows from \eqref{eq:layered-separation} with index $j+1$;
otherwise $\Omega_{j+2}$ is empty.
At level zero, rows lie in $\Omega_0$ and their collar is considered
only inside $\Omega_0$. For the same reason that collar cannot reach
$\Omega_2$. This proves scale comparability on the stated carriers.

A top $y$ of positive level $j$ is outside $\Omega_{j+1}$.
If $j<k$, its unique aligned
depth-$(j+1)$ cell misses $\Omega_{j+1}$ and has $L\ge3$ level-$j$ sites in each
coordinate direction. Along the line through $y$ in that cell, at
least one of the two adjacent sites is still in the cell. At the
highest level $j=k$, use either neighbour without a further tiling.
The chosen neighbour is in $\Omega_{j-1}$ by separation, so the incident row
belongs to \eqref{eq:layered-rows}. For each fixed $j$, the row carriers
have diameter $O(h_j)$, their origins lie on the level-$j$ grid, and
there are at most eight incident oriented choices per origin. They
therefore have bounded overlap at that level. A fine site belongs to
carriers at at most three levels by scale comparability. Periodic
counting uses the rectangular lifts, with bounded multiplicity
under the standing compatibility condition.
\end{proof}

Consider a level-$j$ mixed plaquette, $j\ge1$, with a corner in
$\Omega_j$. Suppose its sides are retained rows or links with both
ends outside $\Omega_j$ that are completed by one averaging step.
A completed side reads the path links and projected branch-test
plaquettes in the union of its two endpoint blocks at the preceding
lattice level, including the links joining those blocks.
Each endpoint is within $h_j$ in
maximum distance of the corner in $\Omega_j$, so the full blocks are
within $2h_j$ of that region. Separation places them in $\Omega_{j-1}$
and alignment places them outside $\Omega_j$. The completed side is
therefore the actual average of the recorded finer-layer stock, and
transforms at its two endpoints. The mixed plaquette transforms by
conjugation at its base.
For real fields, changing any of its four link transports changes
the plaquette by at most the sum of their changes, by unitary
telescoping. In particular the one-step derivative estimates can be
used for a completed side on its full admitted stock.

\subsection{Two-field comparison after axial alignment}\label{sec:layered-rigidity}

We use the real curvature bound of Section~\ref{sec:energy} in its
local form. To bound an individual level-$n$ plaquette it suffices
to bound the bare plaquettes in the four complete backward blocks
at its corners. At one step these blocks contain every path and
every real branch test used in the four output links. They contain
the flux support in \eqref{eq:energy-flux} as well. Applying the
curvature estimate on these stocks and then proceeding through the
backward blocks gives the same curvature recurrence. The local
claim includes the tested periodic plaquettes whenever the projected
stock has an additional face. Equivalently, the recurrence can be
applied on lifts once its branch tests have been checked there.

\begin{proposition}[Depth-uniform raw rigidity]\label{prop:layered-rigidity}
Let $\Omega_0\supset\cdots\supset\Omega_k$ be nested regions,
each a union of aligned depth-$j$ blocks at its respective level,
satisfying \eqref{eq:layered-separation}. Use the determining rows
\eqref{eq:layered-rows} and the associated pillar spacings $\xi$.
Let $U$ and $V$ be real fields
whose bare plaquettes obey
\begin{equation}\label{eq:layered-rigidity-curvature}
 \max\big\{|U(\partial p)-I|,|V(\partial p)-I|\big\}
       \le \alpha(h_0/\xi(p))^2
\end{equation}
for every physical fine plaquette whose four vertices lie in
\(D(c)=B^j(c_-)\cup B^j(c_+)\) for some \(c\in\Gamma_j\),
\(1\le j\le k\). This includes opposite-face plaquettes at a
short periodic seam. Assume \(0\le\alpha\le c_L\) for a sufficiently
small positive constant depending only on \(L\).
Let $d\ge0$ and suppose the raw difference of their averages on each
$c\in\Gamma_j$, $0\le j\le k$, is at most $d$. There is a real
gauge $g$ equal to $I$ at all tops such that
\begin{equation}\label{eq:layered-rigidity-result}
 \big|\mathcal M^n(U^g)(b)-\mathcal M^n(V)(b)\big|
       \le C_L(d+\alpha)
\end{equation}
for all level-$n$ links $b$, \(0\le n\le k\), with both endpoints in $\Omega_0$ and
at least one in $\Omega_n$. The constants are independent of the
depth and volume. The inequality concerns link matrices in the
aligned frame, before dividing by a lattice length.
\end{proposition}

\begin{proof}
Write $U_n=\mathcal M^n(U)$ and $V_n=\mathcal M^n(V)$.
By Lemma~\ref{lem:layered-carriers}, every positive-depth pillar is
contained in a carrier at its top level. Its bare plaquette bound
admits the lower rows and supplies small curvature for each contour
family by Corollary~\ref{cor:avg-local-real-iterate}. Thus the means
in the descent exist. Starting at the tops, set $g(y)=I$ there
and descend to finer centres. In the block
of a level-$(n+1)$ centre $y\in\Omega_{n+1}$ set
\begin{equation}\label{eq:layered-rigidity-gauge}
 g(x)=T_{V_n}(y,x)^{-1}g(y)T_{U_n}(y,x).
\end{equation}
At the centre this retains the value already assigned. The blocks
of such centres and the level-$n$ tops partition the level-$n$
sites of $\Omega_n$. The recursion assigns a single real fine gauge,
which can be extended by the identity outside $\Omega_0$. Covariance
proves that the two block contours agree in every block where the
recursion was used. The averages on determining rows retain their
original difference since both endpoints are tops.

For a carried level-$n$ link $b$ let $j$ be the larger of its endpoint
pillar levels. If \(n=j\), the link or its inverse is already a
determining row, so the asserted comparison holds. Assume \(n<j\).

For \(j\ge1\), the depth-\(j\) ancestor of the inner endpoint is its
level-\(j\) top. The other ancestor is within \(h_n+h_j/2\) of
\(\Omega_j\), hence lies in \(\Omega_{j-1}\) by separation. Both
ancestors are outside \(\Omega_{j+1}\) by alignment. If distinct,
they form a directed level-\(j\) determining row. If they coincide,
their common pillar is contained in a level-\(j\) determining
carrier by Lemma~\ref{lem:layered-carriers}. Each block at an
intermediate comparison is contained in this carrier. The same is
true after projection of the complete backward blocks for its
curvature reads and tests.

The carrier is contained in \(\Omega_{j-1}\), so each of its bare
plaquette logarithms has norm at most \(2\alpha L^2L^{-2j}\), by
\eqref{eq:layered-rigidity-curvature} and the real small-log bound.
Apply Corollary~\ref{cor:avg-local-real-iterate} on the carrier with
parameter \(2L^2\alpha\). The level-\(n\) group defects in the
comparison are at most \(4L^2\alpha L^{2(n-j)}\). This uses only a
fixed decrease of the ceiling for \(\alpha\), depending on \(L\).

For $n<j-1$, both endpoints are in $\Omega_{n+1}$.
Indeed the endpoint in $\Omega_j$ has distance greater than
$\rho_{\rm gap} h_{j-1}$ from $\Omega_{j-1}^c$ and the link has length at most
$h_{j-1}$. The comparison in the next paragraph also applies at
$n=j-1$ if both endpoints are in $\Omega_j$.

If the endpoints belong to the same depth-$(n+1)$ block,
their aligned contour values determine the link up to curvature by
Lemma~\ref{lem:gauge-rigidity}, using just that block. Otherwise their
parent centres are endpoints of a level-$(n+1)$ row. The difference
of that row bounds the difference of $b$ with coefficient one,
again by Lemma~\ref{lem:gauge-rigidity}, because the two endpoint
contour differences vanish.


If the descent comparison reaches an inward level-$(j-1)$ link,
its parent row is an outward member of $\Gamma_j$. The inner
endpoint contour has been aligned. The block at the outer endpoint
of that row lies in layer $j-1$, by Lemma~\ref{lem:layered-carriers}.
Each of its internal paths has length at most $2(L-1)$ and uses
only its determining rows. Comparing these path products in the two
fields costs at most $2(L-1)d$. Replacing a comb transport by its
mean in either field costs at most $C_L\alpha$ by the small-curvature
path comparison. The outer contour discrepancy is therefore at
most $2(L-1)d+C_L\alpha$. The outward row itself contributes $d$.
The proof of Lemma~\ref{lem:gauge-rigidity}, using the separate
contour errors and the aligned inner endpoint, bounds the inward discrepancy by
$(2L-1)d+C_L\alpha$.

Thus the last comparison terminates either within a block, on a
determining row, or on this inward link. Each preceding comparison
adds only $C_L\alpha L^{2(n-j)}$ at its own level, since the coefficient
of the parent row error is one. The sum over those levels is bounded
by $C_L\alpha/(L^2-1)$, which proves
\eqref{eq:layered-rigidity-result}.
\end{proof}

\subsection{Curvature and frame changes on a fixed box}\label{sec:box-frames}

For a coordinate comb with order $4,3,2,1$ rooted at $y$, set
$s_U(x)=U(\Gamma_{y,x})$. Use fine spacing one for the next estimate.
In this frame write $\widetilde F_{\rho\mu}(z)$ for the logarithm of
the positively oriented $\rho\mu$ plaquette at $z$, transported to
$y$ by $s_U(z)$, with $\rho<\mu$. Suppose $U$ is real and the
plaquette defect is at most $\epsilon<1/2$ on the box. Assume also that
for some $0<\gamma\le1$ and $H\ge0$,
\[
 |\widetilde F_{\rho\mu}(z)-\widetilde F_{\rho\mu}(y)|
       \le H|z-y|_1^\gamma
\]
for the plaquettes in the comparison strips, with the base plaquette
present in each orientation. For $r=|x-y|_1$ and $r\epsilon$ small,
\begin{equation}\label{eq:box-comb-first-jet}
 \left|\log U^{s_U}(x,\mu)-
          \sum_{\rho<\mu}(x_\rho-y_\rho)\widetilde F_{\rho\mu}(y)\right|
    \le C\bigl(Hr^{1+\gamma}+r^2\epsilon^2\bigr).
\end{equation}
To prove the estimate, append the positive $\mu$ link to the comb
and move it left past each step in a direction $\rho<\mu$. The
prefix before a swap is the original comb to that point. For a
negative $\rho$ step the resulting inverse plaquette has its
positive base one step farther along that same prefix. Each signed
swap therefore contributes the corresponding
$\pm\widetilde F_{\rho\mu}(z)$ in the root frame. The signed number
in direction $\rho$ is $x_\rho-y_\rho$. The logarithm of the product
differs from the sum of at most $r$ small logarithms by
$C r^2\epsilon^2$, from the exponential and logarithm series on
$r\epsilon$ small. Replacing each transported curvature by its base
value costs at most $Hr^{1+\gamma}$. This proves
\eqref{eq:box-comb-first-jet}.

\begin{proposition}[Bounded-box complex frame]\label{prop:box-complex-frame}
Let $D$ be a rectangular box of side lengths at most $Mh_j$, $M\ge1$, with
root $y$ and a complete stock of bare bonds and plaquettes on an
embedded lift. Let \(\alpha_0,\alpha_1\ge0\), and let \(A'\)
be a complex Lie algebra link field
per unit length. Consider $X=e^{h_0A'}U$, where $U$ is real,
$h_j\|A'\|_{\infty,D}\le\alpha_1$, and both fields have plaquette
defect at most $\alpha_0(h_0/h_j)^2$ on this stock. Suppose
\(M(\alpha_0+\alpha_1)\le c_L\), for a sufficiently small
positive constant depending only on \(L\). Let
$s_X(x)=X(\Gamma_{y,x})$ and $v=s_X^{-1}$ on the level-$j$ sites.
On level-$j$ rows with complete backward reads in $D$, including
projected branch-test plaquettes for real arguments,
\begin{equation}\label{eq:box-complex-frame}
 \mathcal M_{\rm an}^j(X)=V^v,\qquad
 V=\mathcal M_{\rm an}^j(X^{s_X}),\qquad
 \|V-I\|_\infty\le C_L M\alpha_0.
\end{equation}
When the real first branches are admitted throughout the reads,
these identities also hold for the real Borel map.
The factorization $v=v'v_r$ has $v_r=s_U^{-1}$ real and independent
of $A'$, and $v'$ holomorphic with
\[
 \|v'-I\|_\infty\le C M\alpha_1,\qquad
 \sup_x|v(x)|\,|v(x)^{-1}|\le \exp(CM\alpha_1).
\]
\end{proposition}

\begin{proof}
Writing $d=s_Xs_U^{-1}$ expresses $d(x)$ as a product of at most
$4Mh_j/h_0$ link perturbations, conjugated by real prefix transports.
For a reversed link the perturbation is inverted and conjugated by
its real link, with unchanged norm. Consequently
\[
 \|d-I\|_\infty\le e^{CM\alpha_1}-1,
 \qquad \sup_x\max\{|d(x)|,|d(x)^{-1}|\}\le e^{CM\alpha_1}.
\]
The same estimate bounds the condition numbers of all prefixes in
the comb comparison. Lemma~\ref{lem:avg-comb} in that complex frame
gives $|X^{s_X}(b)-I|\le C_LM\alpha_0 h_0/h_j$ on the box. Indeed
its number of plaquette factors is $O(Mh_j/h_0)$ and their defects
are at most $\alpha_0(h_0/h_j)^2$. The small logarithms have the same
bound. Extend these logarithms by zero on a compatible larger torus
to use the flat supremum-ball estimate in
Proposition~\ref{prop:avg-iteration} at final spacing $h_j$. The
complete reads in $D$ are unchanged. Their normalized fine amplitude
is $O_L(M\alpha_0)$, so the output logarithms are $O_L(M\alpha_0)$.
The uniform chart reserve in the proof of that proposition admits
conjugation back by $s_X^{-1}$, because its condition numbers are
bounded as above. This proves \eqref{eq:box-complex-frame}, including
its holomorphic interpretation on the corresponding admitted patch.

Finally $v=s_U^{-1}d^{-1}$, so $v'=s_U^{-1}d^{-1}s_U$. The estimates
for $d$ and $d^{-1}$ prove the asserted bounds. A pure gauge $U=I^u$
with $A'=0$ gives $v(x)=u(x)u(y)^{-1}$. This can be far from $I$,
even though its condition number is one and $v'=I$.
\end{proof}

\section{Gauge identities and site and link rows}\label{sec:gauge-rows}

\subsection{The effective gauge transformation}

Suppose that \(e^aW\) and \((e^aW)^g=e^{a^g}W\) are in the
one-step chart. For complex gauges require the conjugations of
Proposition~\ref{prop:avg-map} to be admitted. With the reference
\(W\) held fixed, define
\begin{equation}\label{eq:gauge-effective}
 r_{W,a}(g)=u_W(a^g)^{-1}\,(g|_{\mathrm{coarse}})\,u_W(a).
\end{equation}
Substitution in \eqref{eq:gauge-straighten} gives
\[
 e^{F_W(a^g)}\mathcal M_{\rm an}(W)
    =\big(e^{F_W(a)}\mathcal M_{\rm an}(W)\big)^{r_{W,a}(g)} .
\]
Suppose \(V=e^aW\) is contour axial relative to \(W\), including
\(a=0\). Put
\[
 C_x=T_W(y,x)g(x)T_W(y,x)^{-1}.
\]
Then \(u_W(a)=I\), and the contour ratio in \(u_W(a^g)\) is
\(g(y)C_x^{-1}\). For a complex gauge, also require the arguments
\(g(y)^{-1}C_x\) of \(R_Wg\) to lie in its logarithm chart.
On these charts, inversion and conjugation of the logarithm give
\[
 -\operatorname{Ad}(g(y)^{-1})\log(g(y)C_x^{-1})
              =\log(g(y)^{-1}C_x).
\]
Substitution in \eqref{eq:gauge-effective} gives
\(r_{W,a}(g)=R_Wg\). If the plain iterates are mean tower axial
relative to the reference, each \(u_{W_i}(a_i)=I\), by
\eqref{eq:gauge-recursion}. At successive stages the same
identification gives \(r_j=R_{W_{j-1}}\cdots R_Wg\), provided
the site logarithms at those stages are admitted.

For \(e^{a_0}W\) and its transform by \(g\), let \(a_j\) and
\(a_j^g\) be the two recursions \eqref{eq:gauge-recursion}, with
accumulated factors \(v_j\) and \(v_j^g\), respectively. Here
\(a_j^g\) denotes the recursion starting from the gauged fine field.
Iterate \eqref{eq:gauge-effective} by setting \(r_0=g\) and
\(r_{j+1}=r_{W_j,a_j}(r_j)\). Induction gives
\begin{equation}\label{eq:gauge-effective-chain}
 e^{a_j^g}W_j=(e^{a_j}W_j)^{r_j},
 \qquad r_j=(v_j^g)^{-1}(g|_{\mathrm{level}\ j})v_j .
\end{equation}
The second identity follows by inserting the two product recursions
for \(v_j\). It specifies the residual transformations of the straight
coordinates: \(r_j=1\) suffices to leave the straightened field fixed,
and at \(a_0=0\), on the admitted site charts, it becomes the
normalization \(R_{W_{j-1}}\cdots R_Wg=1\).
For the plain average, a gauge equal to \(1\) on level \(j\) preserves
its value by covariance.

We record the derivative of \eqref{eq:gauge-effective} at a general
increment. For a Lie algebra element \(z\) put
\[
 E_z b=\int_0^1\operatorname{Ad}(e^{s z})b\,ds .
\]
The left logarithmic derivative of the exponential is \(E_z\), as
follows by differentiating its convergent series. This operator has
a holomorphic inverse near \(z=0\). Write \(V=e^aW\),
\(z_x=\log[T_V(y,x)T_W(y,x)^{-1}]\),
\(m=L^{-4}\sum_{x\in B(y)}z_x\) and \(u=e^m\).
The derivative at \(t=0\) of \(r_{W,a}(e^{t\lambda})(y)\) is
\begin{equation}\label{eq:gauge-tau}
 \begin{split}
 \tau_{W,a}\lambda(y)=\operatorname{Ad}(u^{-1})\Big[
  \lambda(y)-E_m L^{-4}\sum_{x\in B(y)} E_{z_x}^{-1}
  \big(\lambda(y)-\operatorname{Ad}(T_V(y,x))\lambda(x)\big)\Big].
 \end{split}
\end{equation}
Indeed the left derivative of the relative contour ratio is the
expression in the inner parentheses, by covariance of \(T_V\).
Apply \(E_{z_x}^{-1}\), average, and differentiate
\(u_W(a^{e^{t\lambda}})^{-1}e^{t\lambda(y)}u_W(a)\).
At \(a=0\) the centre evaluations cancel and
\begin{equation}\label{eq:gauge-site-one}
 q_W\lambda(y)=\tau_{W,0}\lambda(y)
   =L^{-4}\sum_{x\in B(y)}
        \operatorname{Ad}(T_W(y,x))\lambda(x).
\end{equation}

On the smaller one-step ball in Proposition~\ref{prop:gauge-one-step},
\(\|\tau_{W,a}-q_W\|_{\infty\to\infty}\le C_L\|a\|_\infty\).
This follows from \eqref{eq:gauge-tau}, since \(E_z\),
\(E_z^{-1}\) and \(\operatorname{Ad}(e^z)\) differ from the
identity by \(O(\|z\|)\), and the transport depends holomorphically
on the finite link increments. For a complex reference shifted
from a real one by a raw increment of size \(\delta\),
\(\|q_W\|_{\infty\to\infty}\le 1+C_L\delta\) by the same
one-step transport comparison. Hence the effective site derivative
of the \(j\)-step transformation is
\(\tau_{W_{j-1},a_{j-1}}\cdots\tau_{W,a_0}\).
On the depth-uniform balls of Section~\ref{sec:avg-iteration},
the raw increments and the changes in the reference are summable
over the levels. The product is therefore bounded uniformly, and
telescoping against \(q_{W_{j-1}}\cdots q_W\) bounds the difference
by \(C_L\sum_{r<j}\|a_r\|_\infty\). In particular, the effective
derivative in a fixed chart generally depends on its increment.

For a field \(V\) on links of length \(h\), use the covariant
difference with the sign fixed by \(g=e^{t\lambda}\),
\begin{equation}\label{eq:gauge-D}
 (D_V\lambda)(b)=h^{-1}
     \big[\lambda(b_-)-\operatorname{Ad}(V(b))\lambda(b_+)\big].
\end{equation}
It is the generator in left multiplicative link coordinates per unit
length. In the exponential coordinate \(e^{hA}W\), the exact generator
at \(A\) is \(E_{hA}^{-1}D_{e^{hA}W}\), acting pointwise on links after
the difference is taken. This follows by differentiating
\(e^{hA^g}W=(e^{hA}W)^g\).
Differentiating \eqref{eq:gauge-effective-chain} therefore gives the
nonlinear Ward identity. In raw output coordinates it reads
\begin{equation}\label{eq:gauge-nonlinear-ward}
 D_Aa_j(A)\,E_{h_0 A}^{-1}D_{e^{h_0A}W}
  =h_j E_{a_j(A)}^{-1}D_{e^{a_j(A)}W_j}\,
     \tau_{W_{j-1},a_{j-1}}\cdots\tau_{W,a_0}.
\end{equation}
The formula uses \(a_0=h_0A\), so the input generator is per unit
length and the output generator is raw. It extends by
holomorphy on the same charts.

At zero increment set \(Q'_j(W)=q_{W_{j-1}}\cdots q_W\), with
\(Q'_0=1\). Denote by \(\widehat Q^{\mathrm{st}}_j\) and
\(\widehat Q^{\mathrm{pl}}_j\) the straight and plain link derivatives
with both input and output per unit length. They equal
\(Q^{\mathrm{st}}_j/t_j\) and \(Q^{\mathrm{pl}}_j/t_j\) from
Section~\ref{sec:avg-iteration}. The physical input \(A\) becomes
\(h_kA\) in units of the final spacing, and the raw output is
divided by \(h_j\). Their difference in \eqref{eq:avg-plain-st}
is the covariant difference of the corresponding linear site
factor. Their gauge laws are
\begin{equation}\label{eq:gauge-row-ward}
 \widehat Q^{\mathrm{st}}_jD_W=D_{W_j}Q'_j(W),\qquad
 \widehat Q^{\mathrm{pl}}_jD_W=D_{W_j}(\,\cdot\,|_{\mathrm{level}\ j}).
\end{equation}
The first is \eqref{eq:gauge-nonlinear-ward} at zero. For the second,
differentiate the endpoint covariance of \(\mathcal M_{\rm an}^j\).
In particular the straightened row annihilates \(D_W\ker Q'_j\).

For a family of per-length straightened outputs let \(K(c,b)\)
be the derivative matrix at the admitted increment. The colour
transpose \({}^{\mathsf T}\) is taken in the invariant bilinear
pairing. The transpose for physical link measures is
\begin{equation}\label{eq:gauge-link-transpose}
 (K^\ast\omega)(b)=h_0^{-4}\sum_c h_{j(c)}^4
                         K(c,b)^{\mathsf T}\omega(c),
\end{equation}
where \(j(c)\) is the level of \(c\). This is the result of moving
the kernel across the coarse weighted sum and factoring out the
fine measure. For rows on one scale, the factor \(h_0^{-4}\)
is canceled by the column density from
\eqref{eq:avg-layered-coeffs}. In particular the transpose has
bounded supremum norm and its variation from zero to \(A\) is
bounded by \(C_L h_j\|A\|_\infty\), using the fixed-leg
second derivative in that equation. The same proof holds for the
admitted layered rows by their bounded overlap. The bilinear
transpose is holomorphic in the background. For complex colours,
transposing a coefficient costs at most a fixed factor from the
equivalence of norms on the three-dimensional colour space; this
factor is included in $C_L$. For raw output
coordinates the derivative has the additional factor \(h_{j(c)}\)
on each row, which also multiplies the corresponding summand in
\eqref{eq:gauge-link-transpose}.

On one physical scale $h_j$, write $K_A$ for the per-length
derivative. The same estimates give the joint bound
\begin{equation}\label{eq:gauge-joint-transpose}
 \|K_A^*\omega-K_{\widetilde A}^*\widetilde\omega\|_\infty
 \le C_L\|\omega-\widetilde\omega\|_\infty
   +C_Lh_j\|A-\widetilde A\|_\infty\,
                              \|\widetilde\omega\|_\infty .
\end{equation}
Use $K_A^*(\omega-\widetilde\omega)+
(K_A-K_{\widetilde A})^*\widetilde\omega$ and integrate the
fixed-input-leg second derivative along the segment between the
two increments. The physical transpose cancels its fine-link
density. On an admitted layered collection the same argument is
applied on each complete carrier; the bounded overlap gives its
weighted version with that carrier's $h_j$.

For the nonlinear transpose $N_A^*=(K_A-K_0)^*$ the same argument
retains the radius factor in both inputs:
\begin{equation}\label{eq:gauge-nonlinear-transpose}
 \begin{split}
 \|N_A^*\omega-N_{\widetilde A}^*\widetilde\omega\|_\infty
 \le C_Lh_j\bigl(&\|A\|_\infty\|\omega-\widetilde\omega\|_\infty\\
             &+\|A-\widetilde A\|_\infty\|\widetilde\omega\|_\infty\bigr).
 \end{split}
\end{equation}
Split the difference in the same way and integrate $K_A-K_0$
from zero. On a radius-$r$ ball in both variables this gives
a Lipschitz constant $C_Lh_jr$.

For completeness, physical quadratic forms and currents have
different coordinate powers. On a finite coordinate space let
$D=\operatorname{diag}(\ell_b)$ contain the link lengths and
$M=\operatorname{diag}(m_b)$ the positive physical measure weights.
The raw coordinate is $a=DA$. If
\[
 \mathcal Q(A)=\tfrac12 A^{\mathsf T}MK A,\qquad
 dS(A)[H]=J^{\mathsf T}MH,
\]
then their counting-coordinate expressions are
\begin{equation}\label{eq:physical-raw-conversion}
 \widehat K=D^{-1}MKD^{-1},\qquad
 \widehat J=D^{-1}MJ.
\end{equation}
These identities follow by substituting $A=D^{-1}a$ and
$H=D^{-1}\delta a$. When $\ell_b=\ell$ and $m_b=\ell^4$,
they give $\widehat K=\ell^2 K$ and $\widehat J=\ell^3 J$.
At unequal lengths the displayed diagonal factors retain their
order; no single scale replaces them. The transpose is bilinear
on the complexified coordinates.

\subsection{The relative axial gauge on nested blocks}

Let \(B^j(y)\) be the bare block of side \(L^j\) with centre \(y\).
It is the union of the bare blocks of side \(L^{j-1}\) whose centres
lie in the one-step block below \(y\). For a fine site \(x\in B^j(y)\)
the ancestors \(x=x_0,x_1,\ldots,x_j=y\) are uniquely determined.
For \(V=e^{a_0}W\), the mean tower axial condition on \(g\) is
\[
 T_{\mathcal M_{\rm an}^i(V^g)}(x_{i+1},x_i)=T_{W_i}(x_{i+1},x_i)
      \quad (0\le i<j).
\]
For the normalized mean tower axial gauge impose also
\((R_{W_{j-1}}\cdots R_Wg)(y)=1\) at each level-\(j\) centre.

\begin{proposition}[Normalized tower axialization]\label{prop:gauge-tower}
On the depth-uniform small charts of Proposition~\ref{prop:avg-iteration},
the mean tower
axial condition with this normalization determines a unique gauge on
each tower. Its level-\(i\) transported gauge average is
\((g|_{\mathrm{level}\ i})v_i\), and
\(\mathcal M_{\rm an}^j(V^g)=e^{a_j}W_j\).
On the real ball the same equality holds for the Borel map.
The gauge is a holomorphic function of the admitted increments on
its determining set and is real for real fields. Its distance
from \(1\) is at most \(C_L\sum_{i<j}\|a_i\|_\infty\) on the
depth-uniform ball.
\end{proposition}

\begin{proof}
Set \(Z_i=e^{a_i}W_i\). For an admitted solution \(g\), the
level-\(i\) field has the form \(Z_i^{H_i}\) with
\(H_i=(g|_{\mathrm{level}\ i})v_i\), by covariance and
\eqref{eq:gauge-factorization}. To construct these \(H_i\),
apply \eqref{eq:gauge-relative-axial} to \(Z_i\) with
\(R_{W_i}H_i=H_{i+1}\). The only possible choice is
\[
 H_i(x)=T_{W_i}(y,x)^{-1}H_{i+1}(y)
       u_{W_i}(a_i)(y)^{-1}T_{Z_i}(y,x),
   \qquad x\in B(y).
\]
Since \(H_0=g\), the identities \(R_{W_i}H_i=H_{i+1}\) ensure
that \(H_i\) are precisely its iterated averages. Start at \(H_j=1\)
and descend. At \(x=y\) the displayed formula gives
\(H_i(y)=H_{i+1}(y)u_{W_i}(a_i)(y)^{-1}\).
The product rule for \(v_i\) thus makes
\(H_iv_i^{-1}\) agree with \(H_{i+1}v_{i+1}^{-1}\) at the shared
sites. This verifies the restriction of the resulting fine gauge at
every level. To bound \(H_i-1\), move the
reference transport to the far right of the recursion. The relative
contour ratio and the block factor differ from \(1\) by at most
\(C_L\|a_i\|_\infty\). Conjugation by a real reference transport is
isometric. For a shifted complex reference its norm is at most
\(1+C_L\delta_i\), where \(\delta_i\) is the raw shift at this level.
Consequently each descent step gives
\[
 \|H_i-1\|_\infty\le
 (1+C_L\delta_i+C_L\|a_i\|_\infty)\|H_{i+1}-1\|_\infty
     +C_L\|a_i\|_\infty .
\]
The summable bounds of Section~\ref{sec:avg-iteration} control the
product, and \(g=H_0\). At every level the prospective gauge
\(H_iv_i^{-1}\) thus has condition number at most
\(\exp(C_L\sum_{s<j}(\delta_s+\|a_s\|_\infty))\).
After reducing the ball this fits inside the reserve for the
contour families and logarithms used in the factorization.
Covariance is therefore admitted, so the constructed gauge
satisfies the required conditions and the top field equals \(Z_j\).
The descent uses finite products of holomorphic transports and
site factors, all of which are real on real inputs.
\end{proof}

One can also prescribe \(H_j\) at the top in place of \(1\),
with the same chart requirement for complex gauges.
For a decomposition into disjoint pillars with different top levels,
perform the descent separately in each pillar. If the gauge itself
is required to be \(1\) at each top, prescribe \(H_j=v_j\) there.
Relative contour equations ending at sites with already prescribed
gauges remain equations on the datum. In particular a seam between
different pillars does not create an additional gauge parameter.

For real fields there is a second tower gauge using one prescribed
coordinate comb in each block. Take the plain iterates of the real
Borel map \(\mathcal M\), with \(W_i=\mathcal M^i(W)\) in this
construction. Write \(T^\Gamma_X(y,x)=X(\Gamma_{y,x})\).
Prescribe the value \(g(y)\) at each top and put
\[
 g(x)=T^\Gamma_{W_i}(y,x)^{-1}g(y)
                   T^\Gamma_{\mathcal M^i V}(y,x)
\]
on the next finer level, descending from the top of each pillar.
The formula follows from the same endpoint transport law as
\eqref{eq:gauge-relative-axial}. It fixes all free sites uniquely.
If \(x\) is a pinned node, keep its prescribed gauge and continue
the descent below it. When \(y\) and \(x\) are both pinned to \(1\),
the omitted equality at \(x\) is gauge invariant under those
restrictions and must be supplied as part of the datum for a
representative satisfying every contour equation. This is the
single-comb convention for the tower gauge in \cite[\S B]{Bal85g}.
The identification with \eqref{eq:gauge-recursion} above uses mean
transports at every level, as do its site factors and gauge
normalization.

\subsection{Pinned nodes and the action on the datum}

A tower is a finite rooted forest whose vertices are sites on the
successive coarse lattices. Its edges join a site to the centre of its
parent block. Copies of a centre that coincide at successive levels
are identified, and the resulting zero-length parent edges are omitted.
Let $\mathcal N$ be a set of pinned vertices containing
every root. These are the endpoints of the retained coarse constraints;
they may include junction nodes below a root. Put
\[
 \mathcal G_{\mathcal N}=\{g:\text{sites}\to G:\ g(z)=I\text{ for }z\in\mathcal N\text{ or }z\text{ outside the forest}\}.
\]
For each forest edge $(y,x)$ choose the same transport convention used
in the corresponding average. Impose the axial condition
$T_{\mathcal M^iU}(y,x)=I$ when $x$ is not pinned. Omit it when $x$ is
pinned. All transports in this paragraph are assumed to lie in their
real domains; single-comb transports have no domain restriction.

\begin{proposition}[Pinned axial representative]\label{prop:pinned-representative}
Each $\mathcal G_{\mathcal N}$ orbit in this domain has exactly one
representative satisfying these conditions. For a site function
$v:\mathcal N\to G$, define $\widetilde v(x)$ to be its value at the first
pinned ancestor of $x$, including $x$ itself, and set it to $I$
off the forest. This lift is a group
homomorphism and has the prescribed values on $\mathcal N$.
If $\mathcal S(U)$ denotes the axial representative, then
\begin{equation}\label{eq:pinned-covariance}
 \mathcal S(U^{\widetilde v})=\mathcal S(U)^{\widetilde v}.
\end{equation}
If the retained constraint is $\mathcal M^{j(c)}U(c)=V(c)$, its
endpoints being pinned, the lift sends its fibre onto the fibre with
datum $V^v$.
\end{proposition}

\begin{proof}
Set $g=I$ at pinned vertices and descend from the roots. At a nonpinned
child the required equation is
$g(y)T_{\mathcal M^iU}(y,x)g(x)^{-1}=I$, so it determines
$g(x)=g(y)T_{\mathcal M^iU}(y,x)$. At a pinned child retain its fixed
value and start the descent anew below it. Covariance of the averages
makes these the equations for the transformed fine field. Every free
vertex has one parent, so the construction exists and is unique.
The same descent applied to a transformation between two representatives
forces that transformation to be $I$ at every vertex. This proves
uniqueness of the representative.

The first pinned ancestor is the same for $y$ and a nonpinned child $x$.
Thus $\widetilde v(y)=\widetilde v(x)$ on every imposed edge, and its
identity transport remains the identity. Conjugation by
$\widetilde v$ preserves $\mathcal G_{\mathcal N}$. It consequently
sends an orbit and its unique representative to the transformed orbit
and representative, proving \eqref{eq:pinned-covariance}. The lift
respects products pointwise. Ordinary covariance of each
$\mathcal M^{j(c)}$ proves the last assertion.
\end{proof}

At a pinned child of a pinned centre, both endpoints of the omitted
transport are fixed by $\mathcal G_{\mathcal N}$, so it is constant
on the orbit. Suppose $y$ is pinned and the forest includes every
site of $B(y)\setminus\{y\}$ as a child of $y$ at this level.
The full block contour energy in the representative is then
\begin{equation}\label{eq:pinned-residual-energy}
 \sum_{\substack{x\in B(y)\setminus\{y\}\\x\in\mathcal N}}
       \left(1-\tfrac12\operatorname{tr}T_{\mathcal M^iU}(y,x)\right).
\end{equation}
The other summands vanish in the representative. An omitted pinned
transport need not be determined by the retained link datum alone.

Let \(\mathfrak P\) be the finite group of signed coordinate
permutations preserving the periods, the lattice grids, and the
family of admissible carrier configurations. For a scalar quantity
\(H_\Gamma(U)\) defined using a frame \(\Gamma\), average
\(H_{s\Gamma_0}(U)\) over \(s\in\mathfrak P\), counting each group
element once. If simultaneous relabelling of the field and frame
preserves the construction, the average is invariant under
\(\mathfrak P\) by reindexing the sum. On a symmetric torus and
carrier family this is the full signed-permutation group of order
\(2^4 4!\). No average of the group-valued representatives is required.

The normalization by nonlinear site averages has a different algebraic
property, examined in Section~\ref{sec:gauge-products}. The exact
effective-gauge identities concern transformations based at the current
field; they do not assert a group law for the normalized set.

\subsection{Transported site averages and outward endpoints}

Let \(\mathcal T_j(y,x)\) be the ordered product from \(y=x_j\)
down to \(x=x_0\) of \(T_{W_{i-1}}(x_i,x_{i-1})\),
\(i=j,j-1,\ldots,1\). Expanding \eqref{eq:gauge-site-one}
at each level and using the unique ancestors gives
\begin{equation}\label{eq:gauge-site-row}
 Q'_j(W)\lambda(y)=L^{-4j}\sum_{x\in B^j(y)}
       \operatorname{Ad}(\mathcal T_j(y,x))\lambda(x).
\end{equation}
On disjoint pillars use this formula at each top \(y\) at its
assigned level \(j(y)\), and denote the resulting operator by \(q\).
Use the invariant bilinear colour pairing
\(-\tfrac12\operatorname{tr}(XY)\) and measures \(h_0^4\) on fine
sites and \(h_{j(y)}^4\) at the tops. Define \(q^\ast\) to be the
transpose for these pairings. The site lift is exactly
\begin{equation}\label{eq:gauge-site-lift}
 q^\ast\psi(x)=\operatorname{Ad}(\mathcal T_{j(y)}(y,x))^{-1}\psi(y),
   \quad x\in B^{j(y)}(y),\qquad qq^\ast=1 .
\end{equation}
Invariance of the colour pairing and \(h_0^4L^{4j}=h_j^4\)
give the transpose formula. Inserting it in \eqref{eq:gauge-site-row}
gives the identity. For real transports the lift is an isometry
for these \(L^2\) norms and has supremum norm at most one. The algebraic
identities also hold on complex charts. Norm estimates there use the
condition numbers of the transports. For a complex perturbation of a
real reference on the depth-uniform ball, the difference of the
one-step transport from its real value is bounded by \(C_L\delta_i\),
where \(\delta_i\) is the raw relative increment at that level.
A product telescoping estimate bounds the change in every
\(\mathcal T_j(y,x)\) by \(C_L\sum_{i<j}\delta_i\). This sum is
bounded uniformly by Section~\ref{sec:avg-iteration}. Thus each
coefficient of the forward row \(Q'_j\) is bounded by
\(C_L L^{-4j}\) and has a bounded holomorphic continuation.
The coefficients of the physical transpose \(q^\ast\) are instead
bounded by \(C_L\), as \eqref{eq:gauge-site-lift} shows.
For reference derivatives use the normalized input \(B\) in
\(W'=e^{\eta B}W\), with \(\eta=L^{-k}\) as in
Proposition~\ref{prop:avg-iteration}. Cauchy's formula on a smaller
fixed \(B\)-ball bounds an order-\(m\) derivative of a forward-row
coefficient by \(C_{L,m}L^{-4j}\prod_{r=1}^m\|B_r\|_\infty\).
Here \(B_r\) are the derivative directions. Derivatives in raw
fine-link logarithms require the corresponding factors of
\(\eta^{-1}\).

On a single top scale, let \(\tau\) be the effective site derivative
in \eqref{eq:gauge-nonlinear-ward}. For a prescribed vector
\(\lambda_0\) of fine site values the correction
\begin{equation}\label{eq:gauge-effective-kernel}
 \lambda=\lambda_0-q^\ast(\tau q^\ast)^{-1}\tau\lambda_0
\end{equation}
has \(\tau\lambda=0\). Indeed \(qq^\ast=1\) and the uniform
bound on \(\tau-q\) proves invertibility of \(\tau q^\ast\)
by its Neumann series on a smaller ball. When \(q\lambda_0=0\)
the correction has size at most
\(C_L\|\tau-q\|_{\infty\to\infty}\|\lambda_0\|_\infty\).
This is an algebraic construction in the site variables.

For a level-\(j\) determining link, the right side of
\eqref{eq:gauge-row-ward} uses \(Q'_j\lambda\) at both endpoints.
If one endpoint \(z\) is in the next finer layer, assume its full
backward block \(B^j(z)\) is a disjoint union of the \(L^4\)
level-\((j-1)\) pillars below it. The outer endpoint is then read as
\begin{equation}\label{eq:gauge-outward-site}
 Q'_j(W)\lambda(z)=L^{-4}\sum_{x\in B(z)}
   \operatorname{Ad}(T_{W_{j-1}}(z,x))\,q\lambda(x).
\end{equation}
This is \eqref{eq:gauge-site-one} on the last subdivision.
On an interior row the two endpoint values are the assigned values
of \(q\). Thus the same equation \(\widehat Q^{\mathrm{st}}_jD_W=
D_{W_j}Q'_j\) on an outward row uses the extra averaging
\eqref{eq:gauge-outward-site} before taking the covariant difference.
It annihilates \(D_W\ker q\) on both kinds of rows.

\subsection{Products of transported gauge averages}
\label{sec:gauge-products}

The nonlinear site averages nearly preserve multiplication on slowly
varying gauge transformations. Their product defect is governed by
variation within the blocks, which can be much smaller than the
distances of the site values from the identity. The same distinction
controls the effect of multiplying by a complex gauge transformation.

Let \(W\) be a real reference and \(W_j=\mathcal M^j W\) for
\(0\leq j\leq k\). Write \(t_j=h_j/h_k=L^{j-k}\), and set
\(\mathcal R_0 u=u\) and
\(\mathcal R_{j+1}u=R_{W_j}(\mathcal R_j u)\), with \(R_{W_j}\)
from \eqref{eq:gauge-site-nonlinear}. For a site field \(u_j\) on
level \(j\), its transported variation within a block below \(y\)
is
\[
 \omega_j(u_j;y)=
 \max_{x\in B(y)}
  \left\|\operatorname{Ad}(T_{W_j}(y,x))u_j(x)-u_j(y)\right\|.
\]
Write \(\omega_j(u_j)=\sup_y\omega_j(u_j;y)\).
For a real gauge transformation \(u\) we say that \(u\) satisfies
the bounds of size \(\alpha\) if its iterates \(u_j=\mathcal R_j u\)
are defined and
\begin{equation}\label{eq:gauge-product-class}
 \|u_j-I\|_\infty\leq\alpha\quad (0\leq j\leq k),
 \qquad \omega_j(u_j)\leq\alpha t_{j+1}\quad (0\leq j<k).
\end{equation}

For real inputs, the normalized tower gauge $g$ of
Proposition~\ref{prop:gauge-tower} at depth $k$ satisfies
\eqref{eq:gauge-product-class} with size $C_L\|A\|_\infty$.
Its level-$i$ iterate is $H_i$ in that proof. At a block below
$y$, the descent gives
\[
 \operatorname{Ad}(T_{W_i}(y,x))H_i(x)-H_i(y)
  =H_i(y)\bigl(T_{Z_i}(y,x)T_{W_i}(y,x)^{-1}-I\bigr).
\]
The right side has norm at most $C_L\|a_i\|_\infty$ by the
one-step contour estimate. This is at most
$C_Lt_i\|A\|_\infty$ by \eqref{eq:avg-raw-radii}.
The same descent bounds $\|H_i-I\|_\infty$ by
$C_L\sum_{s=i}^{k-1}\|a_s\|_\infty\leq C_L\|A\|_\infty$.
Define the unscaled covariant difference of a matrix-valued site
field \(f\) at a level-\(j\) bond \(b\) by
\[
 (\nabla^{\rm raw}_{W_j} f)(b)
      =f(b_-)-\operatorname{Ad}(W_j(b))f(b_+).
\]
For Lie algebra fields this is \(h_jD_{W_j}\) in the convention
of \eqref{eq:gauge-D}.

\begin{proposition}[Products and perturbations of gauge averages]
\label{prop:gauge-products}
Suppose \(W\) is real on a torus admitting \(k\) blocking steps.
Put \(t_j=L^{j-k}\) for \(0\le j\le k\), and assume the fine
plaquette defect is at most \(\kappa t_0^2\). There are
\(c_L>0\) and \(C_L<\infty\), independent of \(k\), mesh and
volume, for which the following hold when
\(0\leq\kappa,\alpha,\rho<c_L\).
\begin{enumerate}
\item If two real gauges \(u,v\) satisfy
\eqref{eq:gauge-product-class}, then their pointwise product satisfies
the bounds of size \(2\alpha+C_L\alpha^2\). In addition,
\begin{equation}\label{eq:gauge-product-defect}
 \mathcal R_j(uv)=(\mathcal R_j u)(\mathcal R_j v)e^{z_j},
 \qquad
 \|z_j\|_\infty\leq C_L\alpha^2t_j^2
 \quad (0\leq j\leq k).
\end{equation}
\item Let \(v\) be a real gauge satisfying
\eqref{eq:gauge-product-class}. For a fine complex Lie algebra
field \(\lambda\) with
\[
 \|\lambda\|_\infty\leq\rho,\qquad
 \|\nabla^{\rm raw}_W\lambda\|_\infty\leq\rho t_0,
\]
the quotients
\(f_j=\mathcal R_j(e^\lambda v)(\mathcal R_jv)^{-1}\)
are defined and holomorphic in \(\lambda\) on the corresponding
open ball. They satisfy
\begin{equation}\label{eq:gauge-product-quotient}
 \begin{aligned}
 \|f_j-I\|_\infty&\leq C_L\rho,&
 \|\nabla^{\rm raw}_{W_j} f_j\|_\infty&\leq C_L\rho t_j,\\
 \|\log f_j-Q'_j(W)\lambda\|_\infty
    &\leq C_L(\alpha+\rho)\rho t_j.
 \end{aligned}
\end{equation}
Here the final term is the site row from
\eqref{eq:gauge-site-row}, and the logarithm is the branch at \(I\).
\end{enumerate}
\end{proposition}

\begin{proof}
We first compute a one-step quotient. On a block below \(y\) write
\(P_x=\operatorname{Ad}(T_{W_j}(y,x))\). For site fields \(f,v\)
on level \(j\), define
\[
 X_x=\log\bigl(f(y)^{-1}P_x f(x)\bigr),\qquad
 Y_x=\log\bigl(v(y)^{-1}P_x v(x)\bigr),\qquad
 \widetilde Y_x=\operatorname{Ad}(v(y))Y_x .
\]
In the one-step computation \(v\) is real and \(f\) is either real
or in a fixed small complex ball about \(I\). Take the variations
\(\omega_f,\omega_v\) of these fields in this block small enough
that \(\|X_x\|\leq C\omega_f\) and
\(\|\widetilde Y_x\|\leq C\omega_v\). Denote the average over
\(x\in B(y)\) by brackets, and let
\(\mathcal L(X,Y)=\log(e^Xe^Y)\).
Transport is multiplicative on matrix products. Thus
\[
 (f(y)v(y))^{-1}P_x(f(x)v(x))
        = v(y)^{-1}e^{X_x}v(y)e^{Y_x}.
\]
Conjugating the defining logarithms in
\eqref{eq:gauge-site-nonlinear} by \(v(y)\) gives
\begin{equation}\label{eq:gauge-product-one-step}
 \begin{split}
 R_{W_j}(fv)(y)R_{W_j}(v)(y)^{-1}
   &= f(y)\exp\bigl(\langle\mathcal L(X_x,\widetilde Y_x)\rangle\bigr)
                    \exp(-\langle\widetilde Y_x\rangle)\\
   &= f(y)\exp\bigl(\langle X_x\rangle+e(y)\bigr),
       \qquad \|e(y)\|\leq C\omega_f\omega_v .
 \end{split}
\end{equation}
To see the bound, \(\mathcal L(X,Y)-X-Y\) is analytic on a
fixed small ball and vanishes when either input is zero. Integrating
its bounded mixed derivative gives \(C\|X\|\|Y\|\).
The final logarithm of the product of exponentials in the first
line has the same property after subtracting \(\langle X_x\rangle\).
Indeed, put
\(A=\langle X_x\rangle\), \(B=\langle\widetilde Y_x\rangle\)
and \(E=\langle\mathcal L(X_x,\widetilde Y_x)-X_x-
\widetilde Y_x\rangle\). The logarithm is
\(\mathcal L(A+B+E,-B)\). Changing \(E\) to zero costs
\(C\|E\|\), and \(\mathcal L(A+B,-B)-A\) vanishes on both
coordinate axes and has a bounded mixed derivative. This proves
\eqref{eq:gauge-product-one-step}. The argument uses only fixed
matrix logarithm balls. A sufficiently small bound on
\(\omega_f+\omega_v\) admits all its relative logarithms.

Two consequences will be used. First,
\(R_{W_j}(f)(y)=f(y)e^{\langle X_x\rangle}\), so the quotient in
\eqref{eq:gauge-product-one-step} differs from \(R_{W_j}(f)(y)\)
by at most \(C\omega_f\omega_v\). Second, if
\(F(y)=\langle P_xf(x)\rangle\), the same quotient differs from
\(F(y)\) by at most
\begin{equation}\label{eq:gauge-product-affine}
 C\omega_f(\omega_f+\omega_v).
\end{equation}
Indeed \(F(y)=f(y)\langle e^{X_x}\rangle\) and the constant and
linear terms in its expansion coincide with those of
\(f(y)e^{\langle X_x\rangle}\).

For the first assertion, set \(u_j=\mathcal R_j u\),
\(v_j=\mathcal R_jv\), \(p_j=u_jv_j\) and
\(d_j=\|\mathcal R_j(uv)-p_j\|_\infty\) whenever the iterate
is defined. Since the two fields are unitary,
\(\omega_j(p_j)\leq 2\alpha t_{j+1}\).
If \(\mathcal R_j(uv)=g_jp_j\), then \(g_j\) is unitary,
\(\|g_j-I\|_\infty=d_j\) and
\(\omega_j(g_j)\leq 2d_j\).
Formula \eqref{eq:gauge-product-affine} for \(f=g_j\), \(v=p_j\)
bounds the next quotient's distance from \(I\) by
\(d_j+C d_j(d_j+\alpha t_{j+1})\). The first consequence
of \eqref{eq:gauge-product-one-step}, applied to \(u_j,v_j\),
then gives
\begin{equation}\label{eq:gauge-product-recurrence}
 d_{j+1}\leq d_j(1+C(d_j+\alpha t_{j+1}))
                   +C\alpha^2t_{j+1}^2,\qquad d_0=0 .
\end{equation}
Choose \(K\) so \(2K L^{-2}+C\leq K\), then decrease the
ceiling on \(\alpha\) so the coefficient of \(d_j\) is at most
two under \(d_j\leq K\alpha^2t_j^2\). The same ceiling admits
the logarithms at step \(j+1\), since the variation of \(g_jp_j\)
is bounded by \(2d_j+2\alpha t_{j+1}\).
Induction proves that every iterate is defined and
\(d_j\leq K\alpha^2t_j^2\). Taking the series logarithm of
\(p_j^{-1}\mathcal R_j(uv)\) proves \eqref{eq:gauge-product-defect}.
The distance to \(I\) is at most \(2\alpha+d_j\) and the
variation at most \(2\alpha t_{j+1}+2d_j\). These are
\eqref{eq:gauge-product-class} with the asserted enlarged size.

For the second assertion, set \(v_j=\mathcal R_jv\).
At level zero \(f_0=e^\lambda\). A recursive description of
the quotient is
\[
 f_{j+1}=R_{W_j}(f_jv_j)R_{W_j}(v_j)^{-1}.
\]
To estimate its covariant differences, let
\[
 m_j=\|f_j-I\|_\infty,\qquad
 d_j=\|\nabla^{\rm raw}_{W_j} f_j\|_\infty,\qquad
 \epsilon_j=\sup_p\|W_j(\partial p)-I\|.
\]
By Corollary~\ref{cor:avg-real-iterate},
\(\epsilon_j\leq C_L\kappa t_j^2\), since the small logarithm
and group distance are comparable. For the mean transport
\(T_{W_j}(y,x)\), any of its coordinate paths is within
\(C_L\epsilon_j\) of that mean in operator norm, by the path
interchange estimate and the group mean bound. Transport
\(f_j(x)-I\) along such a path and telescope its covariant
differences. The path has at most \(2L\) links, so
\begin{equation}\label{eq:gauge-product-block-variation}
 \omega_j(f_j)\leq C_L(d_j+\epsilon_jm_j).
\end{equation}

For the arithmetic transported output
\(F(y)=L^{-4}\sum_{x\in B(y)}P_xf_j(x)\), one has
\(\|F-I\|_\infty\leq m_j\).
For a coarse bond \(c=(y,y')\) let \(x'=x+(y'-y)\).
Compare the transport \(T_{W_j}(y,x)W_j([x,x'])\)
with \(W_{j+1}(c)T_{W_j}(y',x')\).
The loop correction in \eqref{eq:avg-map} bounds their
difference by \(C_L\epsilon_j\) for each \(x\).
Thus
\[
 \|F(y)-\operatorname{Ad}(W_{j+1}(c))F(y')\|
                  \leq Ld_j+C_L\epsilon_jm_j .
\]
In fact the change of transport is applied only to
\(f_j(x')-I\), which gives the second term. The first term
comes from the \(L\) adjacent differences along \([x,x']\)
in unitary frames, averaged over \(x\).
Set \(s_j=C_L(d_j+\epsilon_jm_j)\), with \(C_L\) large
enough for \eqref{eq:gauge-product-block-variation}.
By \eqref{eq:gauge-product-affine},
\begin{equation}\label{eq:gauge-product-quotient-recurrence}
 \begin{split}
 m_{j+1}&\leq m_j+C_Ls_j(s_j+\alpha t_{j+1}),\\
 d_{j+1}&\leq Ld_j+C_L\epsilon_jm_j+
                             C_Ls_j(s_j+\alpha t_{j+1}).
 \end{split}
\end{equation}
The coefficient \(L\) here comes from the straight segment count,
even though the individual site transports use the contour mean.

The exponential derivative bound on the \(\rho\)-ball gives
\(m_0\leq e^\rho\rho\), \(d_0\leq e^\rho\rho t_0\).
As long as \(m_j\leq 2\rho\) and \(d_j\leq2\rho t_j\),
one has \(s_j\leq C_L\rho t_j\). The additive terms in
\eqref{eq:gauge-product-quotient-recurrence} are therefore
at most \(C_L(\kappa+\alpha+\rho)\rho t_j^2\).
For \(m_{j+1}\) sum these terms directly. For
\(d_{j+1}/t_{j+1}\) the term \(Ld_j/t_{j+1}\) is \(d_j/t_j\),
and the remaining terms sum to at most
\(C_L(\kappa+\alpha+\rho)\rho\sum_j t_j\).
Both geometric sums are bounded independently of depth.
After reducing the ceilings these estimates give the same two
bounds at the next step, with strict reserve when \(\rho>0\).
The small variations also admit all relative logarithms in
\eqref{eq:gauge-product-one-step}. When \(\rho=0\) the quotient
is \(I\) at every level. Holomorphy follows from the admitted
finite formulas.

It remains to compare the quotient logarithm with the linear
row. At a block centre put \(Z=\log f_j(y)\).
The logarithms of \(P_x f_j(x)\) are
\(\mathcal L(Z,X_x)\). The logarithm of the quotient at the
next level is \(\mathcal L(Z,\langle X_x\rangle+e(y))\)
by \eqref{eq:gauge-product-one-step}. The same mixed
derivative estimate for \(\mathcal L\) now gives
\[
 \|\log f_{j+1}-q_{W_j}\log f_j\|_\infty
   \leq C_L s_j(m_j+\alpha t_{j+1})
   \leq C_L(\alpha+\rho)\rho t_j .
\]
Here the linear terms \(Z+\langle X_x\rangle\) cancel.
The operator \(q_{W_j}\) in \eqref{eq:gauge-site-one}
has supremum norm at most one for a real reference.
Since \(\log f_0=\lambda\), iteration and
\(\sum_{s<j}t_s\leq t_j/(L-1)\) prove the last line of
\eqref{eq:gauge-product-quotient}.
\end{proof}

The bounds adapt the product and quotient estimates for a single
contour in \cite[\S F, Props.~8--10]{Bal85a}. In particular the
defect in \eqref{eq:gauge-product-defect} need not vanish when the
two separately averaged gauges are \(I\). At the trivial
reference, choose two noncentral sites of one block and let the
two gauges take \(e^X,e^{-X}\) and \(e^Y,e^{-Y}\) there,
respectively, with all other values \(I\). Their separate
averages are \(I\). The logarithm of the product average has
leading term \(L^{-4}[X,Y]\), which is nonzero for arbitrarily
small noncommuting \(X,Y\).

The effective gauge about a nonzero field also has a uniform
finite Taylor estimate. Measure lengths in final-scale units,
put $h=L^{-k}$, and let the real reference $W$ have plaquette
defect at most $\kappa h^2$. For a complex site logarithm set
$\|\lambda\|_1=\max\{\|\lambda\|_\infty,
\|D_W\lambda\|_\infty\}$, where $D_W$ divides by $h$.
If $\kappa,\|A\|_\infty,\|\lambda\|_1$ are sufficiently small
depending on $L$, define the transformed fine coordinate by
$(e^{hA}W)^{e^\lambda}=e^{hA^\lambda}W$ and use the effective
gauge $r_k$ in \eqref{eq:gauge-effective-chain}. Then
\begin{equation}\label{eq:gauge-finite-effective-taylor}
 \|\log r_k(A,e^\lambda)-Q'_k(W)\lambda\|_\infty
       \le C_L(\|A\|_\infty+\|\lambda\|_1)\|\lambda\|_1 .
\end{equation}
The map and the logarithm are holomorphic on this joint ball,
uniformly in mesh, depth and volume.

Here is a direct admission and derivative argument. At a link write
$\lambda_+=\operatorname{Ad}(W(b))\lambda(b_+)$ and
$\lambda_- =\lambda(b_-)$. Their difference has norm at most
$h\|\lambda\|_1$. The exact relative transformed link is
$e^{\lambda_-}e^{hA(b)}e^{-\lambda_+}$.
Factoring out the conjugate of $e^{hA(b)}$ and integrating the
exponential between $\lambda_+$ and $\lambda_-$ shows that its
logarithm is $O(h(\|A\|_\infty+\|\lambda\|_1))$ on a fixed
complex ball. Division by $h$ therefore yields a uniformly bounded
holomorphic $A^\lambda$. Proposition~\ref{prop:avg-iteration}
admits both fine fields and bounds their accumulated factors.
The effective factors
\((v_j(A^\lambda))^{-1}(e^\lambda|_{\mathrm{level}\ j})v_j(A)\)
are close to \(I\) at every level, uniformly in depth. Their
endpoint conjugations therefore fit the one-step chart reserves.
The exact formula
$r_k=v_k(A^\lambda)^{-1}e^{\lambda|_{\mathrm{coarse}}}v_k(A)$
then gives a uniformly bounded logarithm on a larger joint ball.
Its value at $\lambda=0$ is zero. Its $\lambda$ derivative at
$(A,\lambda)=(0,0)$ is $Q'_k(W)$ by
\eqref{eq:gauge-site-one} and composition. Cauchy's formula on
the larger ball bounds its second derivatives independently of
dimension. Integrating the $\lambda$ derivative from zero to
$\lambda$ gives \eqref{eq:gauge-finite-effective-taylor}.

\subsection{The flat link rows}

At the trivial reference, every one-step straight derivative is
\(S_1\). For \(n=L^j\) abbreviate
\(\widehat Q_j^{\mathrm{st}}\) by \(\widehat Q_j\) in this
subsection. Iterating \eqref{eq:gauge-straight-row} gives
\begin{equation}\label{eq:gauge-flat-row}
 (\widehat Q_j A)_\mu(y)= n^{-5}\!
    \sum_{\substack{u\in\mathbb Z^4\\|u|_\infty\le(n-1)/2}}
       \sum_{s=0}^{n-1} A_\mu(y+h_0u+h_0s e_\mu).
\end{equation}
For the induction, each coarse start and its finer offset specify
one fine start in the larger block. The segment indices at the two
steps combine as \(n s_{\rm coarse}+s_{\rm fine}\), each occurring
once. This proves the stated normalization and support.
The coefficients are nonnegative, sum to one and are each at most
\(n^{-4}\), since at most \(n\) pairs of longitudinal indices give
any one input. With the fine measure \(h_0^4\), Cauchy's inequality
therefore gives on its double block \(D_c\)
\begin{equation}\label{eq:gauge-flat-row-L2}
 |(\widehat Q_jA)(c)|\le h_j^{-2}\|A_\mu\|_{L^2(D_c)}.
\end{equation}
For a flat row at a seam, assign length \(h_j/L\) to fine links
internal to the tail block and \(h_j\) to the remaining links.
In one longitudinal line there are \(n^2\) segment occurrences,
of which \(\sum_{s=1}^{n-1}s=n(n-1)/2\) are internal to the
tail block. The absolute weighted sum of the raw row coefficients
with weights \(1/\xi(b)\) is therefore
\[
 1+(L-1)\frac{n-1}{2n}.
\]
This exhibits the adjacent-scale cost in the weighted row bound of
Section~\ref{sec:avg-iteration} using the actual flat coefficients.

The exact Fourier multiplier on \(A_\mu(x)=e^{i p\cdot x}a\) is
the product
\[
 \left(\prod_{\nu=1}^4\frac1n\sum_{u=-(n-1)/2}^{(n-1)/2}
                         e^{ih_0p_\nu u}\right)
 \left(\frac1n\sum_{s=0}^{n-1}e^{ih_0p_\mu s}\right),
\]
times \(e^{ip\cdot y}a\). The site multiplier omits the last
factor. Both formulas include zero frequencies by their finite
sums.

The map \(\widehat Q_j\) has an exact right inverse obtained by
assigning \(nB(c)\) to each fine edge crossing the face between the
two bare blocks of \(c\), and zero to the internal edges.
Each segment crosses its designated face once, so
\eqref{eq:gauge-flat-row} sends this lift to \(B\).
With the physical link measures \(h_0^4\) and \(h_j^4\), the
Gram operator \(\widehat Q_j\widehat Q_j^\ast\) has diagonal
coefficient \((2+n^{-2})/3\) and coefficient
\((1-n^{-2})/6\) on each of its two nearest rows in direction \(\mu\).
There is no overlap with transverse shifts or other components.
To see the coefficients, along the segment direction the
multiplicities are \(1,2,\ldots,n-1,n,n-1,\ldots,1\).
Their squared sum is \((2n^3+n)/3\).
The overlap with the list translated by \(n\) is
\(\sum_{s=1}^{n-1}s(n-s)=(n^3-n)/6\).
The transverse multiplicity contributes \(n^3\), and the weighted
transpose contributes \(n^4\). With the two factors \(n^{-5}\)
these reduce to \(n^{-3}\) times the two sums above. Periodization
combines coefficients when neighbors coincide. In particular
\[
 \tfrac13 I\le \widehat Q_j\widehat Q_j^\ast\le I .
\]

\subsection{Comparison with coordinate combs}

The one-step block analysis also applies to a single coordinate
comb in place of the contour mean. To compare with the corner-based
blocks of \cite{Bal85a}, identify the blocks by translating their
representatives at level \(j\) from the corner to the centre.
Relative to the common fine grid this displacement is
\(\frac12(L^j-1)h_0(1,1,1,1)\). Since \(L\) is odd, each
translation aligns the tiling at that level. Formula
\eqref{eq:avg-flat-link} uses the chosen contour sums in the
plain derivative. The subtraction in \eqref{eq:gauge-straighten}
removes both sums for either choice. Thus the flat straight and
site coefficients agree after alignment.

For real \(W\), let \(D\) be the projected vertex set of the two
depth-\(j\) blocks of a row, with an embedded rectangular lift.
Assume every physical fine plaquette with all vertices in \(D\)
has defect at most \(\alpha L^{-2j}\), where \(\alpha\) is sufficiently
small depending on \(L\). This includes opposite-face tests at
the shortest period. By alignment, \(D\) is a union of complete
blocks at each preceding level. The local curvature argument in
Corollary~\ref{cor:avg-local-real-iterate} therefore admits all
intermediate real branch tests. The logarithm and group-distance
bounds are comparable on these small charts.

The coordinate tree gauge on the rectangular lift makes its raw
fine link logarithms bounded by \(C_L\alpha L^{-j}\), by
Lemma~\ref{lem:avg-comb}. Its value at the centre of the tail block
is \(1\). The additional projected edges enter through their
curvature tests; the following link-logarithm bounds concern the
literal path stock. In this gauge the raw reference logarithms of
either scheme at intermediate level \(i\) are bounded by
\(C_L\alpha L^{i-j}\). Indeed the one-step estimates of
Proposition~\ref{prop:gauge-one-step} apply to the comb scheme
by the same finite-path argument. Starting at the trivial reference,
the straight recursion has \(\|a_{i+1}\|\le
L\|a_i\|(1+C_L\|a_i\|)\), with
\(\|u_i-1\|\le C_L\|a_i\|\). The geometric sum of the raw
radii and the factorization bound both the straight and plain
logarithms by \(C_L\alpha L^{i-j}\), as in
Section~\ref{sec:avg-iteration}. Restrict these computations to
the backward determining blocks of the row.

In this common fine gauge a telescoping product of one-step link
derivatives against \(S_1\) gives row error at most \(C_L\alpha\)
and column error at most \(C_L\alpha L^{-4j}\) for the
per-length straight row. At each insertion the relative one-step
error is \(C_L\alpha L^{i-j}\) in both sums. The normalized
straight product has row scale \(1\) and column scale \(L^{-4j}\),
so summing these errors proves the bounds. In particular each entry
error has the column bound, and an input edge appears in only a
fixed number of level-\(j\) rows. In the site product the
transport in this gauge differs from \(1\) by at most
\(C_L\sum_{i<j}\alpha L^{i-j}\), which gives the same error scales
directly from \eqref{eq:gauge-site-row}.

Returning to the original gauge compares the rows in their
different output frames. If \(y,z\) are the centre and corner of
the tail block, the corner row is transported to \(y\) by
\(\operatorname{Ad}(W(\gamma_{yz}))\) before subtraction,
where \(\gamma_{yz}\) is the path of this common coordinate tree.
Both rows acquire the same inverse transformation at each input
site or link tail, which is isometric. Thus the real comparison is
local, with the stated row and column errors, in the specified
output frame. It applies to the straight and site rows used in
the gauge identities above.

\subsection{Smooth extension of a coarse gauge}\label{sec:smooth-node-gauge}

Take a real flat reference \(W\) on the torus, measure lengths
in units of the final coarse spacing, and put \(h=L^{-k}\).
The fine and coarse lattices on the cover are \(h\mathbb Z^4\)
and \(\mathbb Z^4\). Their common periods are \(M_\mu e_\mu\)
with \(M_\mu\) integers at least two. The bounds below for a
single gauge and its extension also apply to every odd mesh
\(h=1/n\), \(n\geq1\).

On the covering lattice put \(P(x)=W(\gamma_{0x})\), where
\(\gamma_{0x}\) is any path from \(0\) to \(x\). The path
interchange identity and flatness give path independence.
In particular \(P(x+h e_\mu)=P(x)W(x,\mu)\), so the gauge \(P\)
makes \(W\) the identity on the cover. For a period vector
\(a\), periodicity gives \(P(x+a)=H_aP(x)\) with \(H_a=P(a)\).
The \(H_a\) form a commuting family of unitary matrices. A site
field, or a link increment attached to its initial point \(x\),
is represented in this frame by conjugation with \(P(x)\).
Its representative satisfies the deck law
\(f(x+a)=\operatorname{Ad}(H_a)f(x)\). All the following
norms and ordinary differences are taken in this frame.

For \(0<\alpha<1\), put
\[
 [f]_{\alpha,\mathrm{loc}}
 =\sup_{\substack{x,x'\in h\mathbb Z^4\\0<|x-x'|\leq1}}
          \frac{|f(x)-f(x')|}{|x-x'|^\alpha},
 \qquad
 \nabla_{h,\mu} f(x)=\frac{f(x+h e_\mu)-f(x)}h .
\]
The distance between sites here is Euclidean. For an array, such as
the four positive link components, take the maximum over components.
Let \(\nabla_h^r\) consist of all products of \(r\) forward
differences. Define
\begin{equation}\label{eq:smooth-gauge-norms}
 \begin{split}
 \|f\|_{X_\alpha}&=\|f\|_\infty+\|\nabla_h f\|_\infty
                         +[\nabla_h f]_{\alpha,\mathrm{loc}},\\
 \|f\|_{Y_\alpha}&=\|f\|_\infty+\|\nabla_h f\|_\infty
               +\|\nabla_h^2 f\|_\infty
                         +[\nabla_h^2 f]_{\alpha,\mathrm{loc}}.
 \end{split}
\end{equation}
Simultaneously translating both sites by a period conjugates their
values by the same unitary matrix and preserves these norms.
To compute the local seminorm on period representatives, choose the
first site in a fundamental domain and retain every lift of the second
site at distance at most one. If several short lifts represent the
same torus pair, all enter the supremum. On the torus each lift uses
parallel transport of \(W\) along that displacement. The
covariant difference \(D_W\) of \eqref{eq:gauge-D} is
\(-\nabla_h\) in the parallel frame.

\begin{lemma}[Extensions of coarse site data]\label{lem:smooth-node-extension}
Fix \(0<\alpha<1\). Let \(W\) be a real flat connection on the
periodic lattices of this subsection, with fine spacing \(h=L^{-k}\),
coarse spacing one, and integer periods \(M_\mu\ge2\). Use the
parallel-frame norms in \eqref{eq:smooth-gauge-norms}.
There is a linear operator \(E\) from coarse \(\mathfrak g_{\mathbb C}\)-valued
site fields to fine site fields with
\begin{equation}\label{eq:smooth-node-extension}
 (E\nu)|_{\mathrm{coarse}}=\nu,\qquad
 \|E\nu\|_{Y_\alpha}\leq C_\alpha\|\nu\|_\infty .
\end{equation}
The column at \(y\) is supported within distance \(1/4\) of \(y\)
in the maximum norm, inside its centered block. The constants are
uniform in the mesh, the periods, and the flat holonomies. For
coarse \(w\) in a fixed logarithm chart at the identity,
\begin{equation}\label{eq:smooth-group-extension}
                    \mathfrak E(w)=\exp(E\log w)
\end{equation}
has restriction \(w\) to the coarse sites. It is holomorphic on
the complex chart and group-valued in \(G\) for real data.
For the transported site average \(q=Q'_k(W)\) of
\eqref{eq:gauge-site-row}, there is a smooth lift \(E_q\) with
the same column supports and
\begin{equation}\label{eq:smooth-average-lift}
 qE_q=I,\qquad
 \|E_q\nu\|_{Y_\alpha}\leq C_\alpha\|\nu\|_\infty,\qquad
 \|E_q\nu\|_{H^1_h}\leq C\|\nu\|_{\ell^2}.
\end{equation}
Here the fine \(H^1_h\) norm squares and sums the field and its
first differences with weight \(h^4\), and the coarse \(\ell^2\)
norm uses counting measure. Each sum is over one period.
\end{lemma}

\begin{proof}
Choose a nonnegative smooth function \(\chi\) supported in
\((-1/4,1/4)^4\) with \(\chi(0)=1\) and a positive lower bound
on \([-1/8,1/8]^4\). In the parallel frame
extend the coarse site function \(\nu\) to a smooth function
on the cover by
\[
                 x\longmapsto\sum_{y\in\mathbb Z^4}\chi(x-y)\nu(y).
\]
The supports are disjoint, and translation by a period gives
the same deck law as for \(\nu\). The sum restricts to \(\nu\)
at every coarse node. Upon sampling on the fine lattice it
defines \(E\nu\). In the original frame its column is
\(\chi(x-y)\operatorname{Ad}(W(\gamma_{yx})^{-1})\nu(y)\).
One may take the path \(\gamma_{yx}\) inside the support box
whenever its coefficient is nonzero.

The suprema of all continuous derivatives of the sum up to
order three are at most \(C\|\nu\|_\infty\), by disjointness
of the supports. A product of \(r\) fine differences of the
sample, \(r\leq2\), is the average of the corresponding
continuous derivative over \(r\) shifts in \([0,h]\).
Subtract the two such averages at \(x,x'\) to bound their
difference for \(r=2\) by \(C\|\nu\|_\infty|x-x'|\), using
the third derivative. This proves \eqref{eq:smooth-node-extension}
at every mesh, including \(h=1\), when sampling gives the
coarse nodal values themselves. The convergent matrix logarithm
and exponential give \eqref{eq:smooth-group-extension}, preserving
the deck law and the real form on the real chart.

At the flat reference all contour transports in the parallel frame
are the identity. The site row is therefore the arithmetic mean
\(qf(y)=h^4\sum_{x\in B^k(y)}f(x)\).
The only bump meeting that block is the bump at \(y\). Thus
\(qE=m_h I\), where \(m_h=h^4\sum_{x\in h\mathbb Z^4}\chi(x)\).
There are at least \((8h)^{-1}\) lattice sites in each coordinate
of \([-1/8,1/8]^4\), and at most \(C h^{-4}\) sites in the support.
This gives a positive lower bound and a finite upper bound for
\(m_h\), uniform in \(h\). Take \(E_q=m_h^{-1}E\).
For the energy bound, each sampled bump and each of its first
differences is bounded by a fixed constant. A difference of
a bump is supported on the fine links with an endpoint in its
support. Such links number at most \(C h^{-4}\) per bump,
and a given link meets the supports of at most two bumps.
Sum the squared values of the columns and use the measure \(h^4\).
This proves the last estimate in \eqref{eq:smooth-average-lift}.
\end{proof}

\begin{proposition}[Finite gauge in the strong norm]\label{prop:smooth-finite-gauge}
Let \(W\) be real and flat, fix \(0<\alpha<1\), and use the
parallel-frame norms \eqref{eq:smooth-gauge-norms} on a periodic mesh
\(0<h\le1\). For a positive-link \(\mathfrak g_{\mathbb C}\)-valued
potential \(A\) per unit length and a \(\mathfrak g_{\mathbb C}\)-valued
site field \(\lambda\), define
\(A^\lambda\) by
\((e^{hA}W)^{e^\lambda}=e^{hA^\lambda}W\), using the small
logarithm on each link. There is \(r_\alpha>0\) such that for
\(\|A\|_{X_\alpha}+\|\lambda\|_{Y_\alpha}<r_\alpha\)
this coordinate is well defined and holomorphic, and
\begin{equation}\label{eq:smooth-finite-gauge}
 \begin{split}
 A^\lambda&=A+D_W\lambda+\mathcal B_h(A,\lambda),\\
 \|\mathcal B_h(A,\lambda)\|_{X_\alpha}
 &\leq C_\alpha(\|A\|_{X_\alpha}+\|\lambda\|_{Y_\alpha})
                                        \|\lambda\|_{Y_\alpha}.
 \end{split}
\end{equation}
For two inputs with norm sums at most \(r<r_\alpha/2\),
\begin{equation}\label{eq:smooth-finite-lipschitz}
 \begin{split}
 &\|\mathcal B_h(A,\lambda)
                -\mathcal B_h(\widetilde A,\widetilde\lambda)\|_{X_\alpha}\\
 &\qquad\leq C_\alpha r\bigl(
      \|A-\widetilde A\|_{X_\alpha}
      +\|\lambda-\widetilde\lambda\|_{Y_\alpha}\bigr).
 \end{split}
\end{equation}
The constants and radius are uniform in \(0<h\leq1\), the periods
and the flat holonomies. The coordinate is real for real inputs.
Its value on a link uses its field value and the two endpoint
gauge values.
\end{proposition}

\begin{proof}
In the parallel frame, for \(b=(x,\mu)\) the exact formula is
\begin{equation}\label{eq:smooth-exact-log}
 A^\lambda_\mu(x)
 =h^{-1}\log\!\left(
       e^{\lambda(x)}e^{hA_\mu(x)}e^{-\lambda(x+h e_\mu)}
           \right).
\end{equation}
The product is precisely the gauge-transformed link in this frame.
Write \(p=\nabla_{h,\mu}\lambda\), so that the last exponent
in \eqref{eq:smooth-exact-log} is \(-\lambda(x)-hp(x)\).
This exposes the factor of \(h\) before differentiation.

For matrix fields \(f\) on the cover, a coordinate path of length
at most \(2|x-x'|\) gives
\([f]_{\alpha,\mathrm{loc}}\leq 2\|\nabla_h f\|_\infty\).
The product rule
\(\nabla_{h,\mu}(fg)=(\nabla_{h,\mu}f)(g(\cdot+h e_\mu))
                          +f\nabla_{h,\mu}g\)
and the product inequality for the local seminorm show that
\[
                    \|fg\|_{X_\alpha}
                    \leq C_\alpha\|f\|_{X_\alpha}\|g\|_{X_\alpha}.
\]
The same bound holds with translated factors. Also
\(\|\lambda\|_{X_\alpha}+\|p\|_{X_\alpha}\leq
C_\alpha\|\lambda\|_{Y_\alpha}\), by applying the path bound
to the first differences of \(\lambda\).

For one component, use the complex scalar \(\zeta\) in the product
\[
               B_\zeta=e^\lambda e^{\zeta A_\mu}
                                           e^{-\lambda-\zeta p}.
\]
Treat \(A_\mu,\lambda,p\) as fixed matrix fields. In the product
norm above, the exponential series converge and \(B_0=I\).
Termwise differentiation of the series bounds
\(\|\partial_\zeta B_\zeta\|_{X_\alpha}\) by
\(C_\alpha(\|A_\mu\|_{X_\alpha}+\|p\|_{X_\alpha})\)
for \(|\zeta|\leq2\) on a fixed small input ball.
Hence \(B_\zeta-I\) is bounded in that norm by this expression
times \(|\zeta|\). On a smaller fixed ball the logarithm series
converges in the product norm. Thus
\(\zeta^{-1}\log B_\zeta\) extends analytically through zero,
with \(X_\alpha\) bound
\(C_\alpha(\|A_\mu\|_{X_\alpha}+\|p\|_{X_\alpha})\).
Take \(\zeta=h\). The resulting small logarithms are traceless,
and are skew-Hermitian for real inputs. Conjugation equivariance
of the product and logarithm gives the deck law and the coordinate
in the original frame.

The bound on a fixed larger complex input ball gives uniform
second derivatives of the coordinate on a smaller ball. For
example, two applications of Cauchy's formula in scalar directions
of norm one give the bilinear derivative bound with radii fixed
by the separation of the two balls. The linearization at zero of
\eqref{eq:smooth-exact-log} is \(A_\mu-p\). At \(\lambda=0\)
the formula is exactly \(A_\mu\). Hence the nonlinear difference
\(\mathcal B_h\) has zero derivative at the origin and vanishes
at every \((A,0)\) in the ball. Integrate its \(\lambda\)-derivative
from zero to \(\lambda\) and use the uniform second derivative
bound. This gives \eqref{eq:smooth-finite-gauge}. On the
radius-\(r\) ball the entire first derivative of \(\mathcal B_h\)
is bounded by \(C_\alpha r\), which yields
\eqref{eq:smooth-finite-lipschitz} by integration on the segment
joining the inputs.
\end{proof}

In particular the actual infinitesimal gauge generator
\(\Gamma_A\lambda=\left.\frac{d}{dt}A^{t\lambda}\right|_{t=0}\)
satisfies
\[
 \|(\Gamma_A-D_W)\lambda\|_{X_\alpha}
                  \leq C_\alpha\|A\|_{X_\alpha}\|\lambda\|_{Y_\alpha}.
\]
This follows by differentiating the remainder with respect to
\(\lambda\) at zero, using its second derivative bound.

The variation of the covariant difference itself has order zero.
For \(U=e^{hA}W\), write \(X=\operatorname{Ad}(W(b))\lambda(b_+)\).
At a positive link \(b\), direct subtraction in \eqref{eq:gauge-D}
gives
\begin{equation}\label{eq:smooth-difference-variation}
 \begin{split}
 ((D_U-D_W)\lambda)(b)
  &=h^{-1}\bigl(I-\operatorname{Ad}(e^{hA(b)})\bigr)X\\
  &=-\int_0^1\operatorname{Ad}(e^{t h A(b)})[A(b),X]\,dt .
 \end{split}
\end{equation}
The integral follows by differentiation of the adjoint action.
It gives \(\|D_U-D_W\|_{\infty\to\infty}\leq C\|A\|_\infty\)
on a fixed potential ball. The same bound holds from physical
site \(L^2\) to physical link \(L^2\), both with measure \(h^4\):
for each direction the endpoint shift preserves the norm, and
there are four directions. In the flat parallel frame the product
estimate in \(X_\alpha\) also gives
\(\|(D_U-D_W)\lambda\|_{X_\alpha}\leq
C_\alpha\|A\|_{X_\alpha}\|\lambda\|_{X_\alpha}\) on a fixed strong
ball. For the derivatives, use
\(\operatorname{Ad}(e^{hA})=\exp(h\operatorname{ad} A)\), where
\(\operatorname{ad} A\) is the commutator with \(A\).
The \(m\)-th derivative, \(m\geq1\), of \(D_U-D_W\) in directions
\(A_1,\ldots,A_m\) has operator norm at most
\[
                   C_m h^{m-1}\prod_{i=1}^m\|A_i\|_\infty
\]
in the two indicated supremum and physical \(L^2\) operator norms.
Indeed the derivatives of the matrix exponential are bounded on
the fixed ball of \(h\operatorname{ad} A\), and the chain rule
gives one factor \(h\) per direction against the single \(h^{-1}\)
in \eqref{eq:smooth-difference-variation}. The same argument in
the product norm gives the \(X_\alpha\) operator bound with
\(X_\alpha\) norms of the directions. These estimates apply to the
holomorphic continuation.

\begin{proposition}[Correction to the effective site kernel]\label{prop:smooth-gauge-kernel}
Let \(W\) be a real flat reference on the periodic lattices of
Lemma~\ref{lem:smooth-node-extension}, with \(h=L^{-k}\), and fix
\(0<\alpha<1\). Let \(q=Q'_k(W)\) and let \(E_q\) be its smooth
right inverse in \eqref{eq:smooth-average-lift}. Denote by \(\tau_A\)
the effective site derivative in \eqref{eq:gauge-nonlinear-ward},
with \(a_0=hA\) and final depth \(k\).
For \(\|A\|_\infty<c_L\), the operator
\(E_{q,A}=E_q(\tau_A E_q)^{-1}\) is a right inverse of \(\tau_A\)
and satisfies the two norm bounds in \eqref{eq:smooth-average-lift},
with constants depending also on \(L\). For a fine site field
\(\lambda_0\in\ker q\), put
\begin{equation}\label{eq:smooth-gauge-kernel}
 \lambda_A=\lambda_0-E_{q,A}\tau_A\lambda_0 .
\end{equation}
Then \(\tau_A\lambda_A=0\) and
\(\|\lambda_A-\lambda_0\|_{Y_\alpha}\leq
C_{L,\alpha}\|A\|_\infty\|\lambda_0\|_{Y_\alpha}\).
The correction is holomorphic in \(A\) on this ball.
\end{proposition}

\begin{proof}
The one-step site estimate following \eqref{eq:gauge-site-one}
and the summable bounds of \cref{prop:avg-iteration} give
\(\|\tau_A-q\|_{\infty\to\infty}\leq C_L\|A\|_\infty\).
Here \(q\) is the product of the linear site rows at zero.
By \eqref{eq:smooth-average-lift}, \(\tau_AE_q\) therefore
differs from \(I\) by less than \(1/2\) on the stated ball
after decreasing \(c_L\). Its inverse is given by the convergent
Neumann series. The site factors in \eqref{eq:gauge-tau} read only
the one-step block in their site argument, so the product
\(\tau_A\) reads its site argument only on \(B^k(y)\) at \(y\).
Since the columns of \(E_q\) lie in their respective blocks,
\(\tau_AE_q\) and its inverse act separately at each coarse site.
The inverse is consequently bounded on coarse \(\ell^2\) as
well as on coarse supremum norm. Both lift estimates follow
from those of \(E_q\). Applying \(\tau_A\) to
\eqref{eq:smooth-gauge-kernel} gives zero, and
\(\tau_A\lambda_0=(\tau_A-q)\lambda_0\) gives the stated error.
The site formulas and the Neumann series give holomorphy.
\end{proof}

The extension converts the straightened mean into a prescribed
literal one without solving a gauge equation. In
\eqref{eq:gauge-factorization}, use \(a_0=hA\), write \(W_k=\mathcal M_{\rm an}^k(W)\)
and denote the top factor by \(v_k(A)\). Set
\begin{equation}\label{eq:smooth-literal-representative}
 \widehat U_A=(e^{hA}W)^{\mathfrak E(v_k(A)^{-1})},
 \qquad
 \mathcal M_{\rm an}^k(\widehat U_A)=e^{a_k(A)}W_k.
\end{equation}
Indeed the gauge on the coarse lattice is \(v_k(A)^{-1}\), which
cancels \(v_k(A)\) in the exact factorization. By
\cref{prop:avg-iteration} and \eqref{eq:smooth-node-extension}, the
fine gauge logarithm has \(Y_\alpha\) norm bounded by
\(C_{L,\alpha}\|A\|_\infty\).
Consequently the fine potential of \(\widehat U_A\) has
\(X_\alpha\) norm at most \(C_{L,\alpha}\|A\|_{X_\alpha}\).
On a smaller fixed \(X_\alpha\) ball both fine potentials are in
the admitted averaging charts. Analytic covariance gives the coarse
identity on this ball. For real fields the iteration's first-branch
admission also gives the identity for the Borel map. The coefficient at a fine link reads the
link and the backward blocks of the nodes whose bumps meet its
endpoints, a region of bounded diameter in coarse units.

% End rows.tex
% Begin fibres.tex
\section{Local variations and changes of frame}\label{sec:local-variations}

\subsection{Compatible local frames}

Nearby real $\operatorname{SU}(2)$ frames can be combined without
leaving the group. Suppose $g_1,\ldots,g_N\in G$ obey
$\|g_i g_1^{-1}-I\|<r<1/2$, and let $\theta_i\ge0$ with
$\sum_i\theta_i=1$. Put $X=\sum_i\theta_i g_i$. Then
\begin{equation}\label{eq:frame-interpolation}
                 g=X(\det X)^{-1/2}\in G,
\end{equation}
where the square root is positive on these real inputs.
In the quaternion representation $X$ is a real quaternion matrix,
so $X^*X=(\det X)I$ and $\det X>0$. Moreover
$\|Xg_1^{-1}-I\|<r$, which prevents $X=0$ and gives a fixed
margin for normalization. This proves the assertion. On a smaller
complex neighborhood the determinant avoids the nonpositive real axis.
Take its principal square root, continuing the positive real value.
Formula \eqref{eq:frame-interpolation}
is then holomorphic, has determinant one, and commutes with common
left and right multiplication by elements of $G_{\mathbb C}$.
For position-dependent weights the assertion applies pointwise;
derivative estimates must also include derivatives of the weights.
It asserts no global gluing when the frames are not close on overlaps.

There is a simple algebraic readout at a retained junction. If a
link with endpoints $y,z$ has actual value $V$ and reduced value
$Z$, and its node factors obey $V=n(y)Z n(z)^{-1}$, then
\begin{equation}\label{eq:junction-node-readout}
                         n(z)=V^{-1}n(y)Z.
\end{equation}
Thus a known factor at one endpoint and the retained link determine
the other factor. This identity is exact. Near-identity or complex
chart admission for its factors is a separate hypothesis.

The normalized axial gauge need not have a bounded fine-scale
derivative. On a flat one-color field $U(x,\mu)=e^{hA_\mu}$,
take only $A_1=X\ne0$, with $X$ fixed in the small logarithm chart,
and a block of physical side one. In its
parallel scalar direction the normalized gauge is
$g(x)=e^{(x_1-y_1)X}$, since the centred mean of $x_1-y_1$ is
zero. It makes internal transports trivial. Across the face of
two blocks its logarithm has jump magnitude $(1-h)|X|$ in the
branch at zero.
The difference divided by $h$ is therefore unbounded as $h\downarrow0$.
The smooth extension in Section~\ref{sec:smooth-node-gauge} has
different defining conditions.

The comparison with coordinate combs uses endpoint covariance,
inversion of paths, bounded path length, and a plaquette filling
of the comparison loops inside the complete carrier. These give
the one-step analytic bounds and their summable composition.
A single prescribed coordinate order does not have the full
signed-coordinate symmetry of the contour mean. Its global
Borel completion also requires its own definition; neither that
completion nor an inverse-operator theorem follows from equality
of the flat straight rows.

\subsection{The local second variation of the Wilson action}

The geometric second variation is useful before a gauge condition or a
constraint is imposed. Fix a real connection $U$ on a periodic fine
lattice of spacing $h$. An increment per unit length $A$ acts by
$U_A(b)=e^{hA(b)}U(b)$ on positive links. On a reversed link use
$A(\bar b)=-\operatorname{Ad}(U(b)^{-1})A(b)$, so that this convention
also describes the inverse of the varied positive link.
For an oriented plaquette with edges $b_1,\ldots,b_4$, write
\begin{equation}\label{eq:wilson-insertions}
 P=U(b_1)\cdots U(b_4),\qquad
 T_{r-1}=U(b_1)\cdots U(b_{r-1}),\qquad
 X_r=h\operatorname{Ad}(T_{r-1})A(b_r).
\end{equation}
The empty prefix is $I$. Define $Y_r$ by replacing $A$ with $B$.

\begin{lemma}[Plaquette Hessian]\label{lem:wilson-local-hessian}
Let \(U\) be a real connection on a periodic lattice of spacing
\(h>0\), and let \(A,B\) be real link increments per unit length,
extended by \(A(\bar b)=-\operatorname{Ad}(U(b)^{-1})A(b)\)
and the same rule for \(B\). For an oriented plaquette
\(p\), use the prefixes and insertions in \eqref{eq:wilson-insertions}.
The Hessian of \(a_p(U)=1-\tfrac12\operatorname{tr}P\) in the
coordinates \(U_A=e^{hA}U\) satisfies
\begin{equation}\label{eq:wilson-local-hessian}
 \begin{split}
 D^2a_p(U)[A,B]
 &=\left\langle\sum_rX_r,\sum_rY_r\right\rangle+R_p(A,B),\\
 |R_p(A,B)|&\le C\|P-I\|
       \left(\sum_r|X_r|^2\right)^{1/2}
       \left(\sum_r|Y_r|^2\right)^{1/2}.
 \end{split}
\end{equation}
Here $\langle X,Y\rangle=-\tfrac12\operatorname{tr}(XY)$.
The constant is absolute. The trace formula extends holomorphically
to the corresponding complex link chart.
\end{lemma}

\begin{proof}
The varied plaquette is
$\prod_{r=1}^4e^{sX_r+tY_r}P$. Its mixed derivative at zero is
\[
 H=\sum_r\tfrac12(X_rY_r+Y_rX_r)
       +\sum_{r<q}(X_rY_q+Y_rX_q),
\]
followed by multiplication by $P$. Cyclicity gives
\[
 \operatorname{tr}H=\operatorname{tr}
                 \left[\left(\sum_rX_r\right)\left(\sum_rY_r\right)\right].
\]
Thus $R_p=-\tfrac12\operatorname{tr}(H(P-I))$.
The trace estimate
$|\operatorname{tr}(CD)|\le2\|C\|\|D\|$, followed by
Cauchy--Schwarz on the four indices, proves the bound. The
formula is a polynomial in the insertions and the plaquette matrix;
it involves no complex conjugation.
\end{proof}

Define the transported plaquette curl
\[
 (\mathcal C_UA)(p)=h^{-1}\sum_{r=1}^4
                  \operatorname{Ad}(T_{r-1})A(b_r).
\]
Suppose each plaquette has an assigned length $\xi_p\ge h$ and
$\|P-I\|\le\varepsilon h^2\xi_p^{-2}$. Assign a length $\xi_b$
to each positive link and require $\xi_p^{-2}\le C_0\xi_b^{-2}$
whenever $b$ belongs to $p$. Set
$\|A\|_\xi^2=\sum_bh^4\xi_b^{-2}|A(b)|^2$ and
$\|F\|_2^2=\sum_ph^4|F(p)|^2$. Summing the lemma gives
\begin{equation}\label{eq:wilson-summed-hessian}
 \left|D^2\sum_pa_p(U)[A,B]
       -\langle\mathcal C_UA,\mathcal C_UB\rangle_2\right|
       \le C C_0\varepsilon\|A\|_\xi\|B\|_\xi .
\end{equation}
Indeed $\sum_r|X_r|^2=h^2\sum_{b\in\partial p}|A(b)|^2$ by
unitarity, and a positive link has six plaquette incidences in four
dimensions, counted with multiplicity on a short torus. The two
factors of $h^2$ give exactly the physical measure $h^4$.

For comparison put
\[
 (\nabla^+_{U,\mu}A_\nu)(x)
 =h^{-1}\bigl[\operatorname{Ad}(U_\mu(x))A_\nu(x+he_\mu)
                         -A_\nu(x)\bigr],\qquad
 (\mathrm c_UA)_{\mu\nu}=\nabla^+_{U,\mu}A_\nu
                              -\nabla^+_{U,\nu}A_\mu.
\]
Multiplication of the four prefixes, with the reverse-link convention,
gives the exact identity
\begin{equation}\label{eq:wilson-curl-comparison}
 h(\mathcal C_UA-\mathrm c_UA)_{\mu\nu}(x)
 =(I-\operatorname{Ad}P)
     \bigl[\operatorname{Ad}(U_\nu(x))A_\mu(x+he_\nu)+A_\nu(x)\bigr].
\end{equation}
For unitary $P$, $\|I-\operatorname{Ad}P\|\le2\|P-I\|$.
Using $h\le\xi_p$ and the same incidence count yields
$\|\mathcal C_UA-\mathrm c_UA\|_2\le C C_0^{1/2}
\varepsilon\|A\|_\xi$. Consequently, for $0<\varepsilon<1/2$,
\begin{equation}\label{eq:wilson-curl-lower}
 D^2\sum_pa_p(U)[A,A]
 \ge(1-C\varepsilon)\|\mathrm c_UA\|_2^2
                          -C(1+C_0)\varepsilon\|A\|_\xi^2.
\end{equation}
This follows by expanding the curl square in
\eqref{eq:wilson-summed-hessian} and using
$2ab\le\varepsilon a^2+\varepsilon^{-1}b^2$.
For positive plaquette weights \(w_p\) fixed as the connection varies,
use \(\|F\|_{2,w}^2=\sum_p h^4w_p|F(p)|^2\). The same pointwise
remainders and incidence count give
\[
 D^2\sum_p w_pa_p(U)[A,A]
 \ge (1-C\varepsilon)\|\mathrm c_UA\|_{2,w}^2
       -C\|w\|_\infty(1+C_0)\varepsilon\|A\|_\xi^2.
\]
Indeed the summed Hessian has leading term
\(\|\mathcal C_UA\|_{2,w}^2\), and the squared weighted curl
error is at most \(C\|w\|_\infty C_0\varepsilon^2\|A\|_\xi^2\).
The same inequality \(2ab\le\varepsilon a^2+\varepsilon^{-1}b^2\)
therefore gives the displayed bound.
These are local Wilson estimates. They do not identify a constrained
Hessian, include multiplier terms, or prove coercivity after a gauge
projection.

\paragraph{Covariant curl and divergence.}
For a real link field $A$ on the full periodic graph, let
\(S_\mu f(x)=\operatorname{Ad}(U_\mu(x))f(x+he_\mu)\)
and \(\delta_\mu=h^{-1}(S_\mu-I)\). Thus
\(\mathrm c_U A\) has components
\(\delta_\mu A_\nu-\delta_\nu A_\mu\),
and \(D_U^\ast A=-\sum_\mu\delta_\mu^\ast A_\mu\) for the
physical adjoint. Each \(S_\mu\) is unitary. Expanding the two
squares and subtracting the component-gradient square gives
\begin{equation}\label{eq:wilson-hodge-commutator}
 \begin{split}
 &\|\mathrm c_U A\|_2^2+\|D_U^\ast A\|_2^2
                        -\sum_{\mu,\nu}\|\delta_\mu A_\nu\|_2^2\\
 &\qquad=2\operatorname{Re}\sum_{\mu<\nu}
       \langle A_\mu,[\delta_\mu,\delta_\nu^\ast]A_\nu\rangle_2 .
 \end{split}
\end{equation}
Indeed the diagonal terms
\(\|\delta_\mu^\ast A_\mu\|_2^2\) and
\(\|\delta_\mu A_\mu\|_2^2\) agree because
\(\delta_\mu\) is normal. In a pair of distinct directions the
divergence contributes the product \(\delta_\mu\delta_\nu^\ast\)
and the curl contributes its reversed order with negative sign.

The commutator is \(h^{-2}[S_\mu,S_\nu^{-1}]\).
Its two paths from \(x\) to \(x+he_\mu-he_\nu\) bound the
backward plaquette in those directions. Their adjoint transports
differ in norm by at most twice the plaquette's group defect,
since both are unitary and their ratio is a based plaquette
holonomy or its inverse. The two link values paired in
\eqref{eq:wilson-hodge-commutator} belong to that plaquette.
Under the inverse-length comparison assumed in
\eqref{eq:wilson-summed-hessian}, each such term is bounded
after summation by \(C C_0\varepsilon\|A\|_\xi^2\).
Here \(h^{-2}\) cancels the \(h^2\) in the curvature bound,
and the finite incidence count and Cauchy--Schwarz give
\begin{equation}\label{eq:wilson-hodge-bound}
 \left|\|\mathrm c_U A\|_2^2+\|D_U^\ast A\|_2^2
            -\sum_{\mu,\nu}\|\delta_\mu A_\nu\|_2^2\right|
       \leq C C_0\varepsilon\|A\|_\xi^2 .
\end{equation}
On an inner graph the adjoint depends on the boundary convention,
so the periodic calculation requires separate boundary terms there.


\subsection{Local curvature coefficients}

Take \(0<h\leq1\) in fixed physical units. In the flat parallel
frame of \cref{sec:smooth-node-gauge}, write
\(U_\mu(x)=e^{hA_\mu(x)}\) for the perturbation \(e^{hA}W\).
The forward difference \(\nabla_{h,\mu}\) divides the step in
direction \(\mu\) by \(h\), and
\((\mathrm d_h A)_{\mu\nu}=\nabla_{h,\mu}A_\nu-\nabla_{h,\nu}A_\mu\).
All link components refer to the positive orientations. The
parallel frame allows these formulas on a periodic lattice with
nontrivial flat holonomy.

\begin{lemma}[Curvature expansion]\label{lem:wilson-curvature-expansion}
Let \(0<h\le1\) and let \(W\) be a real flat periodic connection.
Work in its parallel frame on the cover, with the boundary
identifications determined by its flat holonomy. For a link
potential \(A\) on a fixed small complex supremum ball set
\(U=e^{hA}W\) and \(F_h(A)(p)=h^{-2}\log U(\partial p)\) in
that frame. Then
\begin{equation}\label{eq:wilson-curvature-expansion}
 F_h(A)=\mathrm d_h A+R_h(A),\qquad
 \|R_h(A)\|_\infty\le C\|A\|_\infty^2,\qquad
 \|\nabla_hR_h(A)\|_\infty
            \le C\|A\|_\infty\|\nabla_h A\|_\infty .
\end{equation}
Here \((\mathrm d_hA)_{\mu\nu}=\nabla_{h,\mu}A_\nu-
\nabla_{h,\nu}A_\mu\), and \(\nabla_h\) uses forward
differences divided by \(h\). On plaquette fields it compares
translates of the same orientation. The analytic ball and the
constant are uniform in the mesh, volume and real flat holonomy.
\end{lemma}

\begin{proof}
At \(p=(x,\mu,\nu)\) the ordered link logarithms are \(ha_1,\ldots,ha_4\),
where
\[
 (a_1,a_2,a_3,a_4)
 =\bigl(A_\mu(x),A_\nu(x+he_\mu),
                      -A_\mu(x+he_\nu),-A_\nu(x)\bigr).
\]
Their sum divided by \(h\) is \((\mathrm d_h A)(p)\).
For a complex parameter \(z\), the logarithm
\(\log\prod_{j=1}^4 e^{za_j}\) is analytic for \(|z|\leq2\) on
a sufficiently small fixed ball of four-tuples. Its value and
first derivative at \(z=0\) are \(0\) and \(\sum_j a_j\).
Subtract the linear term and divide by \(z^2\). The result extends
analytically at zero and is bounded on the fixed ball for
\(|z|\leq1\). Its value and differential in the four variables
vanish at their common origin. Cauchy's formula and Taylor's
formula give a bound \(C\max_j|a_j|^2\) for the result and
\(C\max_j|a_j|\) for its differential, uniformly in \(z\).
Use \(z=h\) for the curvature. Between two neighboring plaquettes
of the same orientation the four inputs change by at most
\(h\|\nabla_h A\|_\infty\). Integrating the derivative bound on
the segment of inputs and dividing by \(h\) proves the last
estimate. The same calculation bounds the first and second
derivatives of \(R_h\) in supremum norm by \(C\|A\|_\infty\)
and \(C\), respectively. The first derivative of \(F_h\) in
direction \(B\) is therefore bounded by
\(C(\|\nabla_h B\|_\infty+\|A\|_\infty\|B\|_\infty)\).
\end{proof}

For a plaquette logarithm \(h^2F\), the odd Wilson factor in the
skew-Hermitian convention is
\[
 \mathcal S_h(F)=\frac{e^{h^2F}-e^{-h^2F}}{2h^2}.
\]
It is the matrix sine in the Hermitian convention multiplied by \(i\).
Write \(\mathcal S_h(F)=F+S_h(F)\). For \(\|h^2F\|_\infty\) bounded
by a fixed constant, its series gives
\begin{equation}\label{eq:wilson-sine-remainder}
 \|S_h(F)\|_\infty\leq C h^4\|F\|_\infty^3,\qquad
 \|\nabla_h S_h(F)\|_\infty\leq C h^3\|F\|_\infty^3 .
\end{equation}
For the second inequality subtract two values of the remainder
and use the first inequality at both plaquettes before dividing
by \(h\).

\paragraph{The exact Wilson current.}
For a real connection \(U\) on the full periodic lattice, let
\(J_{\mathrm{left}}(U)\) represent the first variation of the Wilson
energy from Section~\ref{sec:energy}:
\[
 \left.\frac{d}{dt}\mathcal E(e^{thB}U)\right|_{t=0}
                  =\langle J_{\mathrm{left}}(U),B\rangle_2.
\]
Here \(B\) is a real positive-link field per unit length.
Use the bilinear colour pairing with weight \(h^4\) on both links
and plaquettes, and let \({}^{\sharp}\) denote its transpose.
For \(P_p=U(\partial p)\), put
\(\mathcal S_U(p)=(P_p-P_p^{-1})/(2h^2)\), a
\(\mathfrak g\)-valued plaquette field. Then
\begin{equation}\label{eq:wilson-exact-current}
 J_{\mathrm{left}}(U)
    =\mathcal C_U^\sharp\mathcal S_U
    =\mathrm c_U^\sharp\mathcal S_U.
\end{equation}
When \(P_p=e^{h^2F(p)}\), this factor is
\(\mathcal S_U(p)=\mathcal S_h(F(p))\).

To prove the identity, the insertion formula
\eqref{eq:wilson-insertions} gives
\(P_p'=h^2(\mathcal C_UB)(p)P_p\).
The Cayley--Hamilton identity
\(P_p+P_p^{-1}=(\operatorname{tr}P_p)I\) and the tracelessness
of \(\mathcal C_UB\) therefore give
\[
 \left.\frac{d}{dt}\mathcal E(e^{thB}U)\right|_{t=0}
      =h^4\sum_p\langle(\mathcal C_UB)(p),\mathcal S_U(p)\rangle.
\]
This proves the first equality in \eqref{eq:wilson-exact-current}.
For the second, pair \eqref{eq:wilson-curl-comparison} with
\(\mathcal S_U\). Its right side pairs to zero because
\(\operatorname{Ad}(P_p^{-1})\mathcal S_U(p)=\mathcal S_U(p)\).
The same algebra proves the identity on complex link charts,
using the bilinear transpose. No plaquette logarithm is needed
for the formula in terms of \(\mathcal S_U\).

The derivative in normal coordinates at a nonzero potential
uses the exponential differential \(E_z\) from
\eqref{eq:gauge-tau}. Its transpose for the invariant bilinear
colour pairing, extended over \(\mathbb C\), is
\begin{equation}\label{eq:wilson-exponential-differential}
 E_z^\sharp B=\int_0^1\operatorname{Ad}(e^{-t z})B\,dt .
\end{equation}
This follows by moving the adjoint action across the trace
pairing inside the defining integral. The normal-coordinate
variation \(A+sB\) gives left variation \(E_{hA}B\) at \(U=e^{hA}W\)
when both variations are measured per unit length. Consequently,
\eqref{eq:wilson-exact-current} gives the Wilson current in normal
coordinates as
\[
 J_{\mathrm{normal}}(A)
       =E_{hA}^\sharp\mathrm c_U^\sharp\mathcal S_U,
       \qquad U=e^{hA}W.
\]
The same transpose rule converts any differential represented
in the left variable to the normal variable.
The integral formulas give
\[
 \|E_{hA}-I\|+\|E_{hA}^\sharp-I\|\leq C h\|A\|_\infty
\]
on a fixed potential ball, and their first potential derivatives
are bounded by \(C h\) in the corresponding supremum operator norm.

The transported plaquette curl \(\mathcal C_U\) also has local
background derivatives. Treat \(U=e^{hA}W\), with \(W\) fixed
and real. In each coefficient of \(\mathcal C_U\) the adjoint
transport is a product of at most four factors associated with
links or their inverses. For a reversed insertion, the fourth
factor may be the inverse of the positive link at which the
input is stored. With \(U=e^{hA}W\), each varying adjoint factor
is an exponential \(\exp(\pm h\operatorname{ad}A(b))\),
possibly preceded or followed by a fixed unitary adjoint action.
Its first and second derivatives have factors \(h\) and \(h^2\).
The prefactor \(h^{-1}\) in the curl leaves the bounds
\begin{equation}\label{eq:wilson-curl-coeff-jets}
 \begin{split}
 \|\mathcal C_U-\mathcal C_W\|&\leq C\|A\|_\infty,\\
 \|D_A\mathcal C_U[A_1]\|&\leq C\|A_1\|_\infty,\\
 \|D_A^2\mathcal C_U[A_1,A_2]\|
                         &\leq Ch\|A_1\|_\infty\|A_2\|_\infty .
 \end{split}
\end{equation}
The first estimate follows by integration from \(A=0\).
These are coefficient bounds and also operator bounds from links
to plaquettes in supremum norm and in the physical \(L^2\) norms
with measure \(h^4\). Each row involves at most four links, and each
link occurs in at most six plaquette positions, counted with
multiplicity. The same proof applies to complex \(A\) in the fixed
small ball, since the varying factors are uniformly bounded there.
For flat \(W\), \(\mathcal C_W=\mathrm d_h\) in its parallel frame.

\begin{lemma}[Quadratic Wilson current remainder]\label{lem:wilson-current-remainder}
Let \(W\) be real and flat on the full periodic fine lattice, with
\(0<h\le1\), and use its parallel frame with the deck law.
For \(0<\alpha<1\), take \(X_\alpha\) from
\eqref{eq:smooth-gauge-norms}. The normal current of \(e^{hA}W\)
has a holomorphic remainder
\[
 \mathcal N_h(A)=J_{\mathrm{normal}}(A)-\mathrm d_h^\sharp\mathrm d_h A
\]
on a fixed small complex \(X_\alpha\) ball. There are \(r_0>0\)
and \(C\) such that, for \(0<r\le r_0\) and
\(\|A\|_{X_\alpha},\|\widetilde A\|_{X_\alpha}\le r\),
\begin{equation}\label{eq:wilson-current-remainder}
 \|\mathcal N_h(A)\|_\infty\le C\|A\|_{X_\alpha}^2,\qquad
 \|\mathcal N_h(A)-\mathcal N_h(\widetilde A)\|_\infty
       \le Cr\|A-\widetilde A\|_{X_\alpha}.
\end{equation}
The radii and constants are uniform in the mesh, volume and real
flat holonomy.
\end{lemma}

\begin{proof}
Put \(U=e^{hA}W\). The flat curl transpose is a sum of backward differences in this
frame. The formula for the forward covariant curl shows that
\(\|\mathrm c_U^\sharp-\mathrm d_h^\sharp\|_{\infty\to\infty}
\le C\|A\|_\infty\). Each translated coefficient changes by a
factor of the form \(h^{-1}(\operatorname{Ad}(e^{-hA(b)})-I)\).
Write \(F=F_h(A)\) and \(R=R_h(A)\). The exact current identity gives
\[
 J_{\mathrm{left}}(U)-\mathrm d_h^\sharp\mathrm d_h A
   =\mathrm d_h^\sharp R
      +(\mathrm c_U^\sharp-\mathrm d_h^\sharp)F
      +\mathrm c_U^\sharp S_h(F).
\]
The first two terms are bounded by \(C\|A\|_{X_\alpha}^2\) by
\eqref{eq:wilson-curvature-expansion}. Split the last transpose
as its flat value plus the same coefficient difference.
Equation~\ref{eq:wilson-sine-remainder} bounds this term by
\(C(h^3+h^4\|A\|_\infty)\|F\|_\infty^3\), which is absorbed on
the small ball.

The exact normal-coordinate factor is \(E_{hA}^\sharp\).
It is uniformly bounded, and its difference from \(I\) is at
most \(Ch\|A\|_\infty\). The difference applied to
\(\mathrm d_h^\sharp\mathrm d_h A\) is therefore bounded by
\(C\|A\|_\infty\|\nabla_h A\|_\infty\), using
\(\|\mathrm d_h^\sharp\mathrm d_h A\|_\infty
\le Ch^{-1}\|\nabla_h A\|_\infty\). This proves the quadratic
bound. All these expressions and estimates hold on a fixed
complex ball. Cauchy's formula in two scalar directions bounds
the second derivative on a smaller ball. The quadratic bound
gives \(D\mathcal N_h(0)=0\); integrate the derivative along
segments in the radius-\(r\) ball to obtain the Lipschitz estimate.
\end{proof}

\subsection{Averaging penalty ranges}

Choose a real reference \(U\) admitting the row charts. Let \(q\)
and \(Q_U^{\rm st}\) be its transported site row and its straightened link
derivative per unit length at zero perturbation. Denote their
coarse site and link index sets by \(\mathcal Y\) and \(\mathcal L\).
Use the transposes for the physical pairings and write \(D=D_U\).
For the scalar and vector penalties
\[
               V_s=q^\sharp q,\qquad
               V_v=(Q_U^{\rm st})^\sharp Q_U^{\rm st}-Dq^\sharp qD^\sharp
\]
one has
\begin{equation}\label{eq:wilson-penalty-rank}
 \operatorname{rank} V_s\leq 3|\mathcal Y|,\qquad
 \operatorname{rank} V_v\leq 3(|\mathcal Y|+|\mathcal L|).
\end{equation}
Indeed, the scalar range lies in the range of \(q^\sharp\), and
the vector range lies in the sum of the ranges of \((Q_U^{\rm st})^\sharp\)
and \(Dq^\sharp\). Each coarse coordinate has three colour
components. The argument also applies to the bilinear complex
transposes. In a translation-invariant coarse Fourier fibre with
one site and four positive links per cell the two bounds are \(3\)
and \(15\).

\subsection{A constant Abelian curvature}

On the covering lattice, fix \(F\in\mathfrak g\) and \(t\in\mathbb R\).
In the \(1,2\) plane consider the two potentials
\[
 \begin{array}{c|cc}
 & A_1(x)&A_2(x)\\ \hline
 \text{linear}&-t x_2F&0\\
 \text{symmetric}&-\tfrac12 t x_2F&\tfrac12 t x_1F .
 \end{array}
\]
The other two components vanish. All link logarithms lie on the
same Abelian line. With \(g(x)=\exp(-t x_1x_2F/2)\), the finite
gauge action on the linear links \(e^{hA_\mu(x)}\) gives exactly
the symmetric links, since the commuting site difference adds
the forward difference of \(t x_1x_2F/2\) to each potential.
In each gauge the \(1,2\) plaquette holonomy is \(e^{h^2tF}\) by
multiplication of the four commuting link factors. The factor
\(h^2\) is the plaquette area. The adjoint action of \(g\) at the site or link tail preserves
every fixed coordinate support. A covariant endomorphism of
the site or link space therefore transforms by unitary
conjugation, as does its compression to such a support.

There is one plaquette of this orientation per cell of volume
\(h^4\). Since the eigenvalues of \(F\) are \(\pm i|F|\), the
Wilson contribution per unit physical volume is
\begin{equation}\label{eq:wilson-constant-curvature}
 h^{-4}\bigl(1-\cos(h^2t|F|)\bigr)
 =\tfrac12t^2|F|^2+O\bigl(h^4t^4|F|^4\bigr).
\end{equation}
In particular the quadratic coefficient retains its normalization
when converting plaquette angles to physical curvature.

\section{Coordinates on a fixed-average fibre}\label{sec:fibre-chart}

Use units in which the fine spacing of the step is one. Let $E$ be a finite
set of positive fine bonds and $C$ a set of positive coarse bonds whose
complete averaging paths are contained in $E$ with either orientation. For
each $c=(y,y+Le_\mu)\in C$, write $D_c$ for the rectangular patch formed by
its two blocks. On a torus we assume that the vertices of each such rectangle
have distinct projections. We use the bonds of its rectangular lift for all
path calculations, including when the chosen rectangle crosses a periodic
seam. The path support of row $c$ will be denoted $E_c\subset E$. In
particular, $E_c$ need not contain every physical bond joining two projected
vertices of $D_c$.

Fix a background $V$ on $E$. For each $c$, suppose there is a unitary
gauge $s_c$ on the rectangle such that
\begin{equation}\label{eq:fibre-background}
 \max_{b\subset D_c}\|V^{s_c}(b)-I\|\le\epsilon.
\end{equation}
The gauges for different rows can be chosen separately. By the local comb
estimate of Lemma~\ref{lem:avg-comb}, a real field with plaquette defect at
most $\omega$ in these rectangles admits \eqref{eq:fibre-background} with
$\epsilon\le C L\omega$. For a complex perturbation $V=e^XV_0$ of such a real
field, the same real gauges give $\epsilon\le C(L\omega+\|X\|_\infty)$ when
the right side is small. Coarse links in their original frames can have
nontrivial holonomy.

Use the analytic map \(\mathcal M_{\rm an}\) from \eqref{eq:avg-map}
on the paths of these rows. Put
$W=\mathcal M_{\rm an}(V)|_C$ and use the corresponding plain coordinate
on the rows in $C$. Write
\begin{equation}\label{eq:fibre-constraint}
 P_V(A)(c)=\log\big[\mathcal M_{\rm an}(e^A V)(c)W(c)^{-1}\big],\qquad
 K_V=DP_V(0),\qquad N_V(A)=P_V(A)-K_VA.
\end{equation}
Here $A$ takes values in $\mathfrak g_\mathbb C$, or in $\mathfrak g$ on the
real chart, and each exponential multiplies its reference link on the left.
Thus $P_V(0)=0$ throughout the analytic reference domain. All averaging
expressions in this section use $\mathcal M_{\rm an}$ until the real
admission condition below is imposed; on that real chart they agree with
the global Borel map $\mathcal M$. Norms of fields are suprema of the
matrix operator norm.
In this section $C_0,C_1,\ldots$ and $c_0,c_1,\ldots$ are positive constants
independent of the lattice, the row selection, and $L\ge3$. Constants denoted
$C_L$ may depend on $L$.

For a chart of the real Borel map, suppose also that $V$ is defined on every
link used in the first-branch plaquette tests and that
\begin{equation}\label{eq:fibre-real-admission}
 \max_{p\,\text{tested by } C}\|V(\partial p)-I\|
 \le \frac1{400L^2}.
\end{equation}
A test on a projected double block can include plaquettes outside its
rectangular lift. Links of those plaquettes outside $E$ are held fixed.
When every tested plaquette lies in the rectangular lift,
\eqref{eq:fibre-background} gives a defect at most \(4\epsilon\)
on each test, by unitary telescoping. Thus
\eqref{eq:fibre-real-admission} follows on this stock for a
sufficiently small \(\epsilon\) of order \(L^{-5}\).

\subsection{The local Taylor bounds and the pivot coefficient}\label{sec:fibre-pivot}

\begin{lemma}[Taylor bounds in physical bond logarithms]\label{lem:fibre-taylor}
If \eqref{eq:fibre-background} holds in unitary local gauges
\(s_c\) with $\epsilon\le c_0/L$, the expression
\eqref{eq:fibre-constraint} using \eqref{eq:avg-map} is holomorphic for
$\|A\|_\infty<c_0/L$. Each row depends only on $A|_{E_c}$ and $V|_{E_c}$. On
this ball,
\begin{equation}\label{eq:fibre-remainder}
 \|K_V\|\le C_0 L,\qquad
 \|N_V(A)\|_\infty\le C_0 L^2\|A\|_\infty^2,\qquad
 \|DN_V(A)\|\le C_0 L^2\|A\|_\infty.
\end{equation}
In a local unitary frame as in \eqref{eq:fibre-background}, the row of
$K_{V^{s_c}}-K_I$ has operator norm at most $C_0 L^2\epsilon$ on the local
supremum space.
\end{lemma}

\begin{proof}
In the frame $s_c$ put $V^{s_c}=e^Z$, with $\|Z\|_\infty\le 2\epsilon$. A
product of at most $2L$ factors $e^{\pm Z_b}$ or $(e^{A_b}e^{Z_b})^{\pm1}$ is
within $\exp(C_1 L(\|Z\|_\infty+\|A\|_\infty))-1$ of $I$. This follows by
telescoping the product and bounding each factor by its exponential norm
bound. It puts every internal contour family in the fixed local domain of the
contour mean. In \eqref{eq:avg-map} the subsequent operations are products of
a bounded number of such near-identity matrices, local logarithms, and a
normalized sum. Hence \(P_{e^Z}(A)(c)\) is bounded by a universal constant on
joint balls in \(Z,A\) of radius \(c_1/L\).

Apply the scalar Cauchy formula on a smaller ball in two variables,
along directions of supremum norm at most one. The circles can
have radius \(c_2/L\). This bounds \(D_A^2P_{e^Z}\) and
\(D_ZD_AP_{e^Z}\) by \(C_2L^2\). One circle bounds \(D_AP_{e^Z}\)
by \(C_2L\). Integration from zero to \(A\) and from zero to \(Z\)
gives the stated bounds and the comparison with \(K_I\).

By gauge covariance of the average, a left fluctuation transforms by
$A^{s_c}(b)=\operatorname{Ad}(s_c(b_-))A(b)$. The relative logarithm
transforms by $\operatorname{Ad}(s_c(y))$. These conjugations have norm one.
The holomorphic formulas therefore give the same estimates in the original
frame. This also proves locality, since every operation used only the path
variables in the row.
\end{proof}

The pivot of $c=(y,y+Le_\mu)$ is the positive fine bond
\[
 b_c=(y+me_\mu,y+(m+1)e_\mu).
\]
It is the central bond crossing from the first block to the second. An
internal contour never crosses a block face. A parallel segment from $x$ to
$x+Le_\mu$, $x\in B(y)$, crosses this face exactly once. Thus the only pivot
in $E_c$ is $b_c$. This path argument also applies when two opposite faces
have the same pair of neighbouring coarse sites on a short periodic lattice.

Let $P=\{b_c:c\in C\}$ and $E_f=E\setminus P$. The map $c\mapsto b_c$ is
injective, since the direction and initial point determine $c$. Write $\pi$
for restriction to $E_f$ and $\iota$ for extension from $E_f$ by zero. At the
identity, only the parallel segments in the flat derivative
\eqref{eq:avg-flat-link} can contribute a pivot derivative. Exactly $L$ of
them contain $b_c$, one for each longitudinal starting coordinate with fixed
transverse coordinates $y_\nu$ for $\nu\ne\mu$. Each contributes with positive
sign and weight $L^{-4}$. Therefore, with
\begin{equation}\label{eq:fibre-pivot-matrix}
 J_c(V)=D_{A(b_c)}P_V(0)(c),\qquad J_c(I)=L^{-3}I_{\mathfrak g_\mathbb C},
\end{equation}
the restriction of $K_V$ to $P$ is diagonal in the bond index. In a flat field
on the rectangle the coefficient in the original frame is
$L^{-3}\operatorname{Ad}(V[y,b_{c,-}])$, where $V[y,b_{c,-}]$ is the prefix
along the central segment. This follows from \eqref{eq:fibre-pivot-matrix} by
the gauge change in Lemma~\ref{lem:fibre-taylor}.

\begin{lemma}[Right inverse on the pivots]\label{lem:fibre-pivot}
Under \eqref{eq:fibre-background} in unitary local gauges
\(s_c\), with $\epsilon\le c_3L^{-5}$, each $J_c(V)$
is invertible. The maps
\begin{equation}\label{eq:fibre-linear-lift}
 (H_V v)(b_c)=J_c(V)^{-1}v(c),\qquad
 (H_V v)|_{E_f}=0,\qquad \mathsf C_V=\iota-H_VK_V\iota
\end{equation}
satisfy $K_VH_V=I$, $K_V\mathsf C_V=0$, $\pi\mathsf C_V=I$, and
$\|H_V\|\le2L^3$, $\|\mathsf C_V\|\le C_3L^4$. The image of $\mathsf C_V$ is
$\ker K_V$. Both maps are holomorphic in the admitted backgrounds and have
local rows.
\end{lemma}

\begin{proof}
The frame used in \eqref{eq:fibre-background} puts $J_c(V)$ within
$C_0L^2\epsilon$ of $L^{-3}I$, up to the unitary transformations at its source
and target. Choose $c_3$ so that the relative perturbation in that frame is
less than $1/2$. A Neumann series gives the inverse, its bound and holomorphic
dependence. The path support and \eqref{eq:fibre-pivot-matrix} show that each
pivot inverse acts on its own row. The identities and operator bounds follow
from \eqref{eq:fibre-linear-lift} and Lemma~\ref{lem:fibre-taylor}. Finally a
vector in $\ker K_V$ is uniquely determined by its restriction to $E_f$,
because the difference of two such vectors with the same restriction is
supported on $P$, where $K_V$ is invertible.
\end{proof}

\subsection{Exact straightening and the fibre graph}\label{sec:fibre-straightening}

\begin{theorem}[One-step fibre chart]\label{thm:fibre-chart}
Under the hypotheses of Lemma~\ref{lem:fibre-pivot}, there is a holomorphic
change of variables $A=\Psi_V(t)$ for $\|t\|_\infty<c_4L^{-5}$, with
\begin{equation}\label{eq:fibre-straight}
 \pi\Psi_V(t)=\pi t,\qquad P_V(\Psi_V(t))=K_Vt,\qquad
 \|D\Psi_V(t)-I\|\le C_4 L^5\|t\|_\infty.
\end{equation}
The inverse on the image is $A\mapsto A+H_VN_V(A)$. For $\|u\|_\infty<c_5L^{-9}$ the map
\begin{equation}\label{eq:fibre-graph}
 \Phi_V(u)=\Psi_V(\mathsf C_Vu)
\end{equation}
gives the unique solution of $P_V(A)=0$, $\pi A=u$ in a full logarithm ball of
radius $c_6L^{-5}$. It obeys $\|\Phi_V(u)\|_\infty\le C_5L^4\|u\|_\infty$ and
$\|D\Phi_V(u)\|\le C_5L^4$. Its pivot at $b_c$ depends only on $V|_{E_c}$ and
$u|_{E_c\cap E_f}$. The chart is real when $V,u$ are real.
\end{theorem}

\begin{proof}
For $t$ in the stated ball, solve for a field $d$ on $C$ by
\begin{equation}\label{eq:fibre-contraction}
 d=N_V(t-H_Vd),\qquad \Psi_V(t)=t-H_Vd.
\end{equation}
For fixed $t$, restrict to $\|d\|_\infty\le C_6L^2\|t\|_\infty^2$ with
$C_6\ge4C_0$. Decrease $c_4$ so that $\|H_Vd\|_\infty\le\|t\|_\infty$ and
$2\|t\|_\infty<c_0/L$. The second bound in \eqref{eq:fibre-remainder} makes
this ball invariant. The last bound there gives a Lipschitz constant at most
$4C_0 L^5\|t\|_\infty<1/2$. Iteration from zero gives a locally uniform
holomorphic limit in $t,V$. The real subspace is preserved by the iteration.

Substitution into $P_V(A)=K_VA+N_V(A)$ and $K_VH_V=I$ prove the first two
identities in \eqref{eq:fibre-straight}. Differentiating
\eqref{eq:fibre-contraction} gives
\[
 D_td=(I+DN_V(\Psi_V(t))H_V)^{-1}DN_V(\Psi_V(t)).
\]
The inverse in this formula has norm at most two. This proves the derivative
bound. The formula for the inverse of $\Psi_V$ follows because
$d=N_V(\Psi_V(t))$.

Choose $c_6$ so that $\|H_V\|\sup_{\|A\|<c_6L^{-5}}\|DN_V(A)\|<1/2$, and then
choose $c_5$ so that \eqref{eq:fibre-graph} is in this ball and $\mathsf C_Vu$ is in the domain of $\Psi_V$. This is possible by
Lemma~\ref{lem:fibre-pivot} and the bounds above. Two solutions $A,A'$ in the
full ball with the same $u$ have a pivot-supported difference. Their
constraint implies $A-A'=-H_V(N_V(A)-N_V(A'))$. The strict derivative bound on
this convex ball gives $A=A'$. Each equation for $d(c)$ involves only the free
part of $t$ on $E_c$ and the single pivot $b_c$. This proves the stated
support property.

For real arguments satisfying \eqref{eq:fibre-real-admission}, the change in a
tested plaquette product is at most $4\|A\|_\infty$. Indeed, telescope its
four unitary factors and use $\|e^{A(b)}-I\|\le\|A(b)\|$, with the same bound
on a reversed factor. Decrease $c_4,c_6$, and then $c_5$ accordingly, so that
$4\|A\|_\infty<1/(400L^2)$ on
the logarithm balls used here. Then the chart lies in the first branch of the
real map. The larger balls in Lemma~\ref{lem:fibre-taylor} serve to estimate
the analytic formula.
\end{proof}

For nearby targets write \(t=\mathsf C_Vu+H_Vv\) and set
\begin{equation}\label{eq:fibre-product-chart}
 \Theta_V(u,v)=\Psi_V(\mathsf C_Vu+H_Vv).
\end{equation}
For \(\|u\|_\infty<c_7L^{-9}\) and \(\|v\|_\infty<c_7L^{-8}\)
the bound on \(t\) is \((C_3+2)c_7L^{-5}\). Choose \(c_7\leq c_5\)
so that this is below both \(c_4L^{-5}\) and
\(\tfrac12c_6L^{-5}\). The contraction proof gives
\(\|\Theta_V(u,v)\|_\infty\leq2\|t\|_\infty\). Hence
\[
 (\pi,P_V)(\Theta_V(u,v))=(u,v),\qquad
 \mathcal M_{\rm an}(e^{\Theta_V(u,v)}V)(c)=e^{v(c)}W(c).
\]
These equations give the inverse coordinates on the image.
For real \(V,u,v\) with \eqref{eq:fibre-real-admission}, the
four-link bound in the proof of Theorem~\ref{thm:fibre-chart}
puts every tested plaquette in the first branch. The identity
then holds for the real Borel map \(\mathcal M\). At \(v=0\)
the coordinate is \(\Phi_V(u)\). For the fibre graph in Theorem~\ref{thm:fibre-chart}, the substitution
\(u=gB\), \(g>0\), gives radius \(c_5L^{-9}/g\) in \(B\).

The first two derivatives of the graph at the reference point are
\begin{equation}\label{eq:fibre-second-jet}
 D\Phi_V(0)=\mathsf C_V,\qquad
 D^2\Phi_V(0)[u_1,u_2]
  =-H_V D^2P_V(0)[\mathsf C_Vu_1,\mathsf C_Vu_2].
\end{equation}
Indeed, the second derivative vanishes on the free bonds and is determined on
the pivots by twice differentiating the constraint. On smaller uniform balls
all fixed-order derivatives of each output row have bounds depending only on
$L$ and the order. The Cauchy formula on scalar circles gives these bounds
from the holomorphic chart, since each circle can be chosen using the same
margin in the supremum radius. At any fixed order, the number of input bonds
in a row and the number of rows containing an input bond are bounded in terms
of $L$. This gives volume-independent absolute row and column bounds for the
coordinate kernels of the derivatives.

The same differentiation identifies the second variation of a
scalar function along a constrained chart. More generally, let
$P(0)=0$, let $K=DP(0)$ have a specified right inverse $H$,
and let $C$ parametrize $\ker K$. For a chart
$\Phi(u)=Cu-HN(\Phi(u))$ with $\Phi(0)=0$ and $N=P-K$,
differentiation gives
\[
 D\Phi(0)=C,\qquad D^2\Phi(0)[u,v]=-H D^2P(0)[Cu,Cv].
\]
For any twice differentiable scalar $S$,
\begin{equation}\label{eq:fibre-composition-hessian}
 D^2(S\circ\Phi)(0)[u,v]
 =D^2S(0)[Cu,Cv]-DS(0)\bigl[H D^2P(0)[Cu,Cv]\bigr].
\end{equation}
This is the ordinary second chain rule. If $DS(0)=0$, the
second term vanishes. In a
different gauge chart the right inverse $H$ must be supplied
and bounded for that chart. The pivot inverse above supplies
this formula for the one-step plain constraint.

Suppose instead that \eqref{eq:fibre-background} holds in complex
gauges $s_c$ whose condition numbers satisfy
\[
 \|s_c(x)\|\|s_c(x)^{-1}\|\le\kappa,\qquad \kappa\ge1.
\]
Retain \(\epsilon\le c_0/L\) for the Taylor estimates and
\(\epsilon\le c_3L^{-5}\) for the pivot and fibre graph.
Define the continued coarse field in the original frame by
\[
 W(c)=s_c(c_-)^{-1}\mathcal M_{\rm an}(V^{s_c})(c)s_c(c_+).
\]
It agrees with \(\mathcal M_{\rm an}(V)(c)\) when the original row
chart is also admitted. Compute the relative logarithm in each local
frame and conjugate it back by \(s_c(c_-)^{-1}\).
Two near-identity frames give the same result
wherever both admit the calculation. To see this for a contour mean, normalize
the family at its first member and start the fixed-point iteration in
Lemma~\ref{lem:avg-mean} at zero in both frames. The relative matrices are
conjugate at the initial endpoint. The iterates remain conjugate since the
logarithm series converges in each frame. The same series argument applies to
the closed loops and the output ratio. Thus the continuation is holomorphic on
the union of these patches.

Conjugation contributes a factor $\kappa$ for each input direction and for the
output. The radius for the direct Taylor bound becomes $c_0/(\kappa L)$. The
operator bounds are
\[
 \|K_V\|\le C_0\kappa^2 L,\qquad
 \|DN_V(A)\|\le C_0\kappa^3 L^2\|A\|_\infty,\qquad
 \|H_V\|\le2\kappa^2L^3.
\]
The perturbation test for $J_c$ is made inside the near-identity frame, with
$\epsilon\le c_3L^{-5}$. Thus $\|\mathsf C_V\|\le C_3\kappa^4L^4$. The
contraction works for $\|t\|_\infty<c_4\kappa^{-5}L^{-5}$. The resulting graph
admits the free radius $c_5\kappa^{-9}L^{-9}$ after a fixed decrease of $c_5$.

The corresponding product chart is also given by
\eqref{eq:fibre-product-chart}. Its input must satisfy
\[
 \|\mathsf C_Vu+H_Vv\|_\infty
       <\min\{c_4,c_6/2\}\kappa^{-5}L^{-5}.
\]
Sufficient radii in the original coordinates are
\(\|u\|_\infty<c_7\kappa^{-9}L^{-9}\) and
\(\|v\|_\infty<c_7\kappa^{-7}L^{-8}\), after decreasing
\(c_7\) so that \((C_3+2)c_7<\min\{c_4,c_6/2\}\).
The bound \(\|\Psi_V(t)\|_\infty\le2\|t\|_\infty\)
then places the image in the full uniqueness ball of radius
\(c_6\kappa^{-5}L^{-5}\).

\subsection{The real fibre density}\label{sec:fibre-measure}

Use the real colour inner product $\langle
X,Y\rangle=-\tfrac12\operatorname{tr}(XY)$ on $\mathfrak g$ and its product
Lebesgue measure in the bond variables. For $A\in\mathfrak g$, let $r=\|A\|$.
Then $A^2=-r^2I$ and $e^A=\cos r\,I+(\sin r/r)A$. Under the unit-quaternion
identification, normalized Haar measure is round measure on the unit
three-sphere divided by its area $2\pi^2$. Its radial factor is $\sin^2r$,
while the Euclidean radial factor is $r^2$. Consequently, near zero,
\begin{equation}\label{eq:fibre-haar}
 d(e^A V)=\sigma_0 j(A)\,d^3A,
 \qquad \sigma_0=(2\pi^2)^{-1},\qquad j(A)=(\sin r/r)^2,
\end{equation}
with $j(0)=1$. Right translation by a real $V$ preserves Haar measure.

\begin{proposition}[Local disintegration of product Haar measure]\label{prop:fibre-density}
Let \(V\) be real and satisfy \eqref{eq:fibre-real-admission}
and the hypotheses of Theorem~\ref{thm:fibre-chart}. Restrict
fine Haar measure to the group-valued image
\(U=e^{\Theta_V(u,v)}V\) of the product chart in
\eqref{eq:fibre-product-chart}, with
\(\|u\|_\infty<c_7L^{-9}\) and \(\|v\|_\infty<c_7L^{-8}\).
Here \(u\) is a real raw free-link field and \(v\) a real raw
coarse-link field. The target is \(W'=e^vW\).
Relative to Haar measure on this target chart, the unnormalized
fibre kernel has density in the \(u\) coordinates
\[
 \mathfrak d_V(u,v)=
 \frac{\sigma_0^{|E_f|}\prod_{b\in E}j(\Theta_V(u,v)(b))}
      {\prod_{c\in C}j(v(c))\,
       \prod_{c\in C}\det D_{A(b_c)}P_V(\Theta_V(u,v))(c)} .
\]
The density is continuous throughout the real product chart.
At the exact target \(W\), \(v=0\) and this yields
\begin{equation}\label{eq:fibre-density}
 \sigma_0^{|E_f|}
 \prod_{b\in E}j(\Phi_V(u)(b))
 \prod_{c\in C}\det\bigl(D_{A(b_c)}P_V(\Phi_V(u))(c)\bigr)^{-1}
 \prod_{b\in E_f}d^3u_b.
\end{equation}
The pivot determinants are positive. Each determinant and each
Haar factor at a solved pivot reads only its own row. A free-link
Haar factor reads that link.
\end{proposition}

\begin{proof}
The Jacobian of \(A\mapsto(\pi A,P_V(A))\), with the free
variables ordered first, is block triangular. Its free block
is the identity and its pivot block is diagonal in \(c\),
with entries \(D_{A(b_c)}P_V(A)(c)\). At zero these entries have
positive determinant by the invertible straight segment to
\(L^{-3}I\) in their unitary frames. The frame changes have
determinant one. The inequality \(\|DN_V(A)H_V\|<1/2\)
preserves invertibility on the whole logarithm ball, so the
determinants keep their signs along \(A\mapsto sA\),
\(0\leq s\leq1\).

The real change of variables gives, for every bounded Borel
test \(F(U,W')\) on this graph,
\[
 \int_{\rm chart}F(U,\mathcal M(U)|_C)\,\prod_{b\in E}dU_b
 =
 \int_{\rm target\ chart} dW'
   \int_{\rm free\ ball}
      F(e^{\Theta_V(u,v)}V,W')\,\mathfrak d_V(u,v)\,du .
\]
Here \(du=\prod_{b\in E_f}d^3u_b\) and \(dW'\) is product
normalized Haar on the target links. The fine exponential
Jacobian contributes \(\sigma_0^{|E|}\prod_b j(A(b))\).
The triangular Jacobian contributes the reciprocal pivot
determinants. In the coordinates \(W'=e^vW\), the target Haar
differential is \(\sigma_0^{|C|}\prod_c j(v(c))\,dv\), which
gives the denominator in \(\mathfrak d_V\).
This construction specifies the continuous density also at
individual targets. Setting \(v=0\) uses \(\Theta_V(u,0)=\Phi_V(u)\)
and \(j(0)=1\), proving \eqref{eq:fibre-density}.

Let \(\delta_G\) be the unit Dirac distribution at \(I\)
for normalized Haar measure, and let
\(\delta_{\mathbb R^3}\) be the unit Dirac at zero for
Lebesgue measure. In the real logarithm chart
\(\delta_G(e^v)=\sigma_0^{-1}\delta_{\mathbb R^3}(v)\).
Using these factors in the product chart gives the same
fibre kernel at \(W\).
\end{proof}

A Borel restriction \(B\) of the fine chart multiplies the kernel
by \(\mathbf 1_B(e^{\Theta_V(u,v)}V)\). The change-of-variables
identity follows by inserting this factor into the test. At a
fixed target this convention uses the stated pointwise restriction.
It gives a version of the restricted kernel without requiring
continuity of the indicator.

Differentiating $(\pi,P_V)\circ\Psi_V=(\pi,K_V)$ gives another expression for
the determinant factor in \eqref{eq:fibre-density},
\[
 \det D\Psi_V(\mathsf C_Vu)\prod_c\det J_c(V)^{-1}.
\]
In the convention $Q_V=K_V/L$ for a fluctuation inserted as $e^{A/L}$, the
pivot matrix is $J_c(V)/L$. Accordingly,
\[
 \det J_c(V)^{-1}=L^{-3}\det(J_c(V)/L)^{-1}.
\]
The free coordinate substitution $u=gB$ contributes $g^{3|E_f|}$ to \eqref{eq:fibre-density}.

\subsection{Holomorphic continuation of the density}\label{sec:fibre-continuation}

Use $z(A)=-\tfrac12\operatorname{tr}(A^2)$ on $\mathfrak g_\mathbb C$. The
entire function $(\sin\sqrt z/\sqrt z)^2$, interpreted by its power series at
zero, continues $j$. The reciprocal determinants in \eqref{eq:fibre-density}
are holomorphic on the admitted chart. These continuations give a holomorphic
density coefficient. Here a variation of \(V\) keeps the relative
coordinates fixed: at \(v=0\) the target is the moving value
\(W(V)\), using the continued coarse field just defined when complex
frames are needed. To use a fixed absolute target
\(Y\), one must instead admit
\(v(V)=\log(YW^{-1})\) in the product chart
\eqref{eq:fibre-product-chart}.
The measure and its disintegration in
Proposition~\ref{prop:fibre-density} are defined on the real chart.

For local logarithmic densities, each Haar factor has a logarithm with value
zero at $A=0$, from the power series for $\log(1+(j(A)-1))$. For the pivot
factor normalized at $u=0$, put
\[
 B_c(u)=D_{A(b_c)}P_V(\Phi_V(u))(c)\,J_c(V)^{-1}.
\]
In the local frame of row \(c\), denote the normalized matrix
by \(\widetilde B_c(u)\). The local Taylor estimate and the
pivot inverse give
\(\|\widetilde B_c(u)-I\|\leq C L^5
\|\operatorname{Ad}(s_c)\Phi_V(u)\|_{\infty,E_c}\),
where the adjoint acts at each fine link's initial point.
For unitary frames this is at most \(CL^9\|u\|_\infty\).
For frames of condition number at most \(\kappa\) it is at most
\(C\kappa^5 L^9\|u\|_\infty\), by the graph and frame bounds.
Thus a fixed decrease of the free radius puts this matrix
within \(1/2\) of \(I\) in either case. The normalized reciprocal
determinant has logarithm
\(-\operatorname{tr}\log\widetilde B_c(u)\).
Changing frame conjugates this normalized matrix at the coarse
initial site, so the trace defines the same logarithm in every
admitted frame.

At the reference field, the normalized unperturbed pivot
matrix in its local frame is \(L^3J_c=I+S\), with
\(\|S\|<1/2\). The imaginary part of
\(\operatorname{tr}\log(I+S)\) has absolute value less than
\(3\log 2<\pi\). The value of this trace is therefore the
principal logarithm of the normalized determinant and is
independent of the frame. The unnormalized logarithm is
\(-9\log L+\operatorname{tr}\log(I+S)\).
The logarithmic factors normalized at zero have local
fixed-order derivative bounds on smaller balls by Cauchy's
formula. In unitary frames the bounds depend on \(L\) and the
order. With complex frame condition number at most \(\kappa\),
the bounds in the original coordinates may also depend on
\(\kappa\). Derivatives in \(u\) of a pivot Haar factor or a
row determinant read only the free part of that row. A free-link
Haar factor reads its own free link, which need not belong to any
row. Background derivatives of a row factor may read all of
\(V|_{E_c}\), including its pivot. The logarithm of each Haar
factor has zero derivative
at the reference because \(j(A)=1+O(\|A\|^2)\).

% End fibres.tex
% Begin soft.tex
\paragraph{The per-length pivot determinant.}
Write
\(m_c(A)=L^{-1}D_{A(b_c)}P_V(A)(c)\) for the pivot derivative
in the per-length convention. In the local frame at row \(c\),
\eqref{eq:fibre-remainder} and \eqref{eq:fibre-pivot-matrix}
give \(L^4m_c(A)=I+S_c(A)\) with
\[
 \|S_c(A)\|\leq C_0L^5
       \bigl(\epsilon+\|A^{s_c}\|_{\infty,E_c}\bigr).
\]
Shrink the background and fluctuation ceilings by fixed factors
so that the right side is at most \(1/8\). The same shrinking is
available on the complex-frame balls, since those balls control
the increment in each local frame. The product expansion of the
three-dimensional determinant, or its three eigenvalues,
then gives
\begin{equation}\label{eq:fibre-retained-determinant}
 \bigl|L^{12}\det m_c(A)-1\bigr|
     \leq (1+1/8)^3-1<1/2 .
\end{equation}
This determinant is independent of the choice of local frame,
since the adjoint changes of colour coordinates at the source
and target both have determinant one. Thus
\[
 \ell_c(A)=12\log L-\log\bigl(L^{12}\det m_c(A)\bigr)
\]
is holomorphic with the series logarithm about one,
\(\exp(\ell_c(A))=(\det m_c(A))^{-1}\), and
\(|\ell_c(A)-12\log L|\leq C\).
On the real chart this is the positive reciprocal determinant.
The coordinate conversions in the real Haar formula are still
required when using this per-length factor.


\section{Haar normalization of the block gauge}\label{sec:soft-gauge}

The axial condition may be imposed by changing representatives, or
it may be used to weight representatives in a gauge orbit. For the
second operation we need the subgroup that acts on the transports.
Suppose a family of disjoint blocks has fixed centres. At a selection
of sites \(x\ne y\) in these blocks allow independent gauge values
\(g(x)\in G\). The gauge equals \(I\) elsewhere. Let \(\mathcal H\)
be this product group with normalized Haar measure. On the invariant
domain of real fields where the contour means exist, the selected transports obey
\[
 T_{U^g}(y,x)=T_U(y,x)g(x)^{-1}.
\]
Consequently the map from \(\mathcal H\) to their new values pushes
Haar forward to product Haar. A pinned pair, with both endpoints
fixed, has constant transport on each orbit.

\begin{proposition}[Insertion of normalized weights]\label{prop:soft-gauge}
Let \(f_x\) be nonnegative measurable functions on \(G\) of Haar
integral one and let
\(w(U)=\prod_x f_x(T_U(y,x))\) over the selected free sites.
Then \(\int_{\mathcal H}w(U^g)\,\mathrm dg=1\) on the domain above.
For a finite measure \(\nu\) supported on this invariant domain,
drawing \(g\) with density \(w(U^g)\)
and replacing \(U\) by \(U^g\) preserves the integrals of bounded
\(\mathcal H\)-invariant functions. If \(\nu\) is invariant, the
resulting measure is \(w\,\nu\).
\end{proposition}

\begin{proof}
The transport substitutions are independent Haar translations and
inversions. For an invariant function the integrand equals its value
at \(U\) times the density, whose integral is one. For general
bounded test functions and invariant \(\nu\), change \(U\) to \(U^g\)
at fixed \(g\), then integrate over \(g\). Fubini applies since the
weight has orbit integral one.
\end{proof}

\subsection{The quadratic term of the contour weight}

Let $\mathcal P$ be the centres of a collection of disjoint blocks.
Let $W$ be real with $T_W(y,x)=I$ on the contour pairs used below,
and assume their path families have diameter at most $1/80$.
This leaves a uniform reserve in the mean chart. The soft contour
energy is
\[
 \mathcal F(U)=\sum_{y\in\mathcal P}
       \sum_{x\in B(y)\setminus\{y\}}
              \left(1-\tfrac12\operatorname{tr}T_U(y,x)\right).
\]
For a real raw link increment \(a\) and a real parameter \(t\),
define the contour differential by
\[
 C_Wa(y,x)=
  \left.\frac{d}{dt}\log T_{e^{ta}W}(y,x)\right|_{t=0}.
\]
At \(W=I\) it is the arithmetic mean of the oriented path sums.
In the holomorphic continuation the quadratic coefficient is
the bilinear pairing of \(C_Wa\) with itself.
Analyticity of the finite contour map gives, on a fixed smaller ball,
\begin{equation}\label{eq:soft-quadratic}
 \mathcal F(e^{ta}W)=\frac{t^2}{2}
          \sum_{y,x}|C_Wa(y,x)|^2
       +O_L\!\left(|t|^3\sum_{y\in\mathcal P}
                                      \|a\|_{\infty,B(y)}^3\right).
\end{equation}
The constant is uniform over these real backgrounds and over the
number of blocks. To see this, put
$Z_{yx}(t)=\log T_{e^{ta}W}(y,x)$. The finite-path estimates give
$Z_{yx}(t)=tC_Wa(y,x)+O_L(t^2\|a\|_{\infty,B(y)}^2)$.
For $Z\in\mathfrak{su}(2)$,
$1-\tfrac12\operatorname{tr}e^Z=\tfrac12|Z|^2+O(|Z|^4)$.
Substitution proves the formula; the fixed $L^4-1$ contours per block
are absorbed in its constant.

For a site field $\lambda$ supported on the selected free sites
(and hence vanishing at the block centres), put
$a(b)=\lambda(b_-)-\operatorname{Ad}(W(b))\lambda(b_+)$.
Differentiating endpoint covariance gives
$C_Wa(y,x)=\lambda(y)-\operatorname{Ad}(T_W(y,x))\lambda(x)
=-\lambda(x)$. Hence the quadratic form in
\eqref{eq:soft-quadratic} is strictly positive on every nonzero
such gauge direction. If extra nodes are pinned, restrict to
$\lambda=0$ there as well. This is a positivity statement about the
contour weight; identification of all null directions of an action
Hessian requires the corresponding action and boundary hypotheses.

The normalized insertion in Proposition~\ref{prop:soft-gauge} permits
any integrable probability weight on the contour variables. Its
validity therefore does not require matching the parameter in that
weight to the coupling in the Wilson action.

For instance let \(0<\delta\le1\) and put
\[
 z_\delta(s)=\int_G {\bf1}_{\{|u-I|<\delta\}}
              e^{-s(1-\operatorname{tr}u/2)}\,\mathrm du .
\]
For real \(s\) the integrand divided by \(z_\delta(s)\) is such a
weight. To control the denominator on complex coupling balls, set
\(t=1-\cos\theta\), where \(\theta\in[0,\pi]\) is the class angle
of \(u\). The radial density of normalized Haar is
\((2/\pi)\sin^2\theta\,\mathrm d\theta\), by the unit-quaternion
description of \(G\). The cutoff is exactly \(t<\delta^2/2\).
Hence
\begin{equation}\label{eq:soft-z}
 z_\delta(s)=\frac{2}{\pi}\int_0^{\delta^2/2}
                 e^{-st}\sqrt{t(2-t)}\,\mathrm dt .
\end{equation}
The integral is entire. Fix \(s_0\geq1\) with \(s_0\delta^2\geq1\).
For \(a\in[s_0/2,3s_0/2]\) it gives
\(c s_0^{-3/2}\leq z_\delta(a)\leq C s_0^{-3/2}\).
The lower bound follows by integrating over \(0<t<1/(2s_0)\),
which is inside the cutoff. For the upper bound use
\(\sqrt{2-t}\leq\sqrt2\) and extend the resulting integral to
infinity.

For \(s=a+ib\) with \(a>0\), the ratio
\(z_\delta(a+ib)/z_\delta(a)\) is the expectation of \(e^{-ibt}\)
for the probability density proportional to
\(t^{1/2}e^{-at}w(t)\), where
\(w(t)=\sqrt{2-t}\,\mathbf1_{\{0<t<\delta^2/2\}}\)
on \(t>0\), taken to be zero beyond the cutoff.
The function \(w\) is nonincreasing. For the probability
measure \(\mu_a\) with density proportional to \(t^{1/2}e^{-at}\)
on \(t>0\), an increasing function \(f\) satisfies
\[
 \int fw\,d\mu_a-\int f\,d\mu_a\int w\,d\mu_a
  =\frac12\iint(f(t)-f(u))(w(t)-w(u))
                         \,d\mu_a(t)\,d\mu_a(u)\leq0 .
\]
Use \(f(t)=t^2\) and divide by the positive integral of \(w\).
The resulting second moment is at most \(15/(4a^2)\), the
second moment of \(\mu_a\) obtained by two integrations by parts.
Since \(\cos(bt)\geq1-b^2t^2/2\),
\[
 \operatorname{Re}\frac{z_\delta(a+ib)}{z_\delta(a)}
                 \geq 1-\frac{15b^2}{8a^2}.
\]
If \(|s-s_0|\leq s_0/2\), then \(b^2\leq a^2/3\), since
\[
 a^2/3+(a-s_0)^2-s_0^2/4
       =\tfrac43(a-3s_0/4)^2 .
\]
Thus the ratio has real part at least \(3/8\) on this disk.
Together with the real lower bound this gives
\begin{equation}\label{eq:soft-z-complex}
 |z_\delta(s)|^{-1}\leq C s_0^{3/2},
           \qquad |s-s_0|\leq s_0/2 .
\end{equation}
The zero-free disk has a holomorphic logarithm with its real
value at \(s_0\). The cutoff is kept at its fixed real value
throughout this continuation.

The choice of subgroup affects this identity at a seam. Consider an
oriented coarse link leaving a centre \(y\), with target \(V\), and
write one determining product as \(uSP=V\). Here \(S\) is a retained
link, \(P\) is the subsequent transport to the far endpoint, and
\(u\) is the contour variable from \(y\) to the initial endpoint of
\(S\). If that contour variable has weight \(f\), then integration
against the Haar constraint \(uSP=V\) gives exactly
\(f(VP^{-1}S^{-1})\). This follows by right translation in Haar,
or by testing the constraint against a bounded continuous function
of \(V\). 
On a line the centered average on a coarse link is its straight
transport. The path through any shifted site and back to the
terminal centre cancels its backtracks, so every closed loop
in \eqref{eq:avg-map} equals \(I\). For this one-dimensional example,
start from product Haar and insert the normalized contour factors.
For two completed blocks,
take the transports from their centres to their free sites
as rooted tree coordinates. The bond joining the two blocks
is an independent Haar coordinate. If the coarse output for
that internal row is fixed, first integrate this bridge link
against its target constraint. The interval product has
the bridge link between two fixed transports, so this integration
gives one. The only retained bridge is \(S\) on the outer seam,
where the straight product is \(uSP\). Every unused contour
coordinate now contributes the integral of its normalized
weight. The remaining constraint therefore gives
\(f(VP^{-1}S^{-1})\) at fixed \(V,P,S\).
Averaging the gauge at \(y\) alone sends this factor to one by
left Haar translation. Consequently, if the remaining datum measure
\(\nu_0\) is invariant under this node action, the orbit average
of \(f(VP^{-1}S^{-1})\nu_0\) is \(\nu_0\). This follows by
changing the datum by the node gauge in a bounded test integral
and then integrating its normalized factor.

A transformation at \(y\) alone multiplies this argument
on the left. For nonconstant \(f\) the resulting factor is in general
changed. Using the same gauge value at \(y\) and the initial
endpoint of \(S\) conjugates the argument, preserving it for a class
function \(f\). Integrating \(S\) against Haar at fixed \(V,P\)
gives one by the normalization of \(f\).

For an arbitrary finite input measure $\nu$ on this invariant domain, let
$\overline\nu=\int_{\mathcal H}(U\mapsto U^g)_*\nu\,dg$.
The weighted draw in Proposition~\ref{prop:soft-gauge} has measure
\begin{equation}\label{eq:soft-arbitrary-input}
                         w\,\overline\nu .
\end{equation}
To verify it, test against a bounded measurable $F$ and write the
draw as $\int\nu(dU)\int dg\,(wF)(U^g)$. This is the defining
integral of $wF$ against $\overline\nu$. The orbit normalization
preserves total mass. If some contour pairs are outside their
defined domains, omit their weight factors on the invariant
domain where this set of pairs is fixed.

The tree slice can have arbitrary prescribed tree transports.
On a rooted forest, fix $S_e\in G$ for every tree edge from parent
to child. Vary the gauges only at nonroot vertices of the forest,
and fix the root and exterior gauges to \(I\).
For any real field $U$, the equations $(U^g)_e=S_e$ have the unique solution
$g(x)=S_e^{-1}g(y)U(e)$ along $e=(y,x)$.
Let $\pi_T^S U$ denote the resulting field. Conditional on previous
tree variables, each equation is a Haar translation and inversion.
The Haar delta product therefore gives the point mass at the unique
chosen tree gauge. For every bounded Borel test \(\Phi\) with specified
pointwise values, and then for tests integrable against the induced
slice measure,
\begin{equation}\label{eq:tree-target-delta}
 \int\prod_{x\ne\mathrm{roots}}dg(x)
   \prod_{e\in T}\delta_G((U^g)_e S_e^{-1})\,\Phi(U^g)
                       =\Phi(\pi_T^S U).
\end{equation}
There is no determinant from this group-valued tree elimination.
For a gauge-invariant $\Phi$ its value on the right is $\Phi(U)$.
Consequently removing a tree delta gives equality after integration
against invariant tests, rather than equality of the two densities
as functions on all links.

For non-invariant tests the full averaged density can also be identified.
Assume that every tree edge is a variable link. Let \(dU\) be product
normalized Haar on the variable links, with exterior data held fixed
and preserved by the rooted gauges. Write \(\sigma_T^S\) for the
measure that fixes each tree link to \(S_e\) and keeps product
normalized Haar on the off-tree links. For any
\(\Psi\in L^1(\sigma_T^S)\),
\begin{equation}\label{eq:tree-slice-orbit-density}
 \int_{\mathcal H_T}(U\mapsto U^g)_*(\Psi\,\sigma_T^S)\,dg
       =(\Psi\circ\pi_T^S)\,dU,
\end{equation}
where \(\mathcal H_T\) consists of the gauges at nonroot vertices,
fixed to \(I\) at roots and exterior sites. Indeed, each field has a
unique representation \(U=R^g\), with \(R\) in the tree slice.
The triangular Haar substitution of
Section~\ref{sec:gauge-straightening} gives
\(dU=\sigma_T^S(dR)\,dg\), and uniqueness of tree fixing gives
\(\pi_T^S(R^g)=R\). Testing against a bounded Borel function proves
\eqref{eq:tree-slice-orbit-density} by Fubini, also for complex
\(\Psi\). The same substitution shows that the pullback of a
slice-null set is Haar-null, so the density is well defined for
the \(L^1\) class of \(\Psi\).

One must also distinguish two uses of a constraint delta. Integrating
its Haar target gives
$\int_G\delta_G(aY^{-1})\,dY=1$ for every $a\in G$.
Eliminating a fine link through the averaging map instead uses the
Jacobian in Proposition~\ref{prop:fibre-density}. That derivative
is not a Haar translation, and its determinant cannot be discarded.

\subsection{Changes of retained Haar coordinates}
\label{sec:guarded-haar}

When a constraint is integrated over its output links, it is useful
to change the frame of those links while leaving the variables that
parametrize the input unchanged. Such a change can be made before
any support test is evaluated.

\begin{proposition}[Change of target frame]\label{prop:guarded-haar}
Let \(V(c)\), \(c\in C_{\rm int}\), be a finite family of target
variables with product normalized Haar measure. Hold all other data
fixed, including any prescribed targets on a disjoint set
\(C_{\rm fix}\). Let \(u\) and \(\tau\) be real site gauges
depending only on these fixed data. If those data are subsequently
integrated, require this dependence to be measurable. Set \(Y=V^u\) on
\(C_{\rm int}\) and \(Y=V\) on \(C_{\rm fix}\), and require
\(u\) to fix the prescribed target values. The change preserves
product Haar measure and rewrites
the old and mixed target arguments as
\[
 V=Y^{u^{-1}},\qquad V^\tau=Y^{\tau u^{-1}}.
\]
For an averaging graph with output \(a(c)=\mathcal M(v)(c)\),
its target delta on each integrated row \(c\in C_{\rm int}\) satisfies
\begin{equation}\label{eq:guarded-haar-delta}
 \delta_G\bigl(a(c)(Y^{u^{-1}}(c))^{-1}\bigr)
          =\delta_G\bigl(a^u(c)Y(c)^{-1}\bigr).
\end{equation}
In an integrable density on the graph, a factor with old argument
\(F(V,\xi)\), including its indicator or its domain of definition,
becomes \(F(Y^{u^{-1}},\xi)\). A factor with mixed argument keeps
\(Y^{\tau u^{-1}}\) in that argument. The variables \(\xi\) are held
fixed in this change.
\end{proposition}

\begin{proof}
On each integrated link \(Y(c)=u(c_-)V(c)u(c_+)^{-1}\).
Left and right translations preserve Haar measure, and composition
of gauge transformations gives the order \(\tau u^{-1}\).
The argument of the delta on the right of
\eqref{eq:guarded-haar-delta} is the conjugate by \(u(c_-)\)
of the argument on the left. The Haar delta at \(I\) is invariant
under conjugation. Equivalently, with its Haar target differential,
each delta is the point mass at the corresponding output in its
chosen frame. On a prescribed target row the same notation is used
only for an already defined constraint measure, such as the local
fibre kernel of Proposition~\ref{prop:fibre-density}.
This proves the change for the graph measure at fixed input \(v\).
The measurable dependence on any integrated input data then allows
Fubini to give the identity after integrating those variables.
One may multiply that measure by the full measurable density,
including all the tests.

To use this identity as a Haar change of target coordinates in a
resolved local chart, additionally require \(u\) to be independent
of those target coordinates after expressing the input data in that
chart. In particular it must not acquire target dependence through
a reconstructed pivot. Then substitute \(V=Y^{u^{-1}}\) in every
reconstructed quantity, including the pivots, \(\tau\), and all
test and chart domains, while holding the free coordinates fixed.
An input-dependent gauge without this extra property still permits
the ambient graph identity above, but need not induce an invertible
change of target coordinates on the resolved graph. Disjoint target
stocks can be changed separately whenever the gauge for each change
is independent of the targets being integrated at that step.
\end{proof}

The same rule applies to a prescribed local test when its full
domain and its value are gauge invariant. If the test reads all
the endpoints of a determining row, its transformed input uses
\(u\) at both endpoints. Endpoints pinned in the old frame have
\(u=I\) there. If a second reference and a fluctuating field are
written in the same frame, applying the same gauge to both
conjugates their relative link at its initial site. The norm of
this relative link minus \(I\), and the norm of its logarithm
where defined, are unchanged. Consequently the change of frame
does not require any of the tests to pass. It is enough to impose
the frame constraints on their actual read sets, including the
neighborhoods on which the test is defined.

The fine and target coordinates can also change together. The following
identity supplies the corresponding change for the resolved fibre kernel.

\begin{corollary}[Joint gauge covariance]\label{cor:fibre-gauge-covariance}
Let \(V\) and its real product chart satisfy the hypotheses of
Proposition~\ref{prop:fibre-density}. Fix a real fine gauge \(g\),
independent of the chart coordinates being integrated, and let
\(\bar g\) be its restriction to the coarse sites. Transform the
reference and all input data by these gauges; any data required to
remain fixed must be preserved. Put
\[
 V'=V^g,\qquad
 u'(b)=\operatorname{Ad}(g(b_-))u(b),\qquad
 v'(c)=\operatorname{Ad}(\bar g(c_-))v(c).
\]
The transformed chart has the same radii and satisfies
\begin{equation}\label{eq:fibre-joint-covariance}
 \Theta_{V'}(u',v')(b)
   =\operatorname{Ad}(g(b_-))\Theta_V(u,v)(b),\qquad
 \mathfrak d_{V'}(u',v')=\mathfrak d_V(u,v).
\end{equation}
The free Lebesgue measure and target Haar measure are preserved.
Every extra factor, including a test domain, is pulled back through
the inverse change on all of its arguments.
\end{corollary}

\begin{proof}
Replace each row frame \(s_c\) by \(s_cg^{-1}\). This leaves its
smallness bound unchanged. Real branch tests and the supremum balls
are preserved because unitary conjugation is an isometry. Endpoint
covariance gives
\(P_{V'}(\operatorname{Ad}_g A)=\operatorname{Ad}_{\bar g}P_V(A)\),
where each adjoint acts at the tail of its link. Differentiating this
identity and using the same pivot stock gives covariance of
\(K_V,H_V,\mathsf C_V\) and \(N_V\). The unique fixed point in
\eqref{eq:fibre-contraction}, applied to
\(t=\mathsf C_Vu+H_Vv\), therefore gives the first identity in
\eqref{eq:fibre-joint-covariance}, including when \(v\ne0\).

The pivot derivative in row \(c\) is multiplied on the left by
\(\operatorname{Ad}(\bar g(c_-))\) and on the right by
\(\operatorname{Ad}(g(b_{c,-})^{-1})\). Both maps have determinant
one. All Haar factors in Proposition~\ref{prop:fibre-density} are
conjugation invariant. Its density formula proves the second identity.
The same determinant-one property preserves free Lebesgue measure,
and target links change by left and right Haar translations.
Substitution in an integrable test proves the last assertion.
\end{proof}

The identity may be used conditionally when \(g\) depends measurably
on exterior data held fixed during this change. Its coordinate
independence is still required after reconstruction of the fine links.

\subsection{Local reads and plaquette comparison}
\label{sec:guarded-plaquettes}

Consider a tested level-\(q\) plaquette whose links are either
given at level \(q\) or completed by one step from level \(q-1\).
The test reads the given links and, for each completed link,
the path links and branch-test plaquettes in the projected union
of the two level-\((q-1)\) blocks below its endpoints. This includes
the links joining the blocks and any opposite-face periodic tests.
This follows from Proposition~\ref{prop:avg-map}, including
the branch tests. Its whole footprint is within distance
\(C h_q\) of each corner, using the directed lift if necessary.
In particular, if the test reads a determining row at level
at most \(\ell\) and that row is read either directly or by one
completion, then \(q\leq\ell+1\).

Here is the distance count used with nested boxes. Suppose that
\(C_i\) are increasing rectangular boxes, and every carried
level-\(i\) link in question has one endpoint outside \(C_i\).
If the distance from \(C_{i-1}\) to \(C_i^c\) exceeds \(h_i\),
that link cannot have an endpoint in \(C_{i-1}\).
Indeed its two endpoints are at distance \(h_i\).
At the first outer scale \(i=\ell+1\), the same argument excludes
endpoints in the inset of \(C_{\ell+1}\) of depth \(2h_{\ell+1}\).
Assume the separation \(\operatorname{dist}(C_{i-1},C_i^c)>h_i\)
at every \(i>\ell+1\). The same inset lies in \(C_{i-1}\) at each
such level. A cutoff that hides only links reaching this inset can
therefore miss only levels at most \(\ell\).
For the distance between two test footprints, set
\(w=\operatorname{dist}(C_{\ell+1},C_{\ell+2}^c)\).
Suppose their respective points lie within distances \(a\) of
\(C_{\ell+1}\) and \(b\) of \(C_{\ell+2}^c\). The triangle
inequality bounds the length of any connecting path below by
\(\max\{0,w-a-b\}\). These allowances include any offset of a
footprint from its corresponding set, as well as its diameter
when the bound is given only for one anchor point. Thus the bound
in units of \(h_{\ell+1}\) is \(w/h_{\ell+1}-C'\) if
\(a+b\le C'h_{\ell+1}\). This count concerns the locations read
by the tests.

For real fields \(U,V\) on the same plaquette, unitarity gives
\begin{equation}\label{eq:guarded-four-link}
 \|U(\partial p)V(\partial p)^{-1}-I\|
   \leq \sum_{b\in\partial p}\|U(b)V(b)^{-1}-I\|.
\end{equation}
Replace one factor at a time in the ordered product and use
unitarity of the other factors. Reversing a link leaves its
relative norm unchanged. Hence if the reference plaquette
defect is at most \(\epsilon\) and each of the four relative
link defects is at most \(\delta\), then the perturbed
plaquette defect is at most \(\epsilon+4\delta\).
For a link completed by averaging, a raw perturbation
\(\|a\|_\infty\leq\delta\) on its full double block gives a
relative defect \(C_L\delta\) when the segment \(e^{sa}V\),
\(0\leq s\leq1\), lies in a fixed smaller first-branch chart.
Indeed the first-branch formula is
uniformly holomorphic in the raw perturbation in the local
comb gauge of Lemma~\ref{lem:avg-comb}, and integration of its
bounded derivative gives the estimate. The real chart must
have a buffer on all plaquettes that select the branch.
Thus the same four-factor comparison applies when some of the
links are completed and some are prescribed. These local bounds
apply on a specified support, while a statement about the mass
of that support requires estimates for the density.

\subsection{A common tree and the conditional Haar measure}
\label{sec:guarded-tree}

Let \(P\) be an embedded rectangular lift with at most \(M\) sites
in each direction. Let \(T\) be a rooted spanning tree of its
literal rectangular graph.
Suppose real fields \(U,V\) are both equal to \(I\) on the
tree and their plaquette defects on \(P\) are at most
\(\epsilon_U,\epsilon_V\), respectively. For any off-tree bond
\(b\) of this graph, close \(b\) by the two paths from its endpoints
to the root. Since the tree links are \(I\), the rooted cycle
products are precisely \(U(b)\) and \(V(b)\).
For either field, take its coordinate-comb gauge rooted at the
same site. The cycle product is unchanged, and each link in it
is within \(C M\) times the plaquette defect of \(I\) by
Lemma~\ref{lem:avg-comb}. The rooted cycle has at most
\(2|P|-1\) links, counted with their occurrences.
Telescoping the unitary product gives
\[
 \|U(b)V(b)^{-1}-I\|\leq C_M(\epsilon_U+\epsilon_V).
\]
On the small logarithm branch the same bound holds for
\(\|\log(U(b)V(b)^{-1})\|\) after enlarging \(C_M\).
A prescribed entrance bond common to the two fields may be
included in the tree by holding its exterior endpoint fixed.
The comparison inside \(P\) then uses its internal endpoint
as the root.

For a finite forest of variable links with one fixed root in each
component, the tree-coordinate substitution of
Section~\ref{sec:gauge-straightening} gives product Haar
in the nonroot site gauges and the variable off-tree links of the slice.
A prescribed entrance determines a fixed gauge value at the internal
root and supplies no Haar integration variable.
Indeed each child gauge is a left translate of the link from
its parent, conditional on the preceding tree links. After the
tree variables are fixed, all off-tree links undergo independent
left and right translations. This remains true conditionally
on any fixed exterior data on which the chosen rooted forest
depends. If the exterior data are integrated, the forest and the
prescribed entrance values must be measurable functions of those data.
When the entire integrand is pointwise invariant under these
rooted gauges at that fixed exterior, the integrand in the
new coordinates is its value on the tree slice. The normalized
integrals of the gauge coordinates are then one, giving equality
between the unconstrained integral and the slice integral.
If the integrand also reads variables attached to other scales,
this argument requires invariance of those reads under the
lifted site gauges. In particular a gauge which is constant on
the fine sites of each coarse pillar induces its prescribed
action on the coarse targets by covariance of \(\mathcal M\).
The pins and the support of the gauge must preserve any exterior
or auxiliary read that is held fixed.

% End soft.tex
% Begin back.tex
\clearpage
\appendix
\section{Index of notation and conventions}\label{app:notation}

\begingroup
\setlength{\tabcolsep}{4pt}
\renewcommand{\arraystretch}{1.2}
\begin{longtable}{@{}p{0.19\linewidth}p{0.55\linewidth}p{0.21\linewidth}@{}}
\hline
Symbol & Meaning & Location\\
\hline
\endhead
\(B(y),B^j(y)\) & Centred block and its depth-\(j\) version
    & \S\ref{sec:lattice}\\
\(U^g,F_p\) & Site gauge action and unscaled curvature
    & \S\ref{sec:lattice}\\
\(T_U(y,x)\) & Mean of internal transports from \(y\) to \(x\)
    & \S\ref{sec:avg-map}\\
\(\mathcal M,\mathcal M_{\rm an}\) & Real Borel link average and analytic path formula
    & \S\ref{sec:avg-map}\\
\(Q\) & Flat plain differential in raw logarithms
    & \eqref{eq:avg-flat-link}\\
\(\varepsilon(U),\mathcal E(U)\) & Maximum plaquette angle and Wilson energy
    & \S\ref{sec:energy}\\
\(R_W g\) & Transported site average of a gauge transformation
    & \eqref{eq:gauge-site-nonlinear}\\
\(F_W(a)\) & One-step straightened link coordinate
    & \eqref{eq:gauge-straighten}\\
\(S_W\) & Covariant straight-segment row
    & \eqref{eq:gauge-straight-row}\\
\(a_j,v_j\) & Relative link coordinate and accumulated site factor
    & \eqref{eq:gauge-recursion}\\
\(\Omega_j,\Gamma_j\) & Separated regions and determining coarse rows
    & \S\ref{sec:layered-geometry}\\
\(\xi\) & Largest pillar spacing at the endpoints or corners
    & \S\ref{sec:layered-geometry}\\
\(r_j,\tau_{W,a}\) & Effective finite gauge and its infinitesimal site map
    & \eqref{eq:gauge-effective-chain}, \eqref{eq:gauge-tau}\\
\(D_V\) & Covariant difference with sign from the gauge action
    & \eqref{eq:gauge-D}\\
\(Q'_j(W)\) & Transported site row from depth \(j\)
    & \eqref{eq:gauge-site-row}\\
\(\widehat Q_j\) & Flat straightened link row per unit length
    & \eqref{eq:gauge-flat-row}\\
\(X_\alpha,Y_\alpha\) & Link and site norms controlling differences and H\"older variation
    & \S\ref{sec:smooth-node-gauge}\\
\(E,E_q\) & Smooth node lift and smooth exact-average lift
    & \S\ref{sec:smooth-node-gauge}\\
\(\mathcal C_U,F_h\) & Transported curl and physical plaquette curvature
    & \S\ref{sec:local-variations}\\
\(P_V,K_V\) & Plain fibre constraint and its differential
    & \eqref{eq:fibre-constraint}\\
\(J_c(V)\) & Pivot coefficient of a plain row
    & \eqref{eq:fibre-pivot-matrix}\\
\(\Phi_V\) & Fixed-average graph over the free links
    & \eqref{eq:fibre-graph}\\
\hline
\end{longtable}
\endgroup

\begin{thebibliography}{Ba87}
\small
\bibitem[Ba85a]{Bal85a}
T.~Ba{\l}aban, \emph{Averaging operations for lattice gauge theories},
Comm. Math. Phys. \textbf{98} (1985), 17--51.
\href{https://doi.org/10.1007/BF01211042}{doi:10.1007/BF01211042}.

\bibitem[Ba85b]{Bal85g}
T.~Ba{\l}aban, \emph{Spaces of regular gauge field configurations on
a lattice and gauge fixing conditions},
Comm. Math. Phys. \textbf{99} (1985), 75--102.
\href{https://doi.org/10.1007/BF01466594}{doi:10.1007/BF01466594}.

\bibitem[Ba87]{Bal87}
T.~Ba{\l}aban, \emph{Renormalization group approach to lattice gauge
field theories. I. Generation of effective actions in a small field
approximation and a coupling constant renormalization in four
dimensions}, Comm. Math. Phys. \textbf{109} (1987), 249--301.
\href{https://doi.org/10.1007/BF01215223}{doi:10.1007/BF01215223}.
\end{thebibliography}
\end{document}


% End back.tex
